\documentclass[12pt,a4paper]{article}
\usepackage[utf8]{inputenc}
\usepackage[spanish, es-tabla]{babel} % es-tabla para que diga "Tabla" en vez de "Cuadro"
\usepackage{amsmath, amsfonts, amssymb, amsthm} 
\usepackage{mathtools}
\usepackage{physics} 
\usepackage{tikz} 
\usepackage{tikz-3dplot} % Agregado para mejores gráficos 3D si es necesario, o usando tikz estándar
\usepackage{pgfplots} 
\pgfplotsset{compat=1.18}
\usepgfplotslibrary{fillbetween}
\usepackage[margin=2.5cm]{geometry}
\usepackage{tcolorbox} 
\tcbuselibrary{skins, breakable}
\usepackage{fancyhdr} 
\usepackage{hyperref}

% --- Definición de Cajas ---
\newtcolorbox{mydefinition}[1]{
  colback=blue!5!white, colframe=blue!75!black, fonttitle=\bfseries,
  title=Definición: #1,
  breakable
}

\newtcolorbox{mytheorem}[1]{
  colback=red!5!white, colframe=red!75!black, fonttitle=\bfseries,
  title=Teorema: #1,
  breakable
}

\newtcolorbox{mynote}[1]{
  colback=yellow!10!white, colframe=orange!80!black, fonttitle=\bfseries,
  title=Nota: #1,
  breakable
}

% --- Datos del Documento ---
\pagestyle{fancy}
\fancyhead[L]{Apuntes de Cálculo Multivariable}
\fancyhead[R]{\thepage}

\begin{document}

\title{Cálculo Multivariable: Apuntes de Clase}
\author{Levario Medina Sergio \\ Hernández Jiménez Irmin}
\date{\today}
\maketitle

\tableofcontents
\newpage

% \section{Sistemas coordenados}

El estudio del cálculo multivariable requiere entender de forma clara los espacios geométricos donde se desarrollan las funciones de varias variables. A continuación, se desarrollan los conceptos del plano y el espacio euclidiano de forma detallada.

\subsection{El Plano Cartesiano ($\mathbb{R}^2$)}

El sistema de coordenadas rectangulares en el plano se construye mediante dos rectas reales perpendiculares que se intersectan en un punto llamado \textbf{origen}.

\begin{mydefinition}{El Plano $\mathbb{R}^2$}
El conjunto $\mathbb{R}^2$ es el conjunto de todos los pares ordenados de números reales. Matemáticamente se define como el producto cartesiano $\mathbb{R} \times \mathbb{R}$:
$$ \mathbb{R}^2 = \{ (x, y) \mid x \in \mathbb{R}, y \in \mathbb{R} \} $$
A cada par ordenado $(x, y)$ se le denomina \textbf{punto} en el plano.
\end{mydefinition}

Los elementos del par ordenado tienen los siguientes nombres:
\begin{itemize}
    \item $x$: Coordenada $x$ o \textbf{abscisa}. Representa la distancia dirigida desde el eje vertical.
    \item $y$: Coordenada $y$ u \textbf{ordenada}. Representa la distancia dirigida desde el eje horizontal.
\end{itemize}

Los ejes coordenados dividen al plano en cuatro regiones denominadas \textbf{cuadrantes}, numerados en sentido antihorario comenzando por el superior derecho (donde ambas semiejes son positivas).

\begin{center}
\shorthandoff{>}
\begin{tikzpicture}[scale=1.2]
    % Ejes
    \draw[->, thick] (-3,0) -- (3,0) node[right] {$x$ (Eje de abscisas)};
    \draw[->, thick] (0,-3) -- (0,3) node[above] {$y$ (Eje de ordenadas)};
    
    % Punto P
    \coordinate (P) at (2, 1.5);
    \draw[dashed, color=gray] (2,0) node[below, black] {$a$} -- (P);
    \draw[dashed, color=gray] (0,1.5) node[left, black] {$b$} -- (P);
    \filldraw[blue] (P) circle (2pt) node[anchor=south west] {$P(a,b)$};
    
    % Etiquetas de cuadrantes
    \node[gray] at (2,2) {I $(+,+)$};
    \node[gray] at (-2,2) {II $(-,+)$};
    \node[gray] at (-2,-2) {III $(-,-)$};
    \node[gray] at (2,-2) {IV $(+,-)$};
    
    % Origen
    \node[below left] at (0,0) {$O(0,0)$};
\end{tikzpicture}
\end{center}

\subsection{El Espacio Euclidiano ($\mathbb{R}^3$)}

Para representar puntos en el espacio tridimensional, extendemos el concepto anterior añadiendo un tercer eje perpendicular a los otros dos.

\begin{mydefinition}{El Espacio $\mathbb{R}^3$}
El conjunto $\mathbb{R}^3$ es el conjunto de todas las ternas ordenadas de números reales:
$$ \mathbb{R}^3 = \{ (x, y, z) \mid x, y, z \in \mathbb{R} \} $$
\end{mydefinition}

\subsubsection{Orientación y la Regla de la Mano Derecha}
En ingeniería y física, es estándar utilizar un sistema de coordenadas \textit{dextrógiro} (mano derecha). 
\begin{mynote}{Regla de la Mano Derecha}
Si alineas los dedos de tu mano derecha en la dirección del eje $x$ positivo y los cierras hacia el eje $y$ positivo, tu pulgar apuntará en la dirección del eje $z$ positivo.
\end{mynote}

En este sistema:
\begin{itemize}
    \item El eje $z$ suele representar la altura o la cota vertical.
    \item Los tres ejes definen tres \textbf{planos coordenados}:
    \begin{itemize}
        \item \textbf{Plano $xy$}: Formado por los ejes $x$ e $y$ (donde $z=0$).
        \item \textbf{Plano $xz$}: Formado por los ejes $x$ y $z$ (donde $y=0$).
        \item \textbf{Plano $yz$}: Formado por los ejes $y$ y $z$ (donde $x=0$).
    \end{itemize}
\end{itemize}

Estos planos dividen el espacio en ocho regiones llamadas \textbf{octantes}. El primer octante es aquel donde $x, y, z > 0$.

\subsubsection{Localización de un punto en el espacio}

Para localizar un punto $P(a, b, c)$ en el espacio:
1. Avanzamos $a$ unidades sobre el eje $x$.
2. Nos movemos $b$ unidades paralelamente al eje $y$.
3. Subimos (o bajamos) $c$ unidades paralelamente al eje $z$.

\begin{center}
\shorthandoff{>}
\begin{tikzpicture}[x={(-0.5cm,-0.5cm)}, y={(1cm,0cm)}, z={(0cm,1cm)}, scale=1.5]
    % Ejes (parte negativa punteada)
    \draw[dashed, gray] (0,0,0) -- (-3,0,0); % x back
    \draw[dashed, gray] (0,0,0) -- (0,-2,0); % y back
    \draw[dashed, gray] (0,0,0) -- (0,0,-1); % z back
    
    % Ejes positivos
    \draw[->, thick] (0,0,0) -- (3,0,0) node[below left] {$x$};
    \draw[->, thick] (0,0,0) -- (0,4,0) node[right] {$y$};
    \draw[->, thick] (0,0,0) -- (0,0,3) node[above] {$z$};
    
    % Coordenadas del punto
    \def\x{2}
    \def\y{3}
    \def\z{2}
    
    % Construcción de la caja
    \draw[dashed, blue] (0,0,0) -- (\x,0,0) node[left, black] {$a$};
    \draw[dashed, blue] (\x,0,0) -- (\x,\y,0);
    \draw[dashed, blue] (0,\y,0) -- (\x,\y,0);
    \draw[dashed, blue] (0,0,0) -- (0,\y,0) node[below, black] {$b$};
    
    \draw[dashed, blue] (\x,\y,0) -- (\x,\y,\z); % Línea vertical al punto
    
    \draw[dashed, blue] (0,0,\z) -- (\x,0,\z);
    \draw[dashed, blue] (0,0,\z) -- (0,\y,\z);
    \draw[dashed, blue] (\x,0,\z) -- (\x,\y,\z);
    \draw[dashed, blue] (0,\y,\z) -- (\x,\y,\z);
    \draw[dashed, blue] (0,0,0) -- (0,0,\z) node[left, black] {$c$};
    \draw[dashed, blue] (\x,0,0) -- (\x,0,\z);
    \draw[dashed, blue] (0,\y,0) -- (0,\y,\z);

    % Punto y planos etiqueta
    \filldraw[red] (\x,\y,\z) circle (2pt) node[anchor=south west] {$P(a,b,c)$};
\end{tikzpicture}
\end{center}

% \section{Vectores}

En física e ingeniería, muchas cantidades como la temperatura o la masa pueden describirse completamente con un solo número real; a estas cantidades se les llama \textbf{escalares}. Sin embargo, cantidades como la fuerza, la velocidad o el desplazamiento requieren de una magnitud (tamaño) y una dirección espacial para estar completamente definidas. A estas cantidades las llamamos \textbf{vectores}.
Geométricamente, un vector se representa como un segmento de recta dirigido (una flecha). El punto donde comienza la flecha se llama \textit{punto inicial} (u origen) y la punta de la flecha es el \textit{punto final}.

\subsection{Vectores en $\mathbb{R}^n$}

Aunque es fácil visualizar vectores en 2 o 3 dimensiones, en matemáticas y ciencias de la computación es común trabajar con sistemas que tienen múltiples variables o grados de libertad. Para esto, extendemos la definición geométrica a una definición algebraica de $n$ dimensiones (o $n$ entradas).

\begin{mydefinition}{Vector en $\mathbb{R}^n$}
Un vector $\vec{v}$ en un espacio $n$-dimensional ($\mathbb{R}^n$) se define como una lista ordenada de $n$ números reales, llamados las \textbf{componentes} o \textbf{entradas} del vector. Se denota como:
$$ \vec{v} = (v_1, v_2, v_3, \dots, v_n) $$
Donde cada $v_i \in \mathbb{R}$ representa la magnitud del vector en la dirección del $i$-ésimo eje coordenado.
\end{mydefinition}

\subsection{Vectores libres e invarianza bajo traslación}

En el estudio del álgebra vectorial, es crucial entender que un vector está completamente definido por solo dos atributos: su \textbf{magnitud} y su \textbf{dirección}. El punto específico del espacio donde el vector comienza no altera su identidad matemática.

Esta propiedad significa que podemos tomar un vector y trasladarlo libremente a cualquier otra parte del plano o del espacio, siempre y cuando no sea girado ni estirado. 

\begin{mydefinition}{Vectores libres y equivalentes}
Dos o más vectores son \textbf{equivalentes} si tienen exactamente la misma magnitud y la misma dirección. En matemáticas y en muchas aplicaciones de ingeniería, los vectores se consideran \textbf{vectores libres}, lo que significa que un vector y cualquiera de sus traslaciones en el espacio representan exactamente al mismo objeto matemático.
\end{mydefinition}

\subsubsection{Vectores de posición vs. Vectores libres}
Aunque podemos dibujar un vector comenzando en cualquier lugar, para facilitar los cálculos algebraicos es práctica común trasladarlo para que su punto inicial coincida con el origen del sistema de coordenadas $(0,0)$ o $(0,0,0)$. Cuando un vector libre comienza en el origen, se le llama \textbf{vector de posición}.

\begin{center}
\shorthandoff{>}
\begin{tikzpicture}[scale=1]
    \draw[lightgray, very thin, dashed] (-2,-2) grid (6,4);
    
    \draw[->, gray, thick] (-2.5,0) -- (6.5,0) node[right] {$x$};
    \draw[->, gray, thick] (0,-2.5) -- (0,4.5) node[above] {$y$};
    
    \draw[->, very thick, blue] (0,0) -- (2,1.5) node[midway, above left, font=\boldmath] {$\vec{v}$};
    \filldraw[black] (0,0) circle (1.5pt) node[below right] {\footnotesize Origen};
    
    \draw[->, very thick, blue] (3,2) -- (5,3.5) node[midway, above left, font=\boldmath] {$\vec{v}$};
    \draw[->, very thick, blue] (-1,-1) -- (1,0.5) node[midway, above left, font=\boldmath] {$\vec{v}$};
    \draw[->, very thick, blue] (4,-1) -- (6,0.5) node[midway, above left, font=\boldmath] {$\vec{v}$};
    
    \filldraw[black] (3,2) circle (1.5pt);
    \filldraw[black] (-1,-1) circle (1.5pt);
    \filldraw[black] (4,-1) circle (1.5pt);
    
    \node[fill=white, draw=black, rounded corners, inner sep=4pt] at (2, -1.8) 
    {\footnotesize Todos estos representan al mismo vector $\vec{v}$};
\end{tikzpicture}
\end{center}

\subsection{Multiplicación por un escalar}

La multiplicación escalar es la operación de multiplicar un vector por un número real (escalar) $\kappa$. Geométricamente, esto altera la longitud (magnitud) del vector. 
\begin{itemize}
    \item Si $\kappa > 1$, el vector se \textit{estira} o dilata.
    \item Si $0 < \kappa < 1$, el vector se \textit{contrae}.
    \item Si $\kappa < 0$, el vector cambia a la \textit{dirección opuesta} (invierte su sentido).
\end{itemize}
Sin embargo, el vector resultante siempre es paralelo al vector original (se mantiene sobre la misma línea de acción).

\begin{center}
\shorthandoff{>}
\begin{tikzpicture}[scale=1.2]
    % Cuadrícula y ejes de fondo para referencia
    \draw[lightgray, very thin] (-4,-2) grid (5,3);
    \draw[->, thick, gray] (-4.5,0) -- (5.5,0) node[right] {$x$};
    \draw[->, thick, gray] (0,-2.5) -- (0,3.5) node[above] {$y$};

    % Vector original v
    \draw[->, very thick, blue] (0,0) -- (2,1) node[above left, font=\boldmath] {$\vec{v}$};
    
    % Escalar 2: Estiramiento (dibujado ligeramente desplazado para que se vea)
    \draw[->, thick, red, yshift=-0.1cm] (0,0) -- (4,2) node[below right, font=\boldmath] {$2\vec{v}$};
    
    % Escalar 0.5: Contracción (dibujado ligeramente desplazado hacia arriba)
    \draw[->, thick, green!60!black, yshift=0.1cm] (0,0) -- (1,0.5) node[above left, font=\boldmath] {$0.5\vec{v}$};
    
    % Escalar -1: Inversión de sentido
    \draw[->, very thick, orange] (0,0) -- (-2,-1) node[below right, font=\boldmath] {$-\vec{v}$};
    
\end{tikzpicture}
\end{center}

\subsection{Suma de vectores en $\mathbb{R}^2$}

Para sumar dos vectores $\vec{u}$ y $\vec{v}$ geométricamente, se utiliza la \textbf{Regla del Paralelogramo} o el método de "punta a cola". Si colocamos el punto inicial de $\vec{v}$ en el punto final de $\vec{u}$, el vector suma $\vec{u} + \vec{v}$ es la flecha que va desde el origen de $\vec{u}$ hasta la punta de $\vec{v}$.

\begin{center}
\shorthandoff{>}
\begin{tikzpicture}[scale=1]
    % Ejes
    \draw[->, gray] (-1,0) -- (6,0) node[right] {$x$};
    \draw[->, gray] (0,-1) -- (0,5) node[above] {$y$};

    % Vectores u y v
    \draw[->, very thick, blue] (0,0) -- (4,1) node[midway, below right] {$\vec{u}$};
    \draw[->, very thick, red] (0,0) -- (1,3) node[midway, left] {$\vec{v}$};

    % Proyecciones para el paralelogramo
    \draw[->, dashed, red] (4,1) -- (5,4) node[midway, right] {$\vec{v}$};
    \draw[->, dashed, blue] (1,3) -- (5,4) node[midway, above] {$\vec{u}$};

    % Vector Suma
    \draw[->, very thick, purple] (0,0) -- (5,4) node[above right] {$\vec{u} + \vec{v}$};
\end{tikzpicture}
\end{center}

\subsection{Sustracción de vectores}

La resta de dos vectores se define algebraicamente a partir de la suma y la multiplicación escalar. Restar $\vec{v}$ de $\vec{u}$ es lo mismo que sumarle a $\vec{u}$ el negativo de $\vec{v}$:

\[ \vec{u} - \vec{v} = \vec{u} + (-\vec{v}) \] 

Geométricamente, si graficamos $\vec{u}$ y $\vec{v}$ desde el mismo origen, el vector $\vec{u} - \vec{v}$ es el vector que va desde la punta de $\vec{v}$  hasta la punta de $\vec{u}$.

\begin{center}
\shorthandoff{>}
\begin{tikzpicture}[scale=1]
    \draw[->, gray] (-2,0) -- (5,0) node[right] {$x$};
    \draw[->, gray] (0,-3) -- (0,4) node[above] {$y$};

    \draw[->, very thick, blue] (0,0) -- (4,2) node[above right] {$\vec{u}$};
    \draw[->, very thick, red] (0,0) -- (1,3) node[above left] {$\vec{v}$};

    \draw[->, dashed, orange] (0,0) -- (-1,-3) node[below left] {$-\vec{v}$};
    
    \draw[->, dashed, blue] (-1,-3) -- (3,-1);
    \draw[->, dashed, orange] (4,2) -- (3,-1);
    
    \draw[->, very thick, purple] (0,0) -- (3,-1) node[below right] {$\vec{u} - \vec{v}$};

    \draw[->, very thick, purple, dashed] (1,3) -- (4,2) node[midway, above right] {$\vec{u} - \vec{v}$};
\end{tikzpicture}
\end{center}

\subsection{Generalización de operaciones en $\mathbb{R}^n$}

Las operaciones que visualizamos en $\mathbb{R}^2$ se definen algebraicamente componente a componente para cualquier vector en $\mathbb{R}^n$.

Sean $\vec{u} = (u_1, u_2, \dots, u_n)$ y $\vec{v} = (v_1, v_2, \dots, v_n)$ dos vectores en $\mathbb{R}^n$, y sea $\kappa \in \mathbb{R}$ un escalar. Entonces, las operaciones fundamentales se definen como:

\begin{mydefinition}{Operaciones vectoriales en $\mathbb{R}^n$}
\textbf{1. Suma vectorial:}
\[\vec{u} + \vec{v} = (u_1 + v_1, u_2 + v_2, \dots, u_n + v_n)\]

\textbf{2. Multiplicación por un escalar:}
\[\kappa\vec{u} = (\kappa u_1, \kappa u_2, \dots, \kappa u_n)\]

\textbf{3. Resta vectorial:}
\[\vec{u} - \vec{v} = (u_1 - v_1, u_2 - v_2, \dots, u_n - v_n)\]
\end{mydefinition}

\subsection{Propiedades fundamentales de los vectores}

Las operaciones que definimos anteriormente siguen axiomas específicos. Estas propiedades son la base para manipular ecuaciones vectoriales en el cálculo multivariable y el álgebra lineal.

Sean $\vec{u}$, $\vec{v}$ y $\vec{w}$ vectores en $\mathbb{R}^n$, y sean $\kappa$ y $\lambda$ escalares (números reales). Se cumplen las siguientes propiedades:

\begin{mytheorem}{Propiedades del Álgebra Vectorial}
\begin{enumerate}
    \item \textbf{Conmutatividad de la suma:} $\vec{u} + \vec{v} = \vec{v} + \vec{u}$
    \item \textbf{Asociatividad de la suma:} $(\vec{u} + \vec{v}) + \vec{w} = \vec{u} + (\vec{v} + \vec{w})$
    \item \textbf{Elemento neutro aditivo:} Existe un vector cero $\vec{0} = (0,0,\dots,0)$ tal que $\vec{u} + \vec{0} = \vec{u}$
    \item \textbf{Inverso aditivo:} Para cada $\vec{u}$, existe un vector $-\vec{u}$ tal que $\vec{u} + (-\vec{u}) = \vec{0}$
    \item \textbf{Distributividad del escalar sobre suma vectorial:} $\kappa(\vec{u} + \vec{v}) = \kappa\vec{u} + \kappa\vec{v}$
    \item \textbf{Distributividad del escalar sobre suma escalar:} $(\kappa + \lambda)\vec{u} = \kappa\vec{u} + \lambda\vec{u}$
    \item \textbf{Asociatividad de la multiplicación escalar:} $\kappa(\lambda\vec{u}) = (\kappa\lambda)\vec{u}$
    \item \textbf{Elemento neutro multiplicativo:} $1\vec{u} = \vec{u}$
\end{enumerate}
\end{mytheorem}

\subsection{Magnitud de un vector}

La \textbf{magnitud}, \textbf{longitud} o \textbf{norma} de un vector es una medida escalar que representa el ``tamaño'' del vector. Físicamente, si el vector representa un desplazamiento, su magnitud es la distancia rectilínea entre su punto inicial y su punto final.

\subsubsection{La magnitud en $\mathbb{R}^2$ y el teorema de pitágoras}

En el plano bidimensional ($\mathbb{R}^2$), cualquier vector $\vec{v} = (v_1, v_2)$ puede visualizarse como la hipotenusa de un triángulo rectángulo, donde los catetos están formados por las proyecciones del vector sobre los ejes $x$ e $y$.

Por lo tanto, la magnitud del vector, denotada como $\|\vec{v}\|$, se calcula aplicando directamente el \textbf{teorema de pitágoras}:

\begin{center}
\shorthandoff{>}
\begin{tikzpicture}[scale=1.5]
    \draw[->, gray, thick] (-0.5,0) -- (4,0) node[right] {$x$};
    \draw[->, gray, thick] (0,-0.5) -- (0,3) node[above] {$y$};

    \draw[dashed, blue, thick] (3,0) -- (3,2);
    \draw[dashed, red, thick] (0,0) -- (3,0);
    
    \draw (2.8,0) -- (2.8,0.2) -- (3,0.2);

    \draw[->, very thick, black] (0,0) -- (3,2) node[midway, above left] {$\|\vec{v}\| = \sqrt{v_1^2 + v_2^2}$};
    
    \node[below, red, font=\boldmath] at (1.5,0) {$v_1$ (cateto adyacente)};
    \node[right, blue, font=\boldmath] at (3,1) {$v_2$ (cateto opuesto)};
    \filldraw[black] (3,2) circle (1.5pt) node[above right] {$\vec{v} = (v_1, v_2)$};
\end{tikzpicture}
\end{center}

\subsubsection{Extensión a $\mathbb{R}^n$}

Este mismo principio geométrico se extiende a dimensiones superiores utilizando la distancia euclidiana. 

\begin{mydefinition}{Norma de un vector en $\mathbb{R}^n$}
Sea $\vec{v} = (v_1, v_2, \dots, v_n)$ un vector en el espacio $\mathbb{R}^n$. La magnitud o norma de $\vec{v}$ se define como la raíz cuadrada de la suma de los cuadrados de sus componentes:
\[ \|\vec{v}\| = \sqrt{v_1^2 + v_2^2 + \dots + v_n^2} \]
\end{mydefinition}

\subsection{Distancia entre dos vectores}

En muchas aplicaciones de ingeniería, es necesario calcular la distancia en línea recta entre dos puntos en el espacio. Si consideramos que estos dos puntos están definidos por los extremos de dos vectores que parten del origen $\vec{u}$ y $\vec{v}$, podemos encontrar esta distancia utilizando las herramientas que ya se definieron: la resta de vectores y la magnitud.

Geométricamente, se sabe que la resta vectorial $\vec{v} - \vec{u}$ produce un nuevo vector que va desde la punta de $\vec{u}$ hasta la punta de $\vec{v}$, representando el desplazamiento exacto entre ambos puntos. 

Por lo tanto, la distancia física entre los extremos de $\vec{u}$ y $\vec{v}$ es simplemente la \textbf{magnitud de su vector diferencia}.

\begin{mydefinition}{Distancia entre vectores en $\mathbb{R}^3$}
Sean $\vec{u} = (u_1, u_2, u_3)$ y $\vec{v} = (v_1, v_2, v_3)$ dos vectores en el espacio tridimensional. La distancia $d$ entre sus puntos finales se define como la norma del vector diferencia $\vec{u} - \vec{v}$:
\[ d(\vec{u}, \vec{v}) = \|\vec{u} - \vec{v}\| \]
    
Dado que la resta de vectores se realiza componente a componente, el vector diferencia es $(u_1 - v_1, u_2 - v_2, u_3 - v_3)$. Al aplicarle la definición de magnitud, se obtiene así por la distancia euclidiana:
\[ d(\vec{u}, \vec{v}) = \sqrt{(u_1 - v_1)^2 + (u_2 - v_2)^2 + (u_3 - v_3)^2} \]

\textit{Nota:} El orden de la resta no importa para el cálculo de la distancia, ya que $\|\vec{u} - \vec{v}\| = \|\vec{v} - \vec{u}\|$. Al elevar al cuadrado las diferencias de las componentes, cualquier signo negativo se vuelve positivo.
\end{mydefinition}

\begin{center}
\shorthandoff{>}
\begin{tikzpicture}[x={(-0.5cm,-0.5cm)}, y={(1cm,0cm)}, z={(0cm,1cm)}, scale=1.5]
    % Ejes
    \draw[->, gray] (0,0,0) -- (3,0,0) node[below left] {$x$};
    \draw[->, gray] (0,0,0) -- (0,4,0) node[right] {$y$};
    \draw[->, gray] (0,0,0) -- (0,0,3) node[above] {$z$};
    
    % Definición de las coordenadas de los vectores
    \coordinate (U) at (2, 1, 2);
    \coordinate (V) at (1, 3, 1);
    
    % Vectores u y v desde el origen
    \draw[->, very thick, blue] (0,0,0) -- (U) node[midway, above left] {$\vec{u}$};
    \draw[->, very thick, red] (0,0,0) -- (V) node[midway, below right] {$\vec{v}$};
    
    % Vector diferencia (representa la distancia)
    \draw[->, very thick, purple] (U) -- (V) node[midway, above right] {$\vec{v} - \vec{u}$};
    
    % Puntos en los extremos para dar énfasis
    \filldraw[black] (U) circle (1.2pt);
    \filldraw[black] (V) circle (1.2pt);
    
    % Cuadro explicativo
    \node[fill=white, draw=black, rounded corners, inner sep=4pt] at (2.5, 3.5, 3) 
    {\footnotesize Distancia $= \|\vec{v} - \vec{u}\|$};
\end{tikzpicture}
\end{center}

\subsection{Vectores unitarios}

En muchas aplicaciones de ingeniería, a menudo se requiere únicamente la dirección y el sentido en el que actúa un fenómeno (como la dirección del flujo de un fluido o la línea de acción de una fuerza), independientemente de su magnitud. Para aislar esta información, se utilizan los vectores unitarios.

\begin{mydefinition}{Vector unitario}
Un vector unitario es aquel cuya magnitud (o longitud) es exactamente igual a 1. Si se denota la magnitud de un vector $\vec{v}$ como $\|\vec{v}\|$, entonces $\vec{v}$ es unitario si:
\[ \|\vec{v}\| = 1 \]
Para distinguirlos del resto, a los vectores unitarios se les suele denotar con un acento circunflejo en lugar de una flecha, por ejemplo, $\hat{u}$.
\end{mydefinition}

\subsubsection{Normalización de un vector}

Cualquier vector $\vec{v}$ que no sea el vector cero puede transformarse en un vector unitario que apunte exactamente en su misma dirección. A este proceso geométrico se le llama \textbf{normalización}. 

Para normalizar un vector, simplemente se multiplica por el recíproco de su propia magnitud:
\[ \hat{u} = \frac{1}{\|\vec{v}\|} \vec{v} = \frac{\vec{v}}{\|\vec{v}\|} \]
De esta manera, el vector resultante $\hat{u}$ conserva la dirección original, pero su tamaño se ha escalado para que sea exactamente 1.

\subsubsection{Vectores unitarios canónicos}

En los sistemas de coordenadas rectangulares, definimos vectores unitarios especiales que apuntan a lo largo de las direcciones positivas de los ejes coordenados. Estos son los vectores base fundamentales, ya que nos permiten expresar cualquier vector en el espacio como una combinación de ellos de forma sencilla (aunque existen múltiples bases para un mismo espacio, estos vectores se les considera canónicos ya que son los más comúnmente utilizados).

En el espacio $\mathbb{R}^3$, estos vectores son:
\begin{itemize}
    \item $\hat{\imath} = (1, 0, 0)$ en la dirección del eje $x$.
    \item $\hat{\jmath} = (0, 1, 0)$ en la dirección del eje $y$;
    \item $\hat{k} = (0, 0, 1)$ en la dirección del eje $z$.
\end{itemize}

Gracias a esto, cualquier vector $\vec{v} = (v_1, v_2, v_3)$ se puede escribir algebraicamente separando sus componentes direccionales:
\[ \vec{v} = v_1\hat{\imath} + v_2\hat{\jmath} + v_3\hat{k} \]

\begin{center}
\shorthandoff{>}
\begin{tikzpicture}[x={(-0.5cm,-0.5cm)}, y={(1cm,0cm)}, z={(0cm,1cm)}, scale=2]
    \draw[->, gray] (0,0,0) -- (2,0,0) node[below left] {$x$};
    \draw[->, gray] (0,0,0) -- (0,2,0) node[right] {$y$};
    \draw[->, gray] (0,0,0) -- (0,0,2) node[above] {$z$};
    
    \draw[->, very thick, red] (0,0,0) -- (1,0,0) node[midway, below right, font=\boldmath] {$\hat{\imath}$};
    \draw[->, very thick, green!60!black] (0,0,0) -- (0,1,0) node[midway, below right, font=\boldmath] {$\hat{\jmath}$};
    \draw[->, very thick, blue] (0,0,0) -- (0,0,1) node[midway, left, font=\boldmath] {$\hat{k}$};
    

    \draw[dashed, gray] (1,0,0) -- (1,-0.2,0) node[below=2pt, black] {\footnotesize 1};
    \draw[dashed, gray] (0,1,0) -- (0,1,-0.2) node[below=2pt, black] {\footnotesize 1};
    \draw[dashed, gray] (0,0,1) -- (-0.2,0,1) node[left=2pt, black] {\footnotesize 1};
\end{tikzpicture}
\end{center}

\subsection{Ejercicios}
Realizar los cálculos en los ejercicios del 1 al 3.

\begin{enumerate}
    \item $ (8a, -2b, 13c) = (52, 12, 11) + \frac{1}{2}(?, ?, ?)$
    \item $ (-21, 23) - (?, 6) = (-25, ?) $
    \item $ (2,3,5) - 4 \hat{\imath} + 3\hat{\jmath} = (?,?,?) $
    \item ¿Qué restricciones se deben de tener sobre $x$, $y$ y $z$ de modo que la terna $(x,y,z)$ represente un punto sobre el eje $y$? ¿Y sobre el eje $z$? ¿En el plano $xz$? ¿En el plano $yz$?
    \item Trazar los vectores $\vec{v} = (2,3,-6)$ y $\vec{w} = (-1,1,1)$. En esa figura trazar $-\vec{v}$, $\vec{v} + \vec{w}$, $2\vec{v}$ y $\vec{v} - \vec{w}$.
    \item Dos personas jalan horizontalmente de cuerdas atadas a un poste; el ángulo entre las cuerdas es de $60^\circ$. A jala con una fuerza de $150$ lb y B jala con una fuerza de $100$ lb.
\begin{enumerate}
    \item La fuerza resultante es la suma vectorial de las dos fuerzas en un sistema coordenado escogido de manera conveniente. Trazar una figura a escala que represente graficamente a las tres fuerzas.
    \item Usando trigonometría, determinar fórmulas para las componentes vectoriales de las dos fuerzas en un sistema coordenado escogido de manera conveniente. Efectuar la suma algebraica y hallar el ángulo que la fuerza resultante forma con la cuerda de A.
\end{enumerate}
    \item La velocidad $\vec{v_1}$ del viento es de $ 40 \frac{mi}{h} $ de este a oeste, mientras que un aeroplano viaja con velocidad en el aire $ \vec{v_2} = 100 \frac{mi}{h} $ hacia el norte. La rapidez del aeroplano respecto a la tierra es la magnitud de la suma vectorial $\vec{v_1} + \vec{v_2}$.
\begin{enumerate}
    \item Hallar $\vec{v_1} + \vec{v_2}$.
    \item Trazar una figura a escala.
\end{enumerate}
    \item Dado el vector $\vec{v} = (3, -4)$, calcular su magnitud $\|\vec{v}\|$ y encontrar el vector unitario $\hat{u}$ que apunta en la misma dirección que $\vec{v}$. Verificar que $\|\hat{u}\| = 1$.
    \item Calcular la distancia entre los puntos $P = (1, -2, 4)$ y $Q = (3, 1, -1)$ en $\mathbb{R}^3$. Expresar el vector $\overrightarrow{PQ}$ en términos de los vectores unitarios canónicos $\hat{\imath}$, $\hat{\jmath}$ y $\hat{k}$.
    \item Sean $\vec{u} = (2, -1, 3)$ y $\vec{v} = (-1, 4, 2)$. Calcular cada una de las siguientes expresiones:
\begin{enumerate}
    \item $3\vec{u} - 2\vec{v}$
    \item $\|\vec{u} + \vec{v}\|$
    \item $\|\vec{u}\| - \|\vec{v}\|$
\end{enumerate}
¿Por qué en general $\|\vec{u} + \vec{v}\| \neq \|\vec{u}\| + \|\vec{v}\|$?
    \item Un dron parte del punto $A = (0, 0, 0)$ y se desplaza al punto $B = (3, 4, 0)$, luego sube verticalmente hasta el punto $C = (3, 4, 5)$.
\begin{enumerate}
    \item Hallar los vectores de desplazamiento $\overrightarrow{AB}$ y $\overrightarrow{BC}$.
    \item Determinar el desplazamiento total $\overrightarrow{AC}$ como suma vectorial.
    \item Calcular la distancia total recorrida y la distancia en línea recta entre $A$ y $C$. ¿Son iguales? ¿Por qué?
\end{enumerate}
    \item Encontrar todos los valores del escalar $t \in \mathbb{R}$ para los cuales el vector $\vec{w} = (t, 2t-1, 3)$ tenga magnitud igual a $\sqrt{14}$.
    \item Demostrar algebraicamente, usando la definición componente a componente, las siguientes propiedades del álgebra vectorial para vectores en $\mathbb{R}^n$:
\begin{enumerate}
    \item $\kappa(\vec{u} + \vec{v}) = \kappa\vec{u} + \kappa\vec{v}$ \quad (distributividad del escalar sobre la suma vectorial)
    \item $(\kappa + \lambda)\vec{u} = \kappa\vec{u} + \lambda\vec{u}$ \quad (distributividad del escalar sobre la suma escalar)
\end{enumerate}
\end{enumerate}

% \section{Producto punto}

Hasta ahora se ha explorado la multiplicación escalar, donde al multiplicar un vector por un valor escalar se obtiene como resultado otro vector. Sin embargo, en el álgebra vectorial es posible multiplicar dos vectores entre sí. 

Existen dos formas principales de multiplicar vectores en cálculo multivariable. La primera de ellas es el \textbf{producto punto}, también conocido como \textbf{producto escalar} o \textbf{producto interno}. Su principal característica es que la operación toma dos vectores y devuelve como resultado un valor escalar, no un vector.

En aplicaciones de ingeniería y física, el producto punto es una herramienta indispensable. Permite determinar de forma cuantitativa qué tanto "apuntan" dos vectores en la misma dirección, calcular el trabajo mecánico realizado por una fuerza a lo largo de un desplazamiento, o encontrar las proyecciones de un vector sobre otro.

\subsubsection{Definición Algebraica en $\mathbb{R}^n$}

Desde el punto de vista algebraico, el producto punto es una operación muy directa. Se calcula simplemente multiplicando las componentes correspondientes de ambos vectores y luego sumando esos resultados parciales.

\begin{mydefinition}{Producto Punto en $\mathbb{R}^n$}
Sean $\vec{u} = (u_1, u_2, \dots, u_n)$ y $\vec{v} = (v_1, v_2, \dots, v_n)$ dos vectores en el espacio euclidiano $\mathbb{R}^n$. El producto punto de $\vec{u}$ y $\vec{v}$, denotado por el símbolo de un punto centrado $\vec{u} \cdot \vec{v}$, se define como el escalar resultante de:
\[ \vec{u} \cdot \vec{v} = u_1v_1 + u_2v_2 + \dots + u_nv_n \]

De forma compacta, utilizando la notación de sumatoria, se expresa como:
\[ \vec{u} \cdot \vec{v} = \sum_{i=1}^{n} u_i v_i \]
\end{mydefinition}

\subsubsection{Propiedades del Producto Punto}

Al igual que las operaciones aritméticas básicas, el producto punto obedece a ciertas leyes algebraicas que nos permiten manipular y simplificar ecuaciones vectoriales. Sean $\vec{u}$, $\vec{v}$ y $\vec{w}$ vectores en $\mathbb{R}^n$, y sea $c$ un escalar. Se cumplen las siguientes propiedades:

\begin{mytheorem}{Propiedades Algebraicas del Producto Punto}
\begin{enumerate}
    \item \textbf{Conmutatividad:} $\vec{u} \cdot \vec{v} = \vec{v} \cdot \vec{u}$
    \item \textbf{Distributividad sobre la suma:} $\vec{u} \cdot (\vec{v} + \vec{w}) = \vec{u} \cdot \vec{v} + \vec{u} \cdot \vec{w}$
    \item \textbf{Asociatividad con un escalar:} $c(\vec{u} \cdot \vec{v}) = (c\vec{u}) \cdot \vec{v} = \vec{u} \cdot (c\vec{v})$
    \item \textbf{Propiedad del vector cero:} $\vec{0} \cdot \vec{v} = 0$
    \item \textbf{Relación con la magnitud:} El producto punto de un vector consigo mismo es igual al cuadrado de su magnitud:
    \[ \vec{v} \cdot \vec{v} = \|\vec{v}\|^2 \quad \text{o bien} \quad \|\vec{v}\| = \sqrt{\vec{v} \cdot \vec{v}} \]
\end{enumerate}
\end{mytheorem}

\subsection{Ángulo entre Dos Vectores}

Aunque definimos el producto punto de forma algebraica multiplicando componentes, su verdadera utilidad en ingeniería y física radica en su interpretación geométrica. El producto punto nos revela la relación angular entre dos vectores.

Si colocamos dos vectores $\vec{u}$ y $\vec{v}$ de modo que compartan el mismo origen, estos formarán un ángulo $\theta$ entre ellos, donde $0 \leq \theta \leq \pi$.

\begin{mytheorem}{Ángulo entre vectores}
Si $\theta$ es el ángulo entre dos vectores no nulos $\vec{u}$ y $\vec{v}$, entonces el producto punto está dado geométricamente por:
\[ \vec{u} \cdot \vec{v} = \|\vec{u}\| \|\vec{v}\| \cos(\theta) \]

De esta fórmula se puede despejar el ángulo $\theta$ para calcularlo a partir de las componentes de los vectores:
\[ \cos(\theta) = \frac{\vec{u} \cdot \vec{v}}{\|\vec{u}\| \|\vec{v}\|} \]
\end{mytheorem}

\subsubsection{Visualización y Demostración}

Para demostrar este teorema, se utilizará la geometría del triángulo formado por los vectores y se aplicará la \textbf{Ley de los Cosenos}.

\begin{center}
\shorthandoff{>}
\begin{tikzpicture}[scale=1.5]
    % Vectores
    \draw[->, very thick, blue] (0,0) -- (4,0) node[midway, below] {$\vec{u}$};
    \draw[->, very thick, red] (0,0) -- (2.5,2.5) node[midway, above left] {$\vec{v}$};
    
    % Vector diferencia cerrando el triángulo
    \draw[->, very thick, purple] (4,0) -- (2.5,2.5) node[midway, above right] {$\vec{v} - \vec{u}$};
    
    % Ángulo theta
    \draw[thick] (0.8,0) arc (0:45:0.8);
    \node at (1.1,0.4) {$\theta$};
    
    % Puntos en los extremos
    \filldraw[black] (0,0) circle (1.5pt) node[left] {Origen};
\end{tikzpicture}
\end{center}

\begin{proof}[Demostración mediante la Ley de los Cosenos]

Considerese un triángulo cuyos lados están formados por los vectores $\vec{u}$, $\vec{v}$ y el vector diferencia $\vec{v} - \vec{u}$. Las longitudes de los lados de este triángulo son exactamente las magnitudes de estos vectores: $\|\vec{u}\|$, $\|\vec{v}\|$ y $\|\vec{v} - \vec{u}\|$.

Aplicando la Ley de los Cosenos a este triángulo para relacionar los lados con el ángulo $\theta$, se tiene:
\[ \|\vec{v} - \vec{u}\|^2 = \|\vec{u}\|^2 + \|\vec{v}\|^2 - 2\|\vec{u}\|\|\vec{v}\|\cos(\theta) \]

Ahora, se reescribe el lado izquierdo de la ecuación utilizando la propiedad que relaciona la magnitud con el producto punto ($\|\vec{x}\|^2 = \vec{x} \cdot \vec{x}$):
\[ \|\vec{v} - \vec{u}\|^2 = (\vec{v} - \vec{u}) \cdot (\vec{v} - \vec{u}) \]

Aplicando la propiedad distributiva del producto punto, se desarrolla este binomio:
\[ (\vec{v} - \vec{u}) \cdot (\vec{v} - \vec{u}) = \vec{v} \cdot \vec{v} - \vec{v} \cdot \vec{u} - \vec{u} \cdot \vec{v} + \vec{u} \cdot \vec{u} \]

Usando la propiedad conmutativa ($\vec{u} \cdot \vec{v} = \vec{v} \cdot \vec{u}$) y convirtiendo de vuelta a magnitudes:
\[ \|\vec{v} - \vec{u}\|^2 = \|\vec{v}\|^2 - 2(\vec{u} \cdot \vec{v}) + \|\vec{u}\|^2 \]

Se iguala esta expresión algebraica con la ecuación geométrica obtenida de la Ley de los Cosenos:
\[ \|\vec{v}\|^2 - 2(\vec{u} \cdot \vec{v}) + \|\vec{u}\|^2 = \|\vec{u}\|^2 + \|\vec{v}\|^2 - 2\|\vec{u}\|\|\vec{v}\|\cos(\theta) \]

Restando $\|\vec{u}\|^2$ y $\|\vec{v}\|^2$ de ambos lados, la ecuación se simplifica notablemente a:
\[ -2(\vec{u} \cdot \vec{v}) = -2\|\vec{u}\|\|\vec{v}\|\cos(\theta) \]

Finalmente, dividiendo ambos lados entre $-2$, se llega a que:
\[ \vec{u} \cdot \vec{v} = \|\vec{u}\|\|\vec{v}\|\cos(\theta) \]
\end{proof}

\subsection{Vectores ortogonales}

Una de las consecuencias más útiles del teorema del ángulo entre dos vectores es que nos proporciona una prueba algebraica para determinar si dos vectores son perpendiculares entre sí.

Si dos vectores no nulos $\vec{u}$ y $\vec{v}$ son ortogonales, el ángulo entre ellos es exactamente $\theta = \frac{\pi}{2}$ radianes. 

Sustituyendo este valor en nuestra fórmula geométrica:
\[ \vec{u} \cdot \vec{v} = \|\vec{u}\|\|\vec{v}\|\cos\left(\frac{\pi}{2}\right) \]

Dado que $\cos(\frac{\pi}{2}) = 0$, toda la expresión de la derecha se anula.

\begin{mytheorem}{Condición de ortogonalidad}
Dos vectores no nulos $\vec{u}$ y $\vec{v}$ son ortogonales si y solo si su producto punto es igual a cero:
\[ \vec{u} \cdot \vec{v} = 0 \]
Por convención matemática, se considera que el vector cero ($\vec{0}$) es ortogonal a cualquier otro vector, ya que $\vec{0} \cdot \vec{v} = 0$ para todo $\vec{v}$.
\end{mytheorem}

\subsection{Cosenos directores}

En el espacio $\mathbb{R}^3$, a menudo es necesario describir la dirección exacta de un vector en relación con el sistema de coordenadas. Para ello, se utilizan los \textbf{ángulos directores} $\alpha$, $\beta$ y $\gamma$, que son los ángulos que el vector forma con los ejes positivos $x$, $y$ y $z$, respectivamente.

Si se toma un vector $\vec{v} = (v_1, v_2, v_3)$ y se calcula el producto punto de este vector con cada uno de los vectores unitarios canónicos ($\hat{i}, \hat{j}, \hat{k}$), es posible encontrar los cosenos de estos ángulos, conocidos como \textbf{cosenos directores}.

Por ejemplo, para el ángulo $\alpha$ con el eje $x$:
\[ \vec{v} \cdot \hat{\imath} = \|\vec{v}\|\|\hat{\imath}\|\cos(\alpha) \]
Sabiendo que $\vec{v} \cdot \hat{\imath} = (v_1)(1) + (v_2)(0) + (v_3)(0) = v_1$ y que $\|\hat{\imath}\| = 1$, se obtiene:
\[ v_1 = \|\vec{v}\|\cos(\alpha) \implies \cos(\alpha) = \frac{v_1}{\|\vec{v}\|} \]

\begin{mydefinition}{Cosenos directores}
Para cualquier vector no nulo $\vec{v} = (v_1, v_2, v_3)$, sus cosenos directores son:
\[ \cos(\alpha) = \frac{v_1}{\|\vec{v}\|}, \quad \cos(\beta) = \frac{v_2}{\|\vec{v}\|}, \quad \cos(\gamma) = \frac{v_3}{\|\vec{v}\|} \]

Una propiedad importante es que la suma de los cuadrados de los cosenos directores siempre es igual a 1:
\[ \cos^2(\alpha) + \cos^2(\beta) + \cos^2(\gamma) = 1 \]
\end{mydefinition}

Un hecho interesante, es que los cosenos directores son exactamente las componentes del \textbf{vector unitario} que apunta en la dirección de $\vec{v}$. Es decir:
\[ \frac{\vec{v}}{\|\vec{v}\|} = (\cos\alpha, \cos\beta, \cos\gamma) \]

\begin{center}
\shorthandoff{>}
\begin{tikzpicture}[x={(-0.5cm,-0.5cm)}, y={(1cm,0cm)}, z={(0cm,1cm)}, scale=2]
    % Ejes
    \draw[->, gray] (0,0,0) -- (2,0,0) node[below left] {$x$};
    \draw[->, gray] (0,0,0) -- (0,2,0) node[right] {$y$};
    \draw[->, gray] (0,0,0) -- (0,0,2) node[above] {$z$};
    
    % Vector v
    \coordinate (V) at (1.2, 1.5, 1.5);
    \draw[->, very thick, blue] (0,0,0) -- (V) node[above right] {$\vec{v}$};
    
    % Arcos para los ángulos directores
    % Alfa (con el eje x)
    \draw[thick, red] (0.5,0,0) to[out=90,in=-120] (0.4, 0.5, 0.5);
    \node[red, below=2pt] at (0.6, 0.2, 0.2) {$\alpha$};
    
    % Beta (con el eje y)
    \draw[thick, green!60!black] (0,0.6,0) to[out=90,in=-30] (0.3, 0.5, 0.4);
    \node[green!60!black, right=2pt] at (0.2, 0.6, 0.2) {$\beta$};
    
    % Gamma (con el eje z)
    \draw[thick, orange] (0,0,0.6) to[out=0,in=150] (0.3, 0.4, 0.4);
    \node[orange, left=2pt] at (0.1, 0.2, 0.7) {$\gamma$};
\end{tikzpicture}
\end{center}


\subsection{Proyección Vectorial}

Una aplicación práctica del producto punto es averiguar qué parte de un vector actúa exactamente en la misma dirección que otro vector. Imagina que tienes una fuente de luz que brilla directamente sobre un vector $\vec{u}$, perpendicular a un segundo vector $\vec{v}$. La "sombra" que $\vec{u}$ proyecta sobre la línea de acción de $\vec{v}$ es lo que se conoce como \textbf{proyección vectorial}.

\begin{center}
\shorthandoff{>}
\begin{tikzpicture}[scale=1.5]
    % Línea de acción base
    \draw[->, thick, gray] (0,0) -- (5,0) node[right, black] {Línea de acción de $\vec{v}$};
    \draw[->, very thick, red] (0,0) -- (4,0) node[midway, below] {$\vec{v}$};
    
    % Vector u a proyectar
    \draw[->, very thick, blue] (0,0) -- (3,2) node[midway, above left] {$\vec{u}$};
    
    % Línea de proyección (perpendicular) indicando la "sombra"
    \draw[dashed, thick, blue] (3,2) -- (3,0);
    
    % Símbolo de ángulo recto
    \draw (2.8,0) -- (2.8,0.2) -- (3,0.2);
    
    % Ángulo theta
    \draw[thick] (0.8,0) arc (0:33.69:0.8);
    \node at (1,0.3) {$\theta$};
    
    % Vector Proyección (la "sombra")
    \draw[->, very thick, purple, yshift=-0.05cm] (0,0) -- (3,0) node[midway, below=3pt] {$\text{proj}_{\vec{v}}\vec{u}$};
    
    % Punto de origen
    \filldraw[black] (0,0) circle (1.5pt) node[below left] {Origen};
\end{tikzpicture}
\end{center}

Para encontrar la fórmula matemática de este nuevo vector proyección, $\text{proj}_{\vec{v}}\vec{u}$, debe construirse paso a paso utilizando su geometría. Todo vector requiere dos cosas: una dirección y una magnitud.

\textbf{1. La Dirección:} 
Se sabe que la proyección debe estar sobre la misma línea que el vector $\vec{v}$. Por lo tanto, compartirá su mismo vector unitario direccional.
\[ \text{Dirección} = \frac{\vec{v}}{\|\vec{v}\|} \]

\textbf{2. La Magnitud:} 
Si se observa el triángulo rectángulo formado en el gráfico anterior, la hipotenusa es la magnitud de nuestro vector original ($\|\vec{u}\|$). La longitud de la sombra (nuestra proyección) corresponde al cateto adyacente al ángulo $\theta$. Usando trigonometría básica ($\cos\theta = \frac{\text{adyacente}}{\text{hipotenusa}}$), se tiene que:
\[ \text{Magnitud} = \|\vec{u}\| \cos(\theta) \]

\textbf{3. Construcción del Vector:}
Multiplicando la magnitud por la dirección se obtiene la definición geométrica de la proyección:
\[ \text{proj}_{\vec{v}}\vec{u} = (\|\vec{u}\| \cos\theta) \left(\frac{\vec{v}}{\|\vec{v}\|}\right) \]

Dado que rara vez se conoce el ángulo $\theta$, se puede utilizar la propiedad del producto punto de ángulo entre vectores ($\vec{u} \cdot \vec{v} = \|\vec{u}\|\|\vec{v}\| \cos\theta$) para despejar el término del coseno ($\cos\theta = \frac{\vec{u} \cdot \vec{v}}{\|\vec{u}\|\|\vec{v}\|}$) y sustituirlo en la ecuación:

\[ \text{proj}_{\vec{v}}\vec{u} = \left( \|\vec{u}\| \frac{\vec{u} \cdot \vec{v}}{\|\vec{u}\|\|\vec{v}\|} \right) \frac{\vec{v}}{\|\vec{v}\|} \]

Al simplificar, cancelando $\|\vec{u}\|$ y multiplicando los denominadores, se llega a la fórmula algebraica final que es la que se utiliza para calcular proyecciones en $\mathbb{R}^n$:

\begin{mydefinition}{Proyección Vectorial}
La proyección vectorial de $\vec{u}$ sobre $\vec{v}$ se calcula únicamente usando el producto punto y la magnitud del vector base:
\[ \text{proj}_{\vec{v}}\vec{u} = \left( \frac{\vec{u} \cdot \vec{v}}{\|\vec{v}\|^2} \right) \vec{v} \]
\end{mydefinition}

\subsubsection{Magnitud del vector proyección}

En muchas ocasiones, no es de interés el vector completo resultante de la proyección, sino únicamente su longitud. A esta longitud se le conoce como \textbf{proyección escalar} o \textbf{componente} de $\vec{u}$ a lo largo de $\vec{v}$.

Como se realizó en el desarrollo geométrico, esta magnitud está dada por $\|\vec{u}\| \cos\theta$. Al igual que antes, si sustituimos el coseno utilizando el producto punto, se obtiene su fórmula directa.

\begin{mydefinition}{Proyección Escalar}
La proyección escalar o componente del vector $\vec{u}$ en la dirección del vector $\vec{v}$ se denota como $\text{comp}_{\vec{v}}\vec{u}$ y es un número real definido como:
\[ \text{comp}_{\vec{v}}\vec{u} = \|\vec{u}\| \cos(\theta) = \frac{\vec{u} \cdot \vec{v}}{\|\vec{v}\|} \]
\end{mydefinition}

\textit{Nota:} Aunque se habla de longitud, la componente escalar puede ser un número negativo. Si el ángulo $\theta$ entre los vectores es mayor a $\frac{\pi}{2}$, el coseno será negativo. Geométricamente, esto significa que la proyección sombra está apuntando en la dirección opuesta al vector $\vec{v}$.

\subsection{Ejercicios}

\begin{enumerate}
    \item Calcular el producto punto de cada par de vectores:
\begin{enumerate}
    \item $\vec{u} = (3, -2)$ y $\vec{v} = (1, 4)$
    \item $\vec{a} = (1, -3, 2)$ y $\vec{b} = (4, 1, -2)$
\end{enumerate}
    \item Sea $\vec{v} = (2, -3, 6)$. Verificar la propiedad $\vec{v} \cdot \vec{v} = \|\vec{v}\|^2$ calculando ambos lados de la igualdad por separado.
    \item Sean $\vec{u} = (1, 2, -1)$, $\vec{v} = (3, 0, 2)$ y $c = 4$. Verificar las siguientes propiedades del producto punto usando los valores dados:
\begin{enumerate}
    \item Conmutatividad: $\vec{u} \cdot \vec{v} = \vec{v} \cdot \vec{u}$
    \item Asociatividad con escalar: $c(\vec{u} \cdot \vec{v}) = (c\vec{u}) \cdot \vec{v}$
\end{enumerate}
    \item Determinar cuáles de los siguientes pares de vectores son ortogonales entre sí:
\begin{enumerate}
    \item $\vec{u} = (2, -3)$ y $\vec{v} = (3, 2)$
    \item $\vec{a} = (1, -1, 2)$ y $\vec{b} = (2, 0, -1)$
    \item $\vec{p} = (4, 2, -1)$ y $\vec{q} = (1, -3, -2)$
\end{enumerate}
    \item Encontrar el ángulo $\theta$ (en grados) que forman los siguientes pares de vectores:
\begin{enumerate}
    \item $\vec{u} = (1, 0, 0)$ y $\vec{v} = (1, 1, 0)$
    \item $\vec{a} = (2, -1, 3)$ y $\vec{b} = (1, 2, -1)$
\end{enumerate}
    \item Encontrar el valor o los valores del parámetro $k$ para que los siguientes vectores sean ortogonales:
\begin{enumerate}
    \item $\vec{u} = (k, 3)$ y $\vec{v} = (2, -k)$
    \item $\vec{a} = (1, k, -2)$ y $\vec{b} = (3, -1, k)$
\end{enumerate}
    \item Dado el vector $\vec{v} = (1, 2, -2)$:
\begin{enumerate}
    \item Calcular los ángulos directores $\alpha$, $\beta$ y $\gamma$ (en grados).
    \item Verificar la identidad $\cos^2(\alpha) + \cos^2(\beta) + \cos^2(\gamma) = 1$.
    \item Escribir el vector unitario en la dirección de $\vec{v}$ en términos de sus cosenos directores.
\end{enumerate}
    \item Sean $\vec{u} = (2, 3, -1)$ y $\vec{v} = (4, -1, 2)$. Calcular:
\begin{enumerate}
    \item La proyección escalar $\text{comp}_{\vec{v}}\vec{u}$.
    \item La proyección vectorial $\text{proj}_{\vec{v}}\vec{u}$.
\end{enumerate}
    \item Dado $\vec{u} = (1, -2, 3)$ y $\vec{v} = (2, 1, 1)$, descomponer $\vec{u}$ en dos componentes: una paralela a $\vec{v}$ y otra perpendicular a $\vec{v}$, de modo que:
\[ \vec{u} = \vec{u}_{\parallel} + \vec{u}_{\perp} \]
Verificar que $\vec{u}_\perp \cdot \vec{v} = 0$.
    \item Una fuerza $\vec{F} = (3, -2, 5)$ N actúa sobre un objeto que se desplaza del punto $A = (1, 0, -1)$ al punto $B = (4, 2, 3)$ (distancias en metros).
\begin{enumerate}
    \item Encontrar el vector desplazamiento $\vec{d} = \overrightarrow{AB}$.
    \item Calcular el trabajo $W = \vec{F} \cdot \vec{d}$ realizado por la fuerza (en joules).
    \item Determinar el ángulo que forma la fuerza con el desplazamiento.
\end{enumerate}
    \item Demostrar algebraicamente, usando la definición componente a componente, la propiedad distributiva del producto punto para vectores en $\mathbb{R}^n$:
\[ \vec{u} \cdot (\vec{v} + \vec{w}) = \vec{u} \cdot \vec{v} + \vec{u} \cdot \vec{w} \]
Usar la demostración para verificar con $\vec{u} = (1,2,-1)$, $\vec{v} = (3,0,1)$, $\vec{w} = (-1,2,4)$.
    \item Un sensor de temperatura está ubicado en el punto $P = (2, 1, 3)$. El gradiente de temperatura en ese punto apunta en la dirección $\vec{g} = (4, -2, 1)$. Un robot puede moverse en cualquiera de las siguientes tres direcciones:
\[ \vec{d}_1 = (1, 2, -1), \quad \vec{d}_2 = (-1, 3, 2), \quad \vec{d}_3 = (2, 0, -3) \]
\begin{enumerate}
    \item Calcular la proyección escalar del gradiente sobre cada dirección de movimiento posible.
    \item ¿En qué dirección debe moverse el robot para maximizar el incremento de temperatura? ¿Y para minimizarlo?
\end{enumerate}
\end{enumerate}

% \subsection{El Producto Vectorial de Gibbs (Producto Cruz)}

\subsubsection{Naturaleza y Origen Histórico}

A diferencia del producto escalar, el cual produce un número real, se define una segunda forma de multiplicación vectorial cuyo resultado es un nuevo vector. A esta operación se le denomina \textbf{producto vectorial} o \textbf{producto cruz}. La característica geométrica fundamental del vector resultante es que este es simultáneamente ortogonal a los dos vectores originales que fueron multiplicados.

Históricamente, el producto vectorial moderno fue formulado de manera independiente a finales del siglo XIX por el físico estadounidense Josiah Willard Gibbs y el ingeniero británico Oliver Heaviside. Este operador matemático surge como una extracción y simplificación del álgebra de los cuaterniones (desarrollada previamente por William Rowan Hamilton). Fue diseñado específicamente para facilitar el análisis espacial en áreas como el electromagnetismo y la mecánica clásica.

\subsubsection{Dimensionalidad del Producto Vectorial}

Es importante mencionar una particularidad matemática sobre esta operación. El producto cruz tradicional, entendido como una operación binaria que toma dos vectores y devuelve un tercer vector ortogonal a ambos, es un fenómeno algebraico que está estrictamente definido solo en el espacio tridimensional $\mathbb{R}^3$ (y, de forma excepcional, en $\mathbb{R}^7$). 

No obstante, el concepto de producto vectorial sí existe y es generalizado para vectores en un espacio de $N$ dimensiones ($\mathbb{R}^n$). Esta generalización es formalizada mediante el álgebra exterior (utilizando el producto cuña o \textit{wedge product}, $\wedge$), donde el resultado no es un vector tradicional, sino una entidad matemática de dimensión superior llamada bivector o multivector. 

Sin embargo, para los objetivos y el alcance del cálculo multivariable a este nivel, el estudio del producto vectorial se limitará exclusivamente a la multiplicación de dos vectores en el espacio tridimensional $\mathbb{R}^3$.

\subsubsection{Definición Algebraica en $\mathbb{R}^3$}

El cálculo algebraico del producto cruz se realiza a partir de las componentes rectangulares de los vectores. Se utiliza como herramienta mnemotécnica el cálculo del determinante de una matriz de $3 \times 3$, en cuya primera fila se colocan los vectores unitarios canónicos.

\begin{mydefinition}{Producto Cruz en $\mathbb{R}^3$}
Sean $\vec{u} = (u_1, u_2, u_3)$ y $\vec{v} = (v_1, v_2, v_3)$ dos vectores en $\mathbb{R}^3$. El producto vectorial $\vec{u} \times \vec{v}$ (léase ``$\vec{u}$ cruz $\vec{v}$'') es el vector definido por el siguiente determinante formal:

\[ \vec{u} \times \vec{v} = \begin{vmatrix} 
\hat{\imath} & \hat{\jmath} & \hat{k} \\ 
u_1 & u_2 & u_3 \\ 
v_1 & v_2 & v_3 
\end{vmatrix} \]

Al expandir el determinante por los cofactores de la primera fila, se obtiene la fórmula algebraica:

\[ \vec{u} \times \vec{v} = \begin{vmatrix} u_2 & u_3 \\ v_2 & v_3 \end{vmatrix}\hat{\imath} - \begin{vmatrix} u_1 & u_3 \\ v_1 & v_3 \end{vmatrix}\hat{\jmath} + \begin{vmatrix} u_1 & u_2 \\ v_1 & v_2 \end{vmatrix}\hat{k} \]

\[ \vec{u} \times \vec{v} = (u_2v_3 - u_3v_2)\hat{\imath} - (u_1v_3 - u_3v_1)\hat{\jmath} + (u_1v_2 - u_2v_1)\hat{k} \]
\end{mydefinition}

\subsubsection{Interpretación Geométrica y Regla de la Mano Derecha}

Como se estableció anteriormente, el vector resultante $\vec{u} \times \vec{v}$ es perpendicular tanto al vector $\vec{u}$ como al vector $\vec{v}$. Por consiguiente, es perpendicular al plano que contiene a ambos vectores.

El sentido de este nuevo vector es determinado por la \textbf{Regla de la Mano Derecha}. Si se extienden los dedos de la mano derecha en la dirección del primer vector ($\vec{u}$) y se flexionan hacia el segundo vector ($\vec{v}$), el pulgar extendido indicará la dirección y sentido del vector resultante $\vec{u} \times \vec{v}$. Esta dependencia del orden implica que el producto cruz no es conmutativo, se cumple que $\vec{u} \times \vec{v} = -(\vec{v} \times \vec{u})$.

\begin{center}
\shorthandoff{>}
\begin{tikzpicture}[x={(-0.5cm,-0.4cm)}, y={(1cm,0cm)}, z={(0cm,1cm)}, scale=1.5]
    % Plano que contiene a u y v
    \filldraw[fill=blue!10!white, draw=blue!50!white, thick] 
        (-2,-1,0) -- (3,-1,0) -- (4,3,0) -- (-1,3,0) -- cycle;
    \node[blue!80!black] at (-1.5, 1, 0) {Plano de $\vec{u}$ y $\vec{v}$};
    
    % Vectores base u y v
    \draw[->, very thick, blue] (0,0,0) -- (2,0,0) node[below left] {$\vec{u}$};
    \draw[->, very thick, red] (0,0,0) -- (0,2,0) node[below right] {$\vec{v}$};
    
    % Ángulo theta entre u y v
    \draw[thick] (0.6,0,0) arc (0:90:0.6) node[midway, below=3pt] {$\theta$};
    
    % Vector Producto Cruz Resultante
    \draw[->, very thick, purple] (0,0,0) -- (0,0,3) node[above] {$\vec{u} \times \vec{v}$};
    
    % Indicadores de perpendicularidad
    % Con u
    \draw[thick] (0.3,0,0) -- (0.3,0,0.3) -- (0,0,0.3);
    % Con v
    \draw[thick] (0,0.3,0) -- (0,0.3,0.3) -- (0,0,0.3);
    
    % Punto de origen
    \filldraw[black] (0,0,0) circle (1.5pt) node[left=3pt] {Origen};
\end{tikzpicture}
\end{center}

Para verificar que dos vectores son perpendiculares, basta con demostrar que su producto punto es igual a cero. Por lo tanto, si $\vec{w} = \vec{u} \times \vec{v}$, se debe cumplir de manera simultánea que $\vec{w} \cdot \vec{u} = 0$ y $\vec{w} \cdot \vec{v} = 0$.

\begin{proof}[Demostración de Ortogonalidad]
Sean $\vec{u} = (u_1, u_2, u_3)$ y $\vec{v} = (v_1, v_2, v_3)$. A partir de la definición algebraica del producto cruz obtenida por los cofactores, se tiene el vector resultante:
\[
\vec{u} \times \vec{v} = (u_2v_3 - u_3v_2, \, u_3v_1 - u_1v_3, \, u_1v_2 - u_2v_1)
\]

Se evaluará primero el producto punto entre el vector resultante y el vector original $\vec{u}$:
\[
(\vec{u} \times \vec{v}) \cdot \vec{u} = (u_2v_3 - u_3v_2)(u_1) + (u_3v_1 - u_1v_3)(u_2) + (u_1v_2 - u_2v_1)(u_3)
\]

Distribuyendo cada componente de $\vec{u}$ en los binomios correspondientes:
\[
(\vec{u} \times \vec{v}) \cdot \vec{u} = (u_1u_2v_3 - u_1u_3v_2) + (u_2u_3v_1 - u_1u_2v_3) + (u_1u_3v_2 - u_2u_3v_1)
\]

Al eliminar los paréntesis y reorganizar los términos, se observa que todos los componentes se cancelan mutuamente:
\[
(\vec{u} \times \vec{v}) \cdot \vec{u} = u_1u_2v_3 - u_1u_2v_3 - u_1u_3v_2 + u_1u_3v_2 + u_2u_3v_1 - u_2u_3v_1
\]
\[
(\vec{u} \times \vec{v}) \cdot \vec{u} = 0
\]

Este resultado demuestra algebraicamente que el vector $\vec{u} \times \vec{v}$ es perpendicular al vector $\vec{u}$. Realizando el mismo desarrollo para el producto $(\vec{u} \times \vec{v}) \cdot \vec{v}$ resultará también $0$, demostrando que el vector resultante es ortogonal a ambos vectores base de manera simultánea.
\end{proof}

\subsubsection{Propiedades Algebraicas del Producto Vectorial}

Al igual que las otras operaciones vectoriales, el producto cruz posee propiedades matemáticas definidas. Sean $\vec{u}$, $\vec{v}$ y $\vec{w}$ vectores en $\mathbb{R}^3$, y sea $c$ un escalar. Se cumplen las siguientes propiedades:

\begin{mytheorem}{Propiedades del Producto Vectorial}
\begin{enumerate}
    \item \textbf{Anticonmutatividad:} $\vec{u} \times \vec{v} = -(\vec{v} \times \vec{u})$
    \item \textbf{Distributividad sobre la suma:} $\vec{u} \times (\vec{v} + \vec{w}) = (\vec{u} \times \vec{v}) + (\vec{u} \times \vec{w})$
    \item \textbf{Asociatividad escalar:} $c(\vec{u} \times \vec{v}) = (c\vec{u}) \times \vec{v} = \vec{u} \times (c\vec{v})$
    \item \textbf{Producto con el vector cero:} $\vec{u} \times \vec{0} = \vec{0} \times \vec{u} = \vec{0}$
    \item \textbf{Producto consigo mismo:} $\vec{u} \times \vec{u} = \vec{0}$
\end{enumerate}
\end{mytheorem}

\subsubsection{Magnitud del Producto Vectorial}

De manera análoga al producto punto, existe una relación entre la magnitud del producto vectorial y el ángulo $\theta$ que se forma entre los dos vectores originales (donde $0 \leq \theta \leq \pi$). Mientras que el producto punto se relaciona con el coseno del ángulo, la magnitud del producto cruz está directamente vinculada con la función seno.

\begin{mytheorem}{Magnitud del Producto Cruz}
Si $\theta$ es el ángulo entre dos vectores no nulos $\vec{u}$ y $\vec{v}$, entonces la magnitud de su producto cruz es:
\[ \|\vec{u} \times \vec{v}\| = \|\vec{u}\| \|\vec{v}\| \sin(\theta) \]
\end{mytheorem}

\begin{proof}[Demostración]
Sean $\vec{u} = (u_1, u_2, u_3)$ y $\vec{v} = (v_1, v_2, v_3)$. A partir de la definición algebraica del producto cruz, se calcula el cuadrado de su magnitud elevando al cuadrado cada una de sus componentes:
\[ \|\vec{u} \times \vec{v}\|^2 = (u_2v_3 - u_3v_2)^2 + (u_3v_1 - u_1v_3)^2 + (u_1v_2 - u_2v_1)^2 \]

Se desarrollan los tres binomios al cuadrado:
\[ \|\vec{u} \times \vec{v}\|^2 = (u_2v_3 - u_3v_2)^2 + (u_3v_1 - u_1v_3)^2 + (u_1v_2 - u_2v_1)^2 \]

Para agrupar estos términos, se suma y se resta la misma cantidad, específicamente la suma de los cuadrados de los productos de componentes homólogas $(u_1^2v_1^2 + u_2^2v_2^2 + u_3^2v_3^2)$, lo cual no altera el valor de la ecuación:
\[ \quad + u_1^2v_1^2 + u_2^2v_2^2 + u_3^2v_3^2 - (u_1^2v_1^2 + u_2^2v_2^2 + u_3^2v_3^2) \]

Al reorganizar la expresión, se agrupan todos los términos positivos por un lado y los negativos por otro:
\[ \|\vec{u} \times \vec{v}\|^2 = (u_1^2v_1^2 + u_1^2v_2^2 + u_1^2v_3^2 + u_2^2v_1^2 + u_2^2v_2^2 + u_2^2v_3^2 + u_3^2v_1^2 + u_3^2v_2^2 + u_3^2v_3^2) \]
\[ \quad \quad \quad - (u_1^2v_1^2 + u_2^2v_2^2 + u_3^2v_3^2 + 2u_1v_1u_2v_2 + 2u_1v_1u_3v_3 + 2u_2v_2u_3v_3) \]

El primer grupo de términos positivos corresponde exactamente al desarrollo del producto de las magnitudes al cuadrado de los vectores originales. El segundo grupo corresponde al trinomio cuadrado perfecto del producto punto. Por lo tanto, se factoriza la expresión como:
\[ \|\vec{u} \times \vec{v}\|^2 = (u_1^2 + u_2^2 + u_3^2)(v_1^2 + v_2^2 + v_3^2) - (u_1v_1 + u_2v_2 + u_3v_3)^2 \]

Sustituyendo por las definiciones formales de magnitud y producto escalar:
\[ \|\vec{u} \times \vec{v}\|^2 = \|\vec{u}\|^2 \|\vec{v}\|^2 - (\vec{u} \cdot \vec{v})^2 \]

A partir del teorema del producto punto, se sabe que $\vec{u} \cdot \vec{v} = \|\vec{u}\|\|\vec{v}\|\cos(\theta)$. Sustituyendo esta expresión geométrica en la ecuación:
\[ \|\vec{u} \times \vec{v}\|^2 = \|\vec{u}\|^2 \|\vec{v}\|^2 - (\|\vec{u}\|\|\vec{v}\|\cos(\theta))^2 \]

Distribuyendo el exponente cuadrado en el segundo término y factorizando el elemento común $\|\vec{u}\|^2 \|\vec{v}\|^2$:
\[ \|\vec{u} \times \vec{v}\|^2 = \|\vec{u}\|^2 \|\vec{v}\|^2 - \|\vec{u}\|^2\|\vec{v}\|^2\cos^2(\theta) \]
\[ \|\vec{u} \times \vec{v}\|^2 = \|\vec{u}\|^2 \|\vec{v}\|^2 (1 - \cos^2(\theta)) \]

Aplicando la identidad trigonométrica pitagórica fundamental ($1 - \cos^2\theta = \sin^2\theta$):
\[ \|\vec{u} \times \vec{v}\|^2 = \|\vec{u}\|^2 \|\vec{v}\|^2 \sin^2(\theta) \]

Finalmente, al extraer la raíz cuadrada en ambos lados de la ecuación, y considerando que $\sin(\theta) \geq 0$ para ángulos comprendidos entre $0$ y $\pi$ ($0^\circ \leq \theta \leq 180^\circ$), se concluye la demostración:
\[ \|\vec{u} \times \vec{v}\| = \|\vec{u}\| \|\vec{v}\| \sin(\theta) \]
\end{proof}

\subsubsection{Condición de Paralelismo}

Este teorema sobre la magnitud ofrece una herramienta algebraica directa para determinar si dos vectores son paralelos. Si dos vectores no nulos $\vec{u}$ y $\vec{v}$ son paralelos (o colineales), el ángulo entre ellos es exactamente $\theta = 0$ (si apuntan en la misma dirección) o $\theta = \pi$ (si apuntan en direcciones opuestas).

En ambos casos, el seno del ángulo es cero. Sustituyendo este valor en la fórmula de la magnitud:
\[ \|\vec{u} \times \vec{v}\| = \|\vec{u}\| \|\vec{v}\| (0) = 0 \]

Dado que el único vector cuya magnitud es cero es el vector nulo ($\vec{0}$), se deduce el siguiente criterio:

\begin{mytheorem}{Vectores Paralelos}
Dos vectores no nulos $\vec{u}$ y $\vec{v}$ son paralelos si y solo si su producto vectorial es el vector cero:
\[ \vec{u} \times \vec{v} = \vec{0} \]
\end{mytheorem}

\subsection{Productos Triples}

En el análisis vectorial y sus aplicaciones físicas, es frecuente encontrarse con operaciones que involucran tres vectores simultáneamente. Existen dos formas principales de combinar tres vectores utilizando los productos escalar y vectorial: el triple producto escalar y el triple producto vectorial.

\subsubsection{El Triple Producto Escalar}

El triple producto escalar es una operación que toma tres vectores en $\mathbb{R}^3$ y produce un único número real. Se define combinando un producto vectorial seguido de un producto escalar.

\begin{mydefinition}{Triple Producto Escalar}
Sean $\vec{u}$, $\vec{v}$ y $\vec{w}$ tres vectores en el espacio tridimensional. El triple producto escalar se define como:
\[ \vec{u} \cdot (\vec{v} \times \vec{w}) \]
\end{mydefinition}

Algebraicamente, esta operación puede ser calculada de manera directa mediante el cálculo del determinante de una matriz de $3 \times 3$ formada por las componentes de los tres vectores. Si $\vec{u} = (u_1, u_2, u_3)$, $\vec{v} = (v_1, v_2, v_3)$ y $\vec{w} = (w_1, w_2, w_3)$, entonces:

\[ \vec{u} \cdot (\vec{v} \times \vec{w}) = \begin{vmatrix} 
u_1 & u_2 & u_3 \\ 
v_1 & v_2 & v_3 \\ 
w_1 & w_2 & w_3 
\end{vmatrix} \]

Una propiedad algebraica del triple producto escalar es que los operadores de punto y cruz pueden intercambiarse sin alterar el resultado, siempre y cuando se mantenga el orden de los vectores. Es decir, se cumple que $\vec{u} \cdot (\vec{v} \times \vec{w}) = (\vec{u} \times \vec{v}) \cdot \vec{w}$.

\subsubsection{El Triple Producto Vectorial}

A diferencia del caso anterior, el triple producto vectorial es una operación que devuelve como resultado un nuevo vector en el espacio. Se obtiene al realizar el producto cruz de un vector con el resultado de un producto cruz previo.

\begin{mydefinition}{Triple Producto Vectorial}
Sean $\vec{u}$, $\vec{v}$ y $\vec{w}$ tres vectores en $\mathbb{R}^3$. El triple producto vectorial se denota formalmente como:
\[ \vec{u} \times (\vec{v} \times \vec{w}) \]
\end{mydefinition}

Es crucial recalcar que, debido a que el producto vectorial no goza de la propiedad asociativa, la colocación de los paréntesis es estrictamente obligatoria, pues se cumple que $\vec{u} \times (\vec{v} \times \vec{w}) \neq (\vec{u} \times \vec{v}) \times \vec{w}$.

Desde un análisis geométrico, el vector resultante de $\vec{u} \times (\vec{v} \times \vec{w})$ yace exactamente en el mismo plano que está siendo formado por los vectores $\vec{v}$ y $\vec{w}$. Esto se debe a que el resultado de la operación debe ser ortogonal al vector $(\vec{v} \times \vec{w})$, y dado que este último es ortogonal al plano de $\vec{v}$ y $\vec{w}$, el vector final se ve forzado a existir dentro de dicho plano inicial.

\paragraph{Regla de BAC-CAB}
Para calcular este producto sin necesidad de evaluar dos determinantes matrices de forma consecutiva, se emplea una identidad matemática conocida como la Fórmula de Expulsión o regla BAC-CAB. Esta identidad permite expresar el vector resultante de forma mucho más directa, como una combinación lineal de los vectores $\vec{v}$ y $\vec{w}$, utilizando únicamente operaciones de producto punto y escalado:

\begin{mytheorem}{Identidad del Triple Producto Vectorial}
\[ \vec{u} \times (\vec{v} \times \vec{w}) = \vec{v}(\vec{u} \cdot \vec{w}) - \vec{w}(\vec{u} \cdot \vec{v}) \]
\end{mytheorem}

\paragraph{Demostración mediante Notación Tensorial}

La demostración analítica tradicional de la regla BAC-CAB mediante el cálculo de determinantes resulta en un proceso algebraico sumamente extenso y propenso a errores. Para simplificar esta demostración, se emplea la notación tensorial.

Antes de proceder con la demostración, es necesario definir dos operadores fundamentales en esta notación:

\begin{enumerate}
    \item \textbf{La Delta de Kronecker ($\delta_{ij}$):} Es un operador que vale $1$ si los índices son iguales ($i = j$) y $0$ si son diferentes ($i \neq j$). Al multiplicar un vector por la delta de Kronecker, ocurre un fenómeno llamado contracción de índices, que simplemente renombra el índice: $\delta_{ij}v_j = v_i$. Además, el producto punto se define como: $\vec{u} \cdot \vec{v} = u_i v_i$.
    \item \textbf{El Símbolo de Levi-Civita ($\epsilon_{ijk}$):} Es un tensor alternante que vale $1$ si los índices $(i,j,k)$ son una permutación cíclica de $(1,2,3)$, vale $-1$ si es una permutación anticíclica, y vale $0$ si algún índice se repite. Su utilidad principal es definir el producto cruz en su $i$-ésima componente: 
    \[ (\vec{u} \times \vec{v})_i = \epsilon_{ijk} u_j v_k \]
\end{enumerate}

El elemento clave para esta demostración es una identidad matemática que relaciona el producto de dos símbolos de Levi-Civita con deltas de Kronecker, conocida como la identidad epsilon-delta:
\[ \epsilon_{ijk} \epsilon_{klm} = \epsilon_{kij} \epsilon_{klm} = \delta_{il}\delta_{jm} - \delta_{im}\delta_{jl} \]

Con estas herramientas, se procede a la demostración del teorema.

\begin{proof}[Demostración de la Regla BAC-CAB]
Sea el vector resultante $\vec{R} = \vec{u} \times (\vec{v} \times \vec{w})$. Se buscará aislar la $i$-ésima componente de este vector, es decir, $R_i$.

\textbf{Paso 1: Aplicación del primer producto cruz.}
Se aplica la definición tensorial del producto cruz al operador externo, asignando los índices $i, j, k$:
\[ R_i = [\vec{u} \times (\vec{v} \times \vec{w})]_i = \epsilon_{ijk} u_j (\vec{v} \times \vec{w})_k \]
\textbf{Paso 2: Aplicación del segundo producto cruz.}
Ahora se aplica la misma definición para el producto cruz interno, cuya componente es $k$. Para evitar confusión, se utilizan nuevos índices $l, m$ para los vectores internos:
\[ (\vec{v} \times \vec{w})_k = \epsilon_{klm} v_l w_m \]

\textbf{Paso 3: Sustitución y reordenamiento.}
Se sustituye el resultado del paso 2 en la ecuación del paso 1. Como los escalares conmutan, se agrupan los símbolos de Levi-Civita al inicio de la expresión:
\[ R_i = \epsilon_{ijk} u_j (\epsilon_{klm} v_l w_m) = \epsilon_{ijk} \epsilon_{klm} u_j v_l w_m \]

Debido a que el símbolo de Levi-Civita permite permutaciones cíclicas sin alterar su signo ($\epsilon_{ijk} = \epsilon_{kij}$), se reescribe el primer símbolo para preparar la aplicación de la identidad:
\[ R_i = \epsilon_{kij} \epsilon_{klm} u_j v_l w_m \]

\textbf{Paso 4: Aplicación de la identidad epsilon-delta.}
Se sustituye el producto de los tensores de Levi-Civita por la identidad de deltas de Kronecker ($\epsilon_{kij} \epsilon_{klm} = \delta_{il}\delta_{jm} - \delta_{im}\delta_{jl}$):
\[ R_i = (\delta_{il}\delta_{jm} - \delta_{im}\delta_{jl}) u_j v_l w_m \]

\textbf{Paso 5: Distribución algebraica.}
Se distribuyen los componentes vectoriales dentro del paréntesis:
\[ R_i = (\delta_{il}\delta_{jm} u_j v_l w_m) - (\delta_{im}\delta_{jl} u_j v_l w_m) \]

\textbf{Paso 6: Contracción de índices.}
Se aplica la propiedad de la delta de Kronecker para colapsar los índices. 
En el primer término:
\begin{itemize}
    \item $\delta_{il} v_l$ transforma el índice $l$ en $i$, resultando en $v_i$.
    \item $\delta_{jm} u_j w_m$ filtra las componentes donde $j = m$, resultando en $u_j w_j$.
\end{itemize}
En el segundo término:
\begin{itemize}
    \item $\delta_{im} w_m$ transforma el índice $m$ en $i$, resultando en $w_i$.
    \item $\delta_{jl} u_j v_l$ filtra las componentes donde $j = l$, resultando en $u_j v_j$.
\end{itemize}

Sustituyendo estas contracciones en la ecuación:
\[ R_i = v_i (u_j w_j) - w_i (u_j v_j) \]

\textbf{Paso 7: Conversión a notación vectorial estándar.}
Se reconoce que los términos con índices repetidos ($u_j w_j$ y $u_j v_j$) corresponden a la definición del producto escalar (producto punto). Por lo tanto:
\[ R_i = v_i (\vec{u} \cdot \vec{w}) - w_i (\vec{u} \cdot \vec{v}) \]

Dado que esta igualdad se cumple para cualquier componente $i$ arbitraria, se generaliza para todo el vector, concluyendo así la demostración matemática:
\[ \vec{R} = \vec{v} (\vec{u} \cdot \vec{w}) - \vec{w} (\vec{u} \cdot \vec{v}) \]
\end{proof}

\subsection{Aplicaciones geométricas}

El producto cruz tiene dos aplicaciones geométricas importantes: el cálculo de áreas y el cálculo de volúmenes. En ambos casos, la magnitud del producto cruz se relaciona directamente con la medida geométrica que se desea calcular.

\subsubsection{Área de un paralelogramo}

Una de las aplicaciones geométricas del producto vectorial radica en su capacidad para calcular áreas en el espacio tridimensional. Específicamente, se establece que la magnitud del producto cruz de dos vectores es igual al área del paralelogramo que estos forman cuando se sitúan con un origen común.

Para comprender el origen de esta propiedad, es necesario recurrir a la geometría euclidiana clásica. El área de cualquier paralelogramo se calcula mediante el producto de la longitud de su base por su altura perpendicular:
\[ \text{Área} = (\text{base}) \times (\text{altura}) \]

Supóngase que se tienen dos vectores no nulos $\vec{u}$ y $\vec{v}$ separados por un ángulo $\theta$. Si estos vectores se utilizan para definir los lados adyacentes de un paralelogramo, se puede realizar la siguiente deducción geométrica:

\begin{enumerate}
    \item \textbf{La base:} Se define arbitrariamente al vector $\vec{u}$ como la base de la figura. Por lo tanto, la longitud de la base es exactamente la magnitud de este vector: $\text{base} = \|\vec{u}\|$.
    \item \textbf{La altura ($h$):} La altura del paralelogramo es la distancia perpendicular trazada desde la punta del vector $\vec{v}$ hasta la línea de acción del vector $\vec{u}$. Esto forma un triángulo rectángulo donde la hipotenusa es la magnitud de $\vec{v}$ ($\|\vec{v}\|$). Utilizando la definición de la función seno, se despeja la altura:
    \[ h = \|\vec{v}\| \sin(\theta) \]
\end{enumerate}

Al sustituir estos dos componentes en la fórmula geométrica del área, se obtiene:
\[ \text{Área} = \|\vec{u}\| (\|\vec{v}\| \sin\theta) \]

Como se demostró anteriormente, la expresión $\|\vec{u}\| \|\vec{v}\| \sin\theta$ corresponde a la magnitud del producto cruz entre dichos vectores. Por consiguiente, se tiene el siguiente teorema:

\begin{mytheorem}{Área del Paralelogramo}
El área $A$ del paralelogramo definido por los vectores adyacentes $\vec{u}$ y $\vec{v}$ está dada por la magnitud de su producto vectorial:
\[ A = \|\vec{u} \times \vec{v}\| \]
\end{mytheorem}

\begin{center}
\shorthandoff{>}
\begin{tikzpicture}[scale=1.5]
    % Base vector u
    \draw[->, very thick, blue] (0,0) -- (4,0) node[midway, below] {$\vec{u}$ (Base)};

    % Vector v
    \draw[->, very thick, red] (0,0) -- (2.5,2) node[midway, above left] {$\vec{v}$};

    % Completando el paralelogramo (líneas punteadas)
    \draw[dashed, thick, gray] (4,0) -- (6.5,2);
    \draw[dashed, thick, gray] (2.5,2) -- (6.5,2);

    % Altura h
    \draw[dashed, thick, purple] (2.5,2) -- (2.5,0) node[midway, right] {$h = \|\vec{v}\|\sin\theta$};

    % Símbolo de ángulo recto para la altura
    \draw (2.3,0) -- (2.3,0.2) -- (2.5,0.2);

    % Ángulo theta
    \draw[thick] (0.8,0) arc (0:38.6:0.8) node[midway, right=2pt, yshift=2pt] {$\theta$};

    % Etiqueta del Área
    \node[font=\footnotesize, fill=blue!5!white, draw=blue!30!white, rounded corners, inner sep=4pt] 
        at (3.5,1) {Área $= \|\vec{u} \times \vec{v}\|$};
        
    % Punto de origen
    \filldraw[black] (0,0) circle (1.5pt) node[below left] {Origen};
\end{tikzpicture}
\end{center}

\paragraph{Corolario: Área de un Triángulo}
Una consecuencia inmediata de este teorema es su aplicación para encontrar el área de un triángulo en el espacio tridimensional. Dado que cualquier paralelogramo puede ser dividido en dos triángulos congruentes trazando una diagonal desde el origen hasta el vértice opuesto, el área del triángulo formado por los vectores $\vec{u}$ y $\vec{v}$ es exactamente la mitad del área del paralelogramo:
\[ \text{Área del triángulo} = \frac{1}{2} \|\vec{u} \times \vec{v}\| \]

\subsubsection{Volúmen de un paralelepípedo}

Mientras que la magnitud del producto cruz simple ($\|\vec{v} \times \vec{w}\|$) proporciona el área de un paralelogramo bidimensional, el triple producto escalar añade una tercera dimensión al cálculo. El valor absoluto de esta operación representa geométricamente el volumen del paralelepípedo tridimensional que se forma al considerar a los tres vectores como aristas adyacentes que comparten un mismo origen.

La justificación geométrica de esto radica en la fórmula del volumen de un prisma. 
\begin{itemize}
    \item El área de la base, conformada por los vectores $\vec{v}$ y $\vec{w}$, está dada por $\|\vec{v} \times \vec{w}\|$.
    \item La altura $h$ del sólido corresponde a la componente del vector $\vec{u}$ sobre el eje ortogonal a la base. Esta altura es igual a $\|\vec{u}\||\cos\theta|$, donde $\theta$ es el ángulo entre $\vec{u}$ y el vector resultante de $\vec{v} \times \vec{w}$.
\end{itemize}

Por lo tanto, al multiplicar la base por la altura, se recupera exactamente la definición del producto punto:
\[ \text{Volumen} = \|\vec{v} \times \vec{w}\| (\|\vec{u}\| |\cos\theta|) = |\vec{u} \cdot (\vec{v} \times \vec{w})| \]

\begin{center}
\shorthandoff{>}
\begin{tikzpicture}[x={(-0.5cm,-0.4cm)}, y={(1cm,0cm)}, z={(0cm,1cm)}, scale=1.3]
    % Vectores base (Base del paralelepipedo)
    \coordinate (O) at (0,0,0);
    \coordinate (V) at (3,0,0);
    \coordinate (W) at (0,3,0);
    \coordinate (U) at (1,1.5,2.5);
    
    % Vértices superiores del paralelepipedo
    \coordinate (VU) at (4,1.5,2.5);
    \coordinate (WU) at (1,4.5,2.5);
    \coordinate (VW) at (3,3,0);
    \coordinate (VWU) at (4,4.5,2.5);
    
    % Dibujo de las aristas ocultas (dashed)
    \draw[dashed, gray] (O) -- (W);
    \draw[dashed, gray] (W) -- (VW);
    \draw[dashed, gray] (W) -- (WU);
    
    % Dibujo de las aristas visibles
    \draw[thick, gray] (V) -- (VW);
    \draw[thick, gray] (U) -- (VU);
    \draw[thick, gray] (U) -- (WU);
    \draw[thick, gray] (VU) -- (VWU);
    \draw[thick, gray] (WU) -- (VWU);
    \draw[thick, gray] (V) -- (VU);
    \draw[thick, gray] (VW) -- (VWU);
    
    % Relleno tenue de la base (paralelogramo)
    \filldraw[fill=blue!10!white, opacity=0.7] (O) -- (V) -- (VW) -- (W) -- cycle;
    \node[blue!80!black, font=\footnotesize] at (1.5, 1.5, 0) {Área base $= \|\vec{v} \times \vec{w}\|$};
    
    % Vectores principales
    \draw[->, very thick, red] (O) -- (V) node[below left] {$\vec{v}$};
    \draw[->, very thick, blue] (O) -- (W) node[below right] {$\vec{w}$};
    \draw[->, very thick, green!60!black] (O) -- (U) node[above left] {$\vec{u}$};
    
    % Vector Producto Cruz (Normal a la base)
    \coordinate (Nx) at (0,0,3.5);
    \draw[->, very thick, purple] (O) -- (Nx) node[right] {$\vec{v} \times \vec{w}$};
    
    % Proyección de u sobre el vector normal (Altura h)
    \draw[dashed, thick, black] (U) -- (0,0,2.5);
    \draw[decorate, decoration={brace, amplitude=5pt}, thick] (0.2,0,0) -- (0.2,0,2.5) node[midway, left=6pt] {$h$};
    
    % Ángulo theta
    \draw[thick] (0,0,0.8) arc (0:58:0.4) node[midway, above=2pt] {$\theta$};
\end{tikzpicture}
\end{center}

\subsection{Ejercicios}

\begin{enumerate}
    \item Calcular el producto cruz $\vec{u} \times \vec{v}$ para cada par de vectores:
\begin{enumerate}
    \item $\vec{u} = (1, 2, 3)$ y $\vec{v} = (4, 5, 6)$
    \item $\vec{u} = (2, -1, 3)$ y $\vec{v} = (0, 4, -2)$
\end{enumerate}
    \item Sean $\vec{u} = (3, -1, 2)$ y $\vec{v} = (1, 4, -1)$. Calcular $\vec{u} \times \vec{v}$ y $\vec{v} \times \vec{u}$. Verificar que $\vec{u} \times \vec{v} = -(\vec{v} \times \vec{u})$.
    \item Sean $\vec{u} = (2, 1, -3)$ y $\vec{v} = (-1, 3, 2)$. Calcular $\vec{w} = \vec{u} \times \vec{v}$ y verificar algebraicamente que $\vec{w}$ es ortogonal tanto a $\vec{u}$ como a $\vec{v}$, es decir, que $\vec{w} \cdot \vec{u} = 0$ y $\vec{w} \cdot \vec{v} = 0$.
    \item Determinar, usando el producto cruz, si los siguientes pares de vectores son paralelos:
\begin{enumerate}
    \item $\vec{u} = (2, -4, 6)$ y $\vec{v} = (-1, 2, -3)$
    \item $\vec{a} = (1, 2, 3)$ y $\vec{b} = (2, 5, 3)$
\end{enumerate}
    \item Sean $\vec{u} = (1, 2, 2)$ y $\vec{v} = (3, 0, 4)$. Usando la relación $\|\vec{u} \times \vec{v}\| = \|\vec{u}\|\|\vec{v}\|\sin\theta$, encontrar el ángulo $\theta$ (en grados) entre los dos vectores. Comparar el resultado obtenido usando la relación análoga con el producto punto.
    \item Encontrar el valor del parámetro $k$ tal que los vectores $\vec{u} = (1, k, 2)$ y $\vec{v} = (3, 1, k)$ sean paralelos.
    \item Calcular el área del paralelogramo y del triángulo formados por los vectores $\vec{u} = (3, -1, 4)$ y $\vec{v} = (2, 2, -1)$ cuando se sitúan con un origen común.
    \item Dados los tres puntos $A = (1, 0, 2)$, $B = (3, 1, 0)$ y $C = (0, 2, 1)$ en $\mathbb{R}^3$:
\begin{enumerate}
    \item Encontrar los vectores $\overrightarrow{AB}$ y $\overrightarrow{AC}$.
    \item Calcular el área del triángulo $ABC$.
\end{enumerate}
    \item Sean $\vec{u} = (1, 0, 2)$, $\vec{v} = (-1, 3, 1)$ y $\vec{w} = (2, 1, -1)$.
\begin{enumerate}
    \item Calcular el triple producto escalar $\vec{u} \cdot (\vec{v} \times \vec{w})$ usando el determinante $3\times 3$.
    \item Verificar que $\vec{u} \cdot (\vec{v} \times \vec{w}) = (\vec{u} \times \vec{v}) \cdot \vec{w}$.
    \item Interpretar geométricamente el resultado: ¿cuál es el volumen del paralelepípedo que forman los tres vectores?
\end{enumerate}
    \item Sean $\vec{u} = (2, 0, 1)$, $\vec{v} = (1, -1, 2)$ y $\vec{w} = (3, 1, 0)$. Verificar la regla BAC-CAB calculando ambos lados de la identidad para el triple producto vectorial:
\[ \vec{u} \times (\vec{v} \times \vec{w}) = \vec{v}(\vec{u} \cdot \vec{w}) - \vec{w}(\vec{u} \cdot \vec{v}) \]
    \item Una llave ejerce una fuerza $\vec{F} = (0, -20, 0)$ N sobre una tuerca ubicada en el punto $Q = (0.3, 0, 0)$ m respecto al punto de pivote $P$ (en metros). El torque mecánico es el producto cruz $\vec{\tau} = \overrightarrow{PQ} \times \vec{F}$.
\begin{enumerate}
    \item Calcular el vector de torque $\vec{\tau}$.
    \item Determinar la magnitud del torque en N$\cdot$m.
    \item Describir la dirección de rotación que produce el torque usando la regla de la mano derecha.
\end{enumerate}
    \item Dados los cuatro vértices de un tetraedro $A = (0,0,0)$, $B = (2,0,0)$, $C = (1,2,0)$ y $D = (1,1,2)$:
\begin{enumerate}
    \item Encontrar los vectores $\overrightarrow{AB}$, $\overrightarrow{AC}$ y $\overrightarrow{AD}$.
    \item Calcular el volumen del paralelepípedo definido por estos tres vectores usando el triple producto escalar.
    \item Sabiendo que el volumen de un tetraedro es $\frac{1}{6}$ del volumen del paralelepípedo que lo contiene, calcular el volumen del tetraedro $ABCD$.
\end{enumerate}
\end{enumerate}

% \section{Rectas en el espacio}



% \section{Planos en el Espacio}

Un plano en $\mathbb{R}^3$ se define como una superficie formada por >el conjunto de todos los puntos $(x, y, z)$ tales que la función $f(x, y, z) = ax + by + cz$ toma un valor constante $k$. 

Un plano se define de manera única si se conocen:
\begin{enumerate}
    \item Un punto de referencia $P_0(x_0, y_0, z_0)$ que pertenezca a la superficie.
    \item Un vector no nulo $\vec{n} = (a, b, c)$ que sea perpendicular a la superficie en todos sus puntos, denominado \textbf{vector normal}.
\end{enumerate}

\begin{center}
\shorthandoff{>}
\begin{tikzpicture}[x={(-0.5cm,-0.4cm)}, y={(1cm,0cm)}, z={(0cm,1cm)}, scale=1.5]
    % Dibujo del plano (con un poco de transparencia para estética 3D)
    \filldraw[fill=blue!10!white, draw=blue!60!black, thick, opacity=0.8] 
        (-2,-2,0) -- (4,-2,0) -- (5,3,0) -- (-1,3,0) -- cycle;
    
    % Coordenadas
    \coordinate (P0) at (0,0,0);
    \coordinate (P) at (3,0,0);
    \coordinate (N) at (0,0,2.5);
    
    % Vector Normal
    \draw[->, very thick, red] (P0) -- (N) node[above] {$\vec{n} = (a,b,c)$};
    
    % Vector sobre el plano
    \draw[->, very thick, blue] (P0) -- (P) node[midway, below right] {$\vec{P_0P}$};
    
    % Indicador de perpendicularidad
    \draw[thick] (0.3,0,0) -- (0.3,0,0.3) -- (0,0,0.3);
    
    % Puntos
    \filldraw[black] (P0) circle (1.5pt) node[left=3pt] {$P_0(x_0,y_0,z_0)$};
    \filldraw[black] (P) circle (1.5pt) node[right=3pt] {$P(x,y,z)$};
    
    % Etiqueta del plano
    \node[blue!80!black] at (-1, 2, 0) {Plano};
\end{tikzpicture}
\end{center}

\subsection{Ecuación escalar y general}

Si $P(x, y, z)$ es un punto arbitrario en el plano, el vector $\vec{P_0P} = (x-x_0, y-y_0, z-z_0)$ yace en la superficie. Por definición de vector normal, $\vec{n}$ debe ser perpendicular a $\vec{P_0P}$, lo que implica que su producto punto es $0$:

$$ \vec{n} \cdot \vec{P_0P} = 0 \implies (a, b, c) \cdot (x-x_0, y-y_0, z-z_0) = 0 $$

Al desarrollar el producto punto, se obtiene la ecuación escalar del plano:

$$ a(x - x_0) + b(y - y_0) + c(z - z_0) = 0 $$

Agrupando las constantes en un solo término $d = -(ax_0 + by_0 + cz_0)$, se llega a la ecuación general del plano:
$$ ax + by + cz + d = 0 $$

\subsection{Determinación de un Plano dados Tres Puntos}

Si se conocen tres puntos no colineales $A, B$ y $C$, es posible determinar la ecuación del plano que los contiene. \\

Primero, Se definen dos vectores directores que unan los puntos dados:
    $$ \vec{u} = \vec{AB} = B - A $$
    $$ \vec{v} = \vec{AC} = C - A $$

Puesto que $\vec{u}$ y $\vec{v}$ yacen sobre el plano, su producto vectorial generará un vector perpendicular a ambos y, por lo tanto, normal al plano:
    $$ \vec{n} = \vec{u} \times \vec{v} = \begin{vmatrix} \hat{i} & \hat{j} & \hat{k} \\ u_1 & u_2 & u_3 \\ v_1 & v_2 & v_3 \end{vmatrix} = (a, b, c) $$

Con el vector normal $(a, b, c)$ y utilizando el punto $A(x_1, y_1, z_1)$ como referencia, se aplica la ecuación escalar:
    $$ a(x-x_1) + b(y-y_1) + c(z-z_1) = 0 $$

\begin{center}
\shorthandoff{>}
\begin{tikzpicture}[x={(-0.5cm,-0.4cm)}, y={(1cm,0cm)}, z={(0cm,1cm)}, scale=1.5]
    % Superficie del plano
    \filldraw[fill=green!10!white, draw=green!60!black, thick, opacity=0.8] 
        (-1,-1.5,0) -- (4,-1.5,0) -- (5,3.5,0) -- (0,3.5,0) -- cycle;
    
    % Puntos A, B y C
    \coordinate (A) at (1,0,0);
    \coordinate (B) at (3,-1,0);
    \coordinate (C) at (2.5,2.5,0);
    
    % Vectores directores u y v
    \draw[->, very thick, blue] (A) -- (B) node[midway, above left] {$\vec{u} = \vec{AB}$};
    \draw[->, very thick, blue] (A) -- (C) node[midway, above right] {$\vec{v} = \vec{AC}$};
    
    % Vector Normal (Producto Cruz)
    \draw[->, very thick, red] (A) -- ++(0,0,2.5) node[above] {$\vec{n} = \vec{u} \times \vec{v}$};
    
    % Perpendicularidad
    \draw[thick] (A) ++(0,0,0.3) -- ++(0.3,-0.1,0) -- ++(0,0,-0.3);
    
    % Puntos marcados
    \filldraw[black] (A) circle (1.5pt) node[above right] {$A$};
    \filldraw[black] (B) circle (1.5pt) node[above left] {$B$};
    \filldraw[black] (C) circle (1.5pt) node[above right] {$C$};
\end{tikzpicture}
\end{center}

\subsection{Intersecciones Espaciales}

\subsubsection{Intersección de un Plano con una Recta}

Sea la recta $\vec{r}(t) = \vec{r}_0 + t\vec{v}$ y el plano $ax + by + cz + d = 0$.\\

Expresando las coordenadas de la recta en su forma paramétrica: $x = x_0 + at$, $y = y_0 + bt$, $z = z_0 + ct$.

Sustituyendo estas expresiones en la ecuación del plano:
    $$ a(x_0 + v_1 t) + b(y_0 + v_2 t) + c(z_0 + v_3 t) + d = 0 $$

Despejando el parámetro $t$:
    $$ t = \frac{-(ax_0 + by_0 + cz_0 + d)}{av_1 + bv_2 + cv_3} = \frac{-(\vec{n} \cdot \vec{r}_0 + d)}{\vec{n} \cdot \vec{v}} $$

Si $\vec{n} \cdot \vec{v} = 0$, la recta es paralela al plano, por lo tanto, no hay intersección o está contenida en él. Si hay un valor finito de $t$, se sustituye en la ecuación de la recta para hallar el punto exacto.

\begin{center}
\shorthandoff{>}
\begin{tikzpicture}[x={(-0.5cm,-0.4cm)}, y={(1cm,0cm)}, z={(0cm,1cm)}, scale=1.5]
    % Plano
    \filldraw[fill=orange!10!white, draw=orange!60!black, thick, opacity=0.8] 
        (-2,-2,0) -- (3,-2,0) -- (4,3,0) -- (-1,3,0) -- cycle;
    
    % Coordenadas de la recta
    \coordinate (L_start) at (1, 1, -2.5);
    \coordinate (L_int) at (1, 1, 0); % Punto de intersección
    \coordinate (L_end) at (1, 1, 3);
    
    % Recta debajo del plano (punteada para dar efecto de profundidad)
    \draw[dashed, very thick, black!60] (L_start) -- (L_int);
    
    % Punto de intersección
    \filldraw[red] (L_int) circle (2pt) node[above left, black] {Intersección};
    
    % Recta por encima del plano
    \draw[->, very thick, black] (L_int) -- (L_end) node[right] {Recta $\vec{r}(t)$};
    
    % Vector director de la recta
    \draw[->, thick, purple] (L_int) ++(0,0,1) -- ++(0,0,1.5) node[midway, right] {$\vec{v}$};
    
    % Vector normal del plano (como referencia visual)
    \draw[->, thick, orange] (0,2,0) -- (0,2,1.5) node[above] {$\vec{n}$};
\end{tikzpicture}
\end{center}

\subsubsection{Intersección entre dos Planos}

Dos planos no paralelos se intersectan siempre formando una línea recta. 

Para definir dicha recta $L$:
\begin{enumerate}
    \item La recta debe ser perpendicular a los dos vectores normales $\vec{n}_1$ y $\vec{n}_2$. Por lo tanto, su dirección es:
    $$ \vec{v}_L = \vec{n}_1 \times \vec{n}_2 $$
    \item Se resuelve el sistema de dos ecuaciones con tres incógnitas. Usualmente, se asigna un valor arbitrario a una variable y se resuelve el sistema de $2 \times 2$ resultante para $x$ e $y$ para encontrar un punto en común en ambos planos.
\end{enumerate}

\begin{center}
\shorthandoff{>}
\begin{tikzpicture}[scale=1.5, x={(-0.5cm,-0.4cm)}, y={(1cm,0cm)}, z={(0cm,1cm)}]
    % Plano 1 (Horizontal, dibujado mitad atrás)
    \filldraw[fill=blue!10!white, draw=blue!50!white, thick, opacity=0.7] 
        (-2,-2,0) -- (0,-2,0) -- (0,3,0) -- (-2,3,0) -- cycle;
        
    % Plano 2 (Vertical-inclinado)
    \filldraw[fill=red!10!white, draw=red!50!white, thick, opacity=0.7] 
        (0,-2,-2) -- (0,3,-2) -- (0,3,2) -- (0,-2,2) -- cycle;
        
    % Plano 1 (Horizontal, dibujado mitad adelante para efecto 3D)
    \filldraw[fill=blue!10!white, draw=blue!50!white, thick, opacity=0.7] 
        (0,-2,0) -- (3,-2,0) -- (3,3,0) -- (0,3,0) -- cycle;

    % Línea de intersección
    \draw[very thick, purple] (0,-2,0) -- (0,3,0) node[above right] {$L$};
    
    % Vectores Normales
    \draw[->, very thick, blue] (0,0.5,0) -- (0,0.5,1.5) node[above] {$\vec{n}_1$};
    \draw[->, very thick, red] (0,0.5,0) -- (1.5,0.5,0) node[right] {$\vec{n}_2$};
    
    % Vector Director de la recta de intersección
    \draw[->, very thick, purple] (0,1.5,0) -- (0,2.5,0) node[below right] {$\vec{v}_L = \vec{n}_1 \times \vec{n}_2$};
\end{tikzpicture}
\end{center}

\subsection{Distancia de un Punto a un Plano}

Sea un punto exterior $S(x_1, y_1, z_1)$ y un plano definido por $ax + by + cz + d = 0$. La distancia mínima $D$ entre el punto y el plano es la longitud del segmento perpendicular desde $S$ hasta el plano.

Definase $P(x_0, y_0, z_0)$ como un punto cualquiera que pertenezca al plano. Se define el vector $\vec{PS} = (x_1-x_0, y_1-y_0, z_1-z_0)$.
La distancia $D$ es la magnitud de la proyección escalar de $\vec{PS}$ sobre el vector normal $\vec{n}$:
$$ D = |\text{proy}_{\vec{n}} \vec{PS}| = \frac{|\vec{PS} \cdot \vec{n}|}{\|\vec{n}\|} $$



Al desarrollar el producto punto en el numerador:
$$ \vec{PS} \cdot \vec{n} = a(x_1-x_0) + b(y_1-y_0) + c(z_1-z_0) = ax_1 + by_1 + cz_1 - (ax_0 + by_0 + cz_0) $$
Sustituyendo la condición del punto $P$ ($ax_0 + by_0 + cz_0 = -d$):
$$ \vec{PS} \cdot \vec{n} = ax_1 + by_1 + cz_1 - (-d) = ax_1 + by_1 + cz_1 + d $$

\begin{mydefinition}{Fórmula de Distancia Punto-Plano}
La distancia $D$ del punto $S(x_1, y_1, z_1)$ al plano $ax + by + cz + d = 0$ está dada por:
$$ D = \frac{|ax_1 + by_1 + cz_1 + d|}{\sqrt{a^2 + b^2 + c^2}} $$
\end{mydefinition}

\begin{center}
\shorthandoff{>}
\begin{tikzpicture}[x={(-0.5cm,-0.4cm)}, y={(1cm,0cm)}, z={(0cm,1cm)}, scale=1.5]
    % Superficie del plano
    \filldraw[fill=cyan!10!white, draw=cyan!60!black, thick, opacity=0.8] 
        (-2,-2,0) -- (3,-2,0) -- (4,3,0) -- (-1,3,0) -- cycle;
        
    % Puntos clave
    \coordinate (P) at (0,0,0);     
    \coordinate (S) at (2, 2.5, 2); 
    \coordinate (ProjN) at (0, 0, 2); 
    \coordinate (Proj) at (2,2.5,0);
    
    \draw[->, thick, red] (P) -- (0,0,4.2) node[left] {$\vec{n}$};
    
    \draw[->, line width=2.5pt, purple] (P) -- (ProjN) node[midway, left=2pt] {$D = |\text{proy}_{\vec{n}}\vec{PS}|$};
    
    \draw[->, very thick, blue] (P) -- (S) node[midway, below right] {$\vec{PS}$};
    
    \draw[dashed, very thick, purple] (S) -- (ProjN);
    \draw[dashed, very thick, black] (S) -- (Proj);
    
    \draw (0, 0, 0.3) -- (0, -0.3, 0.3) -- (0, -0.3, 0);
    
    \filldraw[black] (P) circle (1.5pt) node[below left] {$P(x_0,y_0,z_0)$};
    \filldraw[black] (S) circle (1.5pt) node[above right] {$S(x_1,y_1,z_1)$};
    \filldraw[black] (ProjN) circle (1.5pt);
\end{tikzpicture}
\end{center}

% \section{Superficies cilíndricas}

\subsection{Definición}

Una superficie cilíndrica en el contexto del cálculo multivariable tiene un significado más general que el cilindro circular de la geometría básica. Engloba cualquier superficie generada por rectas paralelas que barren una curva.

\begin{mydefinition}{Superficie cilíndrica}
Una superficie cilíndrica es una superficie en $\mathbb{R}^3$ formada por todas las rectas generatrices paralelas a un eje fijo y que pasan a través de una curva plana fija denominada curva directriz.
\end{mydefinition}

Ejemplos de identificación:
\begin{itemize}
    \item $x^2 + y^2 = 1$: falta $z$ $\Rightarrow$ cilindro con eje paralelo al eje $z$ (sección circular).
    \item $z = x^2$: falta $y$ $\Rightarrow$ cilindro con eje paralelo al eje $y$ (sección parabólica).
    \item $y^2 - x^2 = 1$: falta $z$ $\Rightarrow$ cilindro con eje paralelo al eje $z$ (sección hiperbólica).
\end{itemize}

\subsection{Cómo graficar una superficie cilíndrica}

\begin{enumerate}
    \item \textbf{Identificar la variable ausente.} Esta indica la dirección de las generatrices.
    \item \textbf{Trazar la curva directriz.} Dibujar la curva en el plano formado por las otras dos variables.
    \item \textbf{Extender.} Extender la curva directriz en ambas direcciones a lo largo del eje de la variable ausente.
\end{enumerate}

\subsubsection{Cilindro circular: $x^2 + y^2 = 1$}

La variable $z$ está ausente. La curva directriz en el plano $xy$ es la circunferencia unitaria. Al extruir a lo largo del eje $z$, obtenemos el cilindro circular recto.

\begin{center}
\begin{tikzpicture}
\begin{axis}[
    view={60}{20},
    axis lines=center,
    xlabel=$x$, ylabel=$y$, zlabel=$z$,
    xmin=-1.5, xmax=1.5,
    ymin=-1.5, ymax=1.5,
    zmin=-2, zmax=2,
    colormap/viridis,
    title={Cilindro circular: $x^2 + y^2 = 1$}
]
\addplot3[
    surf,
    parametric,
    domain=0:360,
    domain y=-2:2,
    samples=40,
    samples y=5,
    opacity=0.8
] ({cos(x)}, {sin(x)}, {y});
\end{axis}
\end{tikzpicture}
\end{center}

\subsubsection{Cilindro parabólico: $z = x^2$}

La variable $y$ está ausente. La curva directriz en el plano $xz$ es la parábola $z = x^2$. Al extruir a lo largo del eje $y$, obtenemos el cilindro parabólico.

\begin{center}
\begin{tikzpicture}
\begin{axis}[
    view={60}{20},
    axis lines=center,
    xlabel=$x$, ylabel=$y$, zlabel=$z$,
    colormap/viridis,
    title={Cilindro parabólico: $z = x^2$}
]
\addplot3[
    surf,
    domain=-2:2,
    domain y=-2:2,
    samples=20,
    samples y=5,
    opacity=0.8
] {x^2};
\end{axis}
\end{tikzpicture}
\end{center}


\subsection{Ejercicios}

\begin{enumerate}
    \item Identificar el tipo de cada superficie cilíndrica, señalar la variable ausente, describir la curva directriz y el eje a lo largo del cual se extiende la superficie.
\begin{enumerate}
    \item $x^2 + z^2 = 9$
    \item $y = \sin(x)$
    \item $4x^2 + 9y^2 = 36$
    \item $z = e^x$
\end{enumerate}
    \item Determinar la ecuación del cilindro cuyas generatrices son paralelas al eje $z$ y cuya curva directriz en el plano $xy$ es la elipse con semiejes $a = 3$ en la dirección $x$ y $b = 2$ en la dirección $y$, centrada en el origen.

Verificar que los puntos $P_1 = (3, 0, 5)$ y $P_2 = (0, 2, -7)$ pertenecen a la superficie, y determinar si $P_3 = (2, 1, 4)$ pertenece o no.
    \item Para el cilindro parabólico $y = x^2 - 4x + 3$:
\begin{enumerate}
    \item Completar el cuadrado para escribir la ecuación en la forma $y = (x - h)^2 + k$ e identificar el vértice de la parábola directriz.
    \item Señalar cuál es la variable ausente y describir el eje del cilindro.
    \item Encontrar la intersección del cilindro con el plano $z = 2$ y con el plano $x = 1$.
    \item Determinar si el punto $Q = (2, -1, 100)$ pertenece al cilindro.
\end{enumerate}
    \item La recta $L: \vec{r}(t) = (1,\, 0,\, -1) + t(0,\, 1,\, 3)$ y el cilindro circular $x^2 + y^2 = 5$.
\begin{enumerate}
    \item Determinar si la recta es paralela, interseca o está contenida en el cilindro.
    \item Si la recta interseca al cilindro, encontrar todos los puntos de intersección.
    \item Encontrar la distancia del eje del cilindro (el eje $z$) a la recta $L$.
\end{enumerate}

\textit{Sugerencia:} Sustituir las expresiones paramétricas de $x$ e $y$ en la ecuación del cilindro y resolver en $t$.
    \item Encontrar la ecuación del cilindro con generatrices paralelas al eje $y$ que contiene a la curva de intersección de las superficies $z = x^2 + 1$ y $z = 4 - x^2$.

\begin{enumerate}
    \item Encontrar las curvas de intersección igualando las dos ecuaciones.
    \item Eliminar la variable $z$ para obtener la ecuación del cilindro en el plano $xy$.
    \item Describir geométricamente la curva directriz obtenida.
\end{enumerate}
    \item Escribir las ecuaciones paramétricas de cada superficie cilíndrica e identificar el rango de los parámetros para cubrir la superficie completa.
\begin{enumerate}
    \item Cilindro circular $x^2 + z^2 = 4$ (eje $y$), para $-3 \leq y \leq 3$.
    \item Cilindro elíptico $\dfrac{x^2}{9} + \dfrac{z^2}{4} = 1$ (eje $y$), para $y \in \mathbb{R}$.
    \item Cilindro parabólico $z = x^2$ (eje $y$), para $-2 \leq x \leq 2$, $-2 \leq y \leq 2$.
\end{enumerate}

\textit{Sugerencia:} Para cónicas en el plano $xz$, usar una parametrización trigonométrica o explícita según corresponda, y agregar $y = t$ como parámetro libre.
    \item Clasificar cada ecuación como cilindro o no. Si es un cilindro, identificar el tipo y la variable ausente. Si no lo es, justificar brevemente.
\begin{enumerate}
    \item $x^2 + y^2 + z^2 = 16$
    \item $x^2 - y^2 = 1$
    \item $z = \ln(x^2 + 1)$
    \item $xy = 4$
    \item $x^2 + y^2 + z = 0$
    \item $x^2 + z^2 - 4x + 2z - 4 = 0$
\end{enumerate}

Para el inciso (f), completar el cuadrado en las variables correspondientes y escribir la ecuación en forma estándar.
    \item El cilindro hiperbólico $x^2 - y^2 = 1$ tiene generatrices paralelas al eje $z$.
\begin{enumerate}
    \item Identificar los vértices y asíntotas de la hipérbola directriz en el plano $xy$.
    \item Determinar la traza del cilindro con el plano $z = k$ para cualquier valor de $k$.
    \item Encontrar la intersección del cilindro con el plano $y = 2$ y describir qué curva se obtiene.
    \item Verificar que el punto $P = (\sqrt{2},\, 1,\, -3)$ pertenece al cilindro.
    \item Encontrar todos los puntos del cilindro con $x = 2$, $z = 0$.
\end{enumerate}
    \item Determinar la ecuación del cilindro con generatrices paralelas al eje $z$ que pasa por la curva $C$ definida por la intersección de:
\[
x^2 + y^2 + z^2 = 9 \quad \text{y} \quad z = 1
\]

\begin{enumerate}
    \item Encontrar la curva $C$ sustituyendo $z = 1$ en la ecuación de la esfera.
    \item Eliminar la variable $z$ para obtener la ecuación del cilindro proyectado sobre el plano $xy$.
    \item Describir geométricamente la superficie cilíndrica obtenida e indicar su radio.
    \item Encontrar la intersección del cilindro obtenido con el plano $y = 0$.
\end{enumerate}
    \item La región sólida $E$ está acotada por el cilindro circular $x^2 + y^2 = 4$, el plano $z = 0$ y el plano $z = 3$.
\begin{enumerate}
    \item Describir en palabras la forma del sólido $E$.
    \item Determinar si el punto $Q = (1,\, 1,\, 2)$ está dentro, fuera o sobre la frontera de $E$.
    \item Encontrar la longitud del segmento de la recta $\vec{r}(t) = (1,\, 0,\, t)$, $t \in \mathbb{R}$, que se encuentra dentro de $E$.
    \item Dos cilindros, $x^2 + y^2 = 4$ y $x^2 + z^2 = 4$, se intersectan. Encontrar y describir la curva de intersección en la región $x \geq 0$.

    \textit{Sugerencia:} En la curva de intersección se cumple $y^2 = z^2$; usar esto para parametrizar.
\end{enumerate}
\end{enumerate}


% 
\section{Superficies cuádricas}

\subsection{Definición y forma general}

Las superficies cuádricas son la extensión natural de las cónicas (elipse, parábola, hipérbola) al espacio tridimensional.

\begin{mydefinition}{Superficie cuádrica}
Una superficie cuádrica es el conjunto de puntos $(x, y, z) \in \mathbb{R}^3$ que satisfacen una ecuación de la forma:
\[
Ax^2 + By^2 + Cz^2 + Dxy + Exz + Fyz + Gx + Hy + Iz + J = 0
\]
donde $A, B, \ldots, J \in \mathbb{R}$ y al menos uno de los coeficientes cuadráticos es distinto de cero.
\end{mydefinition}

Aplicando traslaciones y rotaciones de los ejes, toda cuádrica no degenerada se reduce a una de las seis formas estándar que se estudian a continuación.

\subsection{Método de las trazas}

\begin{mydefinition}{Traza de una superficie}
La traza de una superficie en un plano $P$ es la curva de intersección entre ambos. Las trazas más útiles son:
\begin{itemize}
    \item $z = k$: intersección con un plano horizontal.
    \item $y = k$: intersección con un plano paralelo al plano $xz$.
    \item $x = k$: intersección con un plano paralelo al plano $yz$.
\end{itemize}
Al variar $k$, se observa cómo cambia la sección transversal, lo que permite reconstruir la forma tridimensional de la superficie.
\end{mydefinition}

\subsection{Elipsoide}

\begin{mydefinition}{Elipsoide}
La ecuación estándar del elipsoide centrado en el origen es:
\[
\frac{x^2}{a^2} + \frac{y^2}{b^2} + \frac{z^2}{c^2} = 1
\]
Los valores $a$, $b$ y $c$ son los semiejes en las direcciones $x$, $y$ y $z$ respectivamente. Si $a = b = c$, el elipsoide es una esfera de radio $a$.
\end{mydefinition}

\textbf{Trazas:} todas las trazas del elipsoide son elipses (o circunferencias en casos especiales).
\begin{itemize}
    \item $z = k$ (con $|k| < c$): $\;\dfrac{x^2}{a^2\left(1-\frac{k^2}{c^2}\right)} + \dfrac{y^2}{b^2\left(1-\frac{k^2}{c^2}\right)} = 1$ \quad (elipse)
    \item $y = 0$: $\;\dfrac{x^2}{a^2} + \dfrac{z^2}{c^2} = 1$ \quad (elipse en el plano $xz$)
    \item $x = 0$: $\;\dfrac{y^2}{b^2} + \dfrac{z^2}{c^2} = 1$ \quad (elipse en el plano $yz$)
\end{itemize}

El elipsoide es la única cuádrica acotada: todos sus puntos están dentro de la caja $[-a,a]\times[-b,b]\times[-c,c]$.

\begin{center}
\begin{tikzpicture}
\begin{axis}[
    view={60}{20},
    axis lines=center,
    xlabel=$x$, ylabel=$y$, zlabel=$z$,
    colormap/viridis,
    title={Elipsoide: $\frac{x^2}{4} + y^2 + \frac{z^2}{4} = 1$\quad ($a=2,\,b=1,\,c=2$)}
]
\addplot3[
    surf,
    parametric,
    domain=0:360,
    domain y=0:180,
    samples=40,
    samples y=20,
    opacity=0.85
] ({2*sin(y)*cos(x)}, {sin(y)*sin(x)}, {2*cos(y)});
\end{axis}
\end{tikzpicture}
\end{center}

\subsection{Paraboloide elíptico}

\begin{mydefinition}{Paraboloide elíptico}
La ecuación estándar del paraboloide elíptico con vértice en el origen es:
\[
z = \frac{x^2}{a^2} + \frac{y^2}{b^2}
\]
La superficie se abre hacia el lado positivo del eje $z$. Si $a = b$, se obtiene el paraboloide circular (de revolución).
\end{mydefinition}

\textbf{Trazas:}
\begin{itemize}
    \item $z = k > 0$: $\;\dfrac{x^2}{a^2 k} + \dfrac{y^2}{b^2 k} = 1$ \quad (elipse que crece con $k$)
    \item $z = 0$: solo el vértice en el origen
    \item $y = 0$: $\;z = \dfrac{x^2}{a^2}$ \quad (parábola en el plano $xz$)
    \item $x = 0$: $\;z = \dfrac{y^2}{b^2}$ \quad (parábola en el plano $yz$)
\end{itemize}

\begin{center}
\begin{tikzpicture}
\begin{axis}[
    view={60}{20},
    axis lines=center,
    xlabel=$x$, ylabel=$y$, zlabel=$z$,
    colormap/viridis,
    zmin=0,
    title={Paraboloide elíptico: $z = x^2 + y^2$}
]
\addplot3[
    surf,
    domain=-2:2,
    domain y=-2:2,
    samples=25,
    opacity=0.85
] {x^2 + y^2};
\end{axis}
\end{tikzpicture}
\end{center}

\subsection{Paraboloide hiperbólico}

\begin{mydefinition}{Paraboloide hiperbólico}
La ecuación estándar del paraboloide hiperbólico con punto de silla en el origen es:
\[
z = \frac{y^2}{b^2} - \frac{x^2}{a^2}
\]
\end{mydefinition}

Esta superficie tiene la característica forma de silla de montar: desde el origen, sube en la dirección $y$ y baja en la dirección $x$.

\textbf{Trazas:}
\begin{itemize}
    \item $z = k > 0$: $\;\dfrac{y^2}{b^2 k} - \dfrac{x^2}{a^2 k} = 1$ \quad (hipérbola abierta en dirección $y$)
    \item $z = k < 0$: $\;\dfrac{x^2}{a^2|k|} - \dfrac{y^2}{b^2|k|} = 1$ \quad (hipérbola abierta en dirección $x$)
    \item $z = 0$: $\;y/b = \pm\, x/a$ \quad (dos rectas que se intersectan --- punto de silla)
    \item $y = 0$: $\;z = -\dfrac{x^2}{a^2}$ \quad (parábola hacia abajo)
    \item $x = 0$: $\;z = \dfrac{y^2}{b^2}$ \quad (parábola hacia arriba)
\end{itemize}

\begin{center}
\begin{tikzpicture}
\begin{axis}[
    view={60}{20},
    axis lines=center,
    xlabel=$x$, ylabel=$y$, zlabel=$z$,
    colormap/viridis,
    title={Paraboloide hiperbólico: $z = y^2 - x^2$}
]
\addplot3[
    surf,
    domain=-2:2,
    domain y=-2:2,
    samples=25,
    opacity=0.85
] {y^2 - x^2};
\end{axis}
\end{tikzpicture}
\end{center}

\subsection{Cono elíptico}

\begin{mydefinition}{Cono elíptico}
La ecuación estándar del cono elíptico con vértice en el origen es:
\[
\frac{x^2}{a^2} + \frac{y^2}{b^2} = \frac{z^2}{c^2}
\]
Cuando $a = b$, se obtiene el cono circular recto. La superficie consta de dos napas simétricas respecto al origen.
\end{mydefinition}

\textbf{Trazas:}
\begin{itemize}
    \item $z = k \neq 0$: $\;\dfrac{x^2}{\left(\frac{ak}{c}\right)^2} + \dfrac{y^2}{\left(\frac{bk}{c}\right)^2} = 1$ \quad (elipse que se expande con $|k|$)
    \item $z = 0$: solo el vértice en el origen
    \item $y = 0$: $\;z = \pm\,\dfrac{c}{a}x$ \quad (dos rectas que se cruzan)
    \item $x = 0$: $\;z = \pm\,\dfrac{c}{b}y$ \quad (dos rectas que se cruzan)
\end{itemize}

\begin{center}
\begin{tikzpicture}
\begin{axis}[
    view={60}{20},
    axis lines=center,
    xlabel=$x$, ylabel=$y$, zlabel=$z$,
    colormap/viridis,
    title={Cono: $z^2 = x^2 + y^2$}
]
\addplot3[
    surf,
    parametric,
    domain=0:360,
    domain y=0:2,
    samples=40,
    samples y=10,
    opacity=0.85
] ({y*cos(x)}, {y*sin(x)}, {y});
\addplot3[
    surf,
    parametric,
    domain=0:360,
    domain y=0:2,
    samples=40,
    samples y=10,
    opacity=0.85
] ({y*cos(x)}, {y*sin(x)}, {-y});
\end{axis}
\end{tikzpicture}
\end{center}

\subsection{Hiperboloide de una hoja}

\begin{mydefinition}{Hiperboloide de una hoja}
La ecuación estándar del hiperboloide de una hoja con eje a lo largo de $z$ es:
\[
\frac{x^2}{a^2} + \frac{y^2}{b^2} - \frac{z^2}{c^2} = 1
\]
\end{mydefinition}

El hiperboloide de una hoja es una superficie conexa (una sola pieza). Su sección transversal horizontal es siempre una elipse, con tamaño mínimo en $z = 0$ que crece conforme $|z|$ aumenta.

\textbf{Trazas:}
\begin{itemize}
    \item $z = k$: $\;\dfrac{x^2}{a^2\left(1+\frac{k^2}{c^2}\right)} + \dfrac{y^2}{b^2\left(1+\frac{k^2}{c^2}\right)} = 1$ \quad (elipse mínima en $k=0$)
    \item $y = 0$: $\;\dfrac{x^2}{a^2} - \dfrac{z^2}{c^2} = 1$ \quad (hipérbola)
    \item $x = 0$: $\;\dfrac{y^2}{b^2} - \dfrac{z^2}{c^2} = 1$ \quad (hipérbola)
\end{itemize}

\begin{mynote}{Relación con el cono}
El cono y el hiperboloide de una hoja solo difieren en el lado derecho de la expresión $\dfrac{x^2}{a^2}+\dfrac{y^2}{b^2}-\dfrac{z^2}{c^2}$: es $0$ para el cono y $1$ para el hiperboloide. Las ``paredes'' del hiperboloide se acercan asintóticamente al cono cuando $|z|\to\infty$.
\end{mynote}

\begin{center}
\begin{tikzpicture}
\begin{axis}[
    view={60}{20},
    axis lines=center,
    xlabel=$x$, ylabel=$y$, zlabel=$z$,
    colormap/viridis,
    title={Hiperboloide de una hoja: $x^2 + y^2 - z^2 = 1$}
]
\addplot3[
    surf,
    parametric,
    domain=0:360,
    domain y=-1.5:1.5,
    samples=40,
    samples y=20,
    opacity=0.85
] ({sqrt(1+y^2)*cos(x)}, {sqrt(1+y^2)*sin(x)}, {y});
\end{axis}
\end{tikzpicture}
\end{center}

\subsection{Hiperboloide de dos hojas}

\begin{mydefinition}{Hiperboloide de dos hojas}
La ecuación estándar del hiperboloide de dos hojas con eje a lo largo de $z$ es:
\[
\frac{z^2}{c^2} - \frac{x^2}{a^2} - \frac{y^2}{b^2} = 1
\]
\end{mydefinition}

Para que la ecuación tenga solución real se requiere $z^2/c^2 \geq 1$, es decir, $|z| \geq c$. La superficie no existe en la región $-c < z < c$ y consta de dos componentes separadas: una hoja para $z \geq c$ y otra para $z \leq -c$.

\textbf{Trazas:}
\begin{itemize}
    \item $z = k$ con $|k| > c$: $\;\dfrac{x^2}{a^2\!\left(\frac{k^2}{c^2}-1\right)} + \dfrac{y^2}{b^2\!\left(\frac{k^2}{c^2}-1\right)} = 1$ \quad (elipse)
    \item $|z| < c$: no existe traza (región vacía)
    \item $y = 0$: $\;\dfrac{z^2}{c^2} - \dfrac{x^2}{a^2} = 1$ \quad (hipérbola abierta en dirección $z$)
    \item $x = 0$: $\;\dfrac{z^2}{c^2} - \dfrac{y^2}{b^2} = 1$ \quad (hipérbola abierta en dirección $z$)
\end{itemize}

\begin{center}
\begin{tikzpicture}
\begin{axis}[
    view={60}{20},
    axis lines=center,
    xlabel=$x$, ylabel=$y$, zlabel=$z$,
    colormap/viridis,
    title={Hiperboloide de dos hojas: $z^2 - x^2 - y^2 = 1$}
]
\addplot3[
    surf,
    domain=-2:2,
    domain y=-2:2,
    samples=25,
    opacity=0.85
] {sqrt(1+x^2+y^2)};
\addplot3[
    surf,
    domain=-2:2,
    domain y=-2:2,
    samples=25,
    opacity=0.85
] {-sqrt(1+x^2+y^2)};
\end{axis}
\end{tikzpicture}
\end{center}

\subsection{Resumen de superficies cuádricas}

\begin{center}
\renewcommand{\arraystretch}{2.2}
\begin{tabular}{|l|c|l|}
\hline
\textbf{Superficie} & \textbf{Ecuación estándar (eje $z$)} & \textbf{Trazas características} \\
\hline
Elipsoide & $\dfrac{x^2}{a^2}+\dfrac{y^2}{b^2}+\dfrac{z^2}{c^2}=1$ & Todas elipses \\
\hline
Paraboloide elíptico & $z=\dfrac{x^2}{a^2}+\dfrac{y^2}{b^2}$ & Elipses (horiz.) / Parábolas (vert.) \\
\hline
Paraboloide hiperbólico & $z=\dfrac{y^2}{b^2}-\dfrac{x^2}{a^2}$ & Hipérbolas / Parábolas \\
\hline
Cono elíptico & $\dfrac{x^2}{a^2}+\dfrac{y^2}{b^2}=\dfrac{z^2}{c^2}$ & Elipses / Rectas cruzadas \\
\hline
Hiperboloide de una hoja & $\dfrac{x^2}{a^2}+\dfrac{y^2}{b^2}-\dfrac{z^2}{c^2}=1$ & Elipses / Hipérbolas \\
\hline
Hiperboloide de dos hojas & $\dfrac{z^2}{c^2}-\dfrac{x^2}{a^2}-\dfrac{y^2}{b^2}=1$ & Elipses (si $|z|>c$) / Hipérbolas \\
\hline
\end{tabular}
\end{center}


\subsection{Ejercicios}

\begin{enumerate}
    \item Identificar el tipo de superficie cuádrica de cada ecuación (elipsoide, paraboloide elíptico, paraboloide hiperbólico, cono, hiperboloide de una hoja o hiperboloide de dos hojas). Justificar la clasificación.
\begin{enumerate}
    \item $\dfrac{x^2}{4} + \dfrac{y^2}{9} + \dfrac{z^2}{1} = 1$
    \item $z = x^2 + 4y^2$
    \item $x^2 + y^2 - z^2 = 1$
    \item $\dfrac{x^2}{4} + \dfrac{y^2}{9} - \dfrac{z^2}{16} = 0$
    \item $z^2 - x^2 - y^2 = 4$
    \item $z = y^2 - x^2$
\end{enumerate}
    \item Completar el cuadrado en las variables correspondientes para llevar cada ecuación a su forma estándar. Identificar el tipo de cuádrica, el centro y los semiejes o vértice según corresponda.
\begin{enumerate}
    \item $4x^2 + y^2 + 9z^2 - 8x + 2y - 36z + 37 = 0$
    \item $z = x^2 - 4x + y^2 + 6y + 13$
    \item $x^2 + y^2 - z^2 - 2x + 4y + 2z - 2 = 0$
\end{enumerate}

\textit{Sugerencia:} Para el inciso (a), completar el cuadrado en $x$, $y$ y $z$ por separado y verificar si el lado derecho resulta positivo.
    \item Para el elipsoide $\dfrac{x^2}{9} + \dfrac{y^2}{4} + \dfrac{z^2}{16} = 1$:
\begin{enumerate}
    \item Identificar los semiejes $a$, $b$ y $c$ en las direcciones $x$, $y$ y $z$.
    \item Encontrar las trazas con los tres planos coordenados ($z=0$, $y=0$, $x=0$) y describir cada cónica obtenida.
    \item Encontrar la traza con el plano $z = 2$ y determinar sus semiejes.
    \item Determinar si los puntos $P_1 = (3, 0, 0)$, $P_2 = (1, 1, 2)$ y $P_3 = (0, 2, 4)$ están sobre, dentro o fuera del elipsoide.
    \item ¿Para qué valores de $k$ el plano $z = k$ interseca al elipsoide?
\end{enumerate}
    \item El paraboloide elíptico $z = \dfrac{x^2}{4} + \dfrac{y^2}{9}$ tiene su vértice en el origen.
\begin{enumerate}
    \item Encontrar la traza con el plano $z = 1$ y calcular el área de la elipse obtenida usando $A = \pi ab$.
    \item Encontrar las trazas con los planos $y = 0$ y $x = 0$. ¿Qué tipo de curva es cada una?
    \item Determinar el punto de la superficie más cercano al origen. ¿Cuál es ese punto?
    \item El plano $z = 4$ corta al paraboloide. Encontrar la ecuación de la elipse de intersección y calcular su área.
    \item Describir cómo cambia la forma de las trazas horizontales $z = k$ al aumentar $k$ desde $0$ hasta $\infty$.
\end{enumerate}
    \item El paraboloide hiperbólico $z = x^2 - y^2$ tiene un punto de silla en el origen.
\begin{enumerate}
    \item Encontrar las trazas con los planos $z = 1$, $z = -1$ y $z = 0$, e identificar las cónicas obtenidas.
    \item Las trazas con $y = 0$ y $x = 0$ son parábolas. Determinar en qué dirección se abren.
    \item Verificar que el origen es un punto de silla: mostrar que a lo largo de la dirección $x$ el origen es un mínimo local de $z$, y a lo largo de la dirección $y$ es un máximo local.
    \item Encontrar todos los puntos de la superficie con $z = 0$ (la curva del ``ecuador'' de la silla). ¿Qué rectas son estas?
    \item Determinar si el punto $P = (3, 2, 5)$ pertenece a la superficie.
\end{enumerate}
    \item Dado el cono elíptico $\dfrac{x^2}{4} + \dfrac{y^2}{9} = z^2$:
\begin{enumerate}
    \item Verificar que el vértice se encuentra en el origen.
    \item Encontrar las trazas con los planos $z = 1$, $z = 2$ y $z = -1$. Identificar qué tipo de curva es cada una e indicar los semiejes.
    \item Encontrar las trazas con $y = 0$ y $x = 0$ y mostrar que son pares de rectas.
    \item ¿Cuántas napas (mitades) tiene el cono? Describir la napa superior ($z > 0$) y la inferior ($z < 0$).
    \item El plano $z = k$ interseca al cono en una elipse para todo $k \neq 0$. Escribir la fórmula de los semiejes en función de $k$ y calcular el área de la sección transversal.
\end{enumerate}
    \item Para el hiperboloide de una hoja $\dfrac{x^2}{4} + \dfrac{y^2}{4} - \dfrac{z^2}{9} = 1$:
\begin{enumerate}
    \item Identificar los semiejes $a$, $b$, $c$ y describir la orientación de la superficie.
    \item Encontrar la traza con $z = 0$ (la ``cintura'' del hiperboloide) e identificar su tipo y radio.
    \item Encontrar la traza con $z = 3$ y comparar su tamaño con la traza en $z = 0$.
    \item Encontrar las trazas con $y = 0$ y $x = 0$. ¿Qué tipo de cónica es cada una?
    \item Determinar para qué valores de $k$ el plano $z = k$ produce una traza válida (es decir, siempre existe intersección). Comparar con el hiperboloide de dos hojas.
\end{enumerate}
    \item Para el hiperboloide de dos hojas $\dfrac{z^2}{4} - \dfrac{x^2}{9} - \dfrac{y^2}{9} = 1$:
\begin{enumerate}
    \item Determinar los valores de $z$ para los cuales existe la superficie. ¿Cuál es la región vacía?
    \item Encontrar los vértices del hiperboloide (los puntos de las hojas más cercanos al origen).
    \item Encontrar la traza con el plano $z = 4$ e identificar su tipo y semiejes.
    \item Encontrar las trazas con $y = 0$ y $x = 0$. ¿Qué hipérbolas se obtienen?
    \item Comparar las ecuaciones del hiperboloide de una hoja y del de dos hojas del ejercicio anterior. ¿Qué diferencia hay en los signos? ¿Cómo afecta esto a la topología (número de componentes conexas)?
\end{enumerate}
    \item Una superficie cuádrica tiene la ecuación general:
\[
9x^2 + 4y^2 - 36z^2 - 18x + 16y + 72z - 11 = 0
\]

\begin{enumerate}
    \item Completar el cuadrado en cada variable para llevar la ecuación a su forma estándar.
    \item Identificar el tipo de superficie cuádrica e indicar el centro.
    \item Encontrar los semiejes de la superficie.
    \item Describir las trazas con cada plano coordenado trasladado al nuevo centro.
\end{enumerate}

\textit{Sugerencia:} Agrupar los términos en $x$, en $y$ y en $z$ por separado antes de completar el cuadrado.
    \item Encontrar la ecuación de la superficie cuádrica que satisface cada descripción. Escribir la ecuación en forma estándar.
\begin{enumerate}
    \item \textbf{Elipsoide} centrado en el origen que pasa por los puntos $(4, 0, 0)$, $(0, 2, 0)$ y $(0, 0, 3)$.
    \item \textbf{Paraboloide elíptico} con vértice en el origen, eje a lo largo de $z$ positivo, que pasa por el punto $(2, 1, 2)$.
    \item \textbf{Cono circular} con vértice en el origen y eje a lo largo de $z$, tal que la traza con el plano $z = 1$ es la circunferencia $x^2 + y^2 = 3$.
    \item \textbf{Hiperboloide de una hoja} con eje a lo largo de $z$, semiejes horizontales $a = b = 2$ y tal que la traza con $z = 2$ tenga semiejes $a' = b' = 2\sqrt{2}$.
    \item \textbf{Hiperboloide de dos hojas} con eje a lo largo de $z$ y vértices en $(0, 0, \pm 3)$, con semiejes horizontales $a = b = 2$.
\end{enumerate}
\end{enumerate}


% \section{Funciones Vectoriales}

Una función vectorial generaliza el concepto de función real: en lugar de asignar un número real a cada valor del parámetro, asigna un \textbf{vector}. Son la herramienta fundamental para describir curvas y trayectorias en el espacio, y conectan el cálculo univariable con el análisis multivariable.

\subsection{Definición y Dominio}

\begin{mydefinition}{Función vectorial}
Una función vectorial en $\mathbb{R}^n$ es una función $\vec{r}: D \subseteq \mathbb{R} \to \mathbb{R}^n$ que asigna a cada escalar $t \in D$ un único vector en $\mathbb{R}^n$:
\[ \vec{r}(t) = \bigl(f_1(t),\; f_2(t),\; \dots,\; f_n(t)\bigr) = f_1(t)\,\hat{\imath}_1 + f_2(t)\,\hat{\imath}_2 + \dots + f_n(t)\,\hat{\imath}_n \]
Las funciones reales $f_1, f_2, \dots, f_n$ se llaman funciones componentes de $\vec{r}$.
\end{mydefinition}

El dominio natural de $\vec{r}$ es la intersección de los dominios de sus componentes:
\[ D = \bigcap_{i=1}^{n} \text{Dom}(f_i) \]

\subsection{Límites y Continuidad}

Los conceptos de límite y continuidad se trasladan de manera directa a las funciones vectoriales, operando componente a componente.

\begin{mydefinition}{Límite de una función vectorial}
Si $\vec{r}(t) = (f_1(t), f_2(t), \dots, f_n(t))$, el límite cuando $t \to a$ se define como:
\[ \lim_{t \to a} \vec{r}(t) = \left(\lim_{t \to a} f_1(t),\; \lim_{t \to a} f_2(t),\; \dots,\; \lim_{t \to a} f_n(t)\right) \]
siempre que los límites escalares existan.
\end{mydefinition}

\begin{mydefinition}{Continuidad}
Una función vectorial $\vec{r}(t)$ es continua en $t = a$ si:
\[ \lim_{t \to a} \vec{r}(t) = \vec{r}(a) \]
Esto ocurre si y solo si cada función componente es continua en $t = a$.
\end{mydefinition}

La continuidad de $\vec{r}(t)$ en su dominio garantiza que la curva trazada en el espacio no tiene saltos ni discontinuidades: es una curva \textit{suave por partes}.


\subsection{Representación Gráfica}

La imagen geométrica de una función vectorial $\vec{r}(t)$, conforme $t$ recorre su dominio, es una curva en $\mathbb{R}^n$. Cada valor de $t$ determina la posición de un punto sobre dicha curva.

\subsubsection{Curvas planas en $\mathbb{R}^2$}

Si la componente $z$ es cero, la curva yace en el plano $xy$. Por ejemplo, sea: $\vec{r}(t) = (\cos t, \sin t, 0)$ con $t \in [0, 2\pi]$ traza el círculo unitario.

\begin{center}
\shorthandoff{>}
\begin{tikzpicture}[scale=2]
    \draw[->, gray, thick] (-1.4,0) -- (1.6,0) node[right] {$x$};
    \draw[->, gray, thick] (0,-1.4) -- (0,1.6) node[above] {$y$};

    \draw[very thick, blue] (1,0) arc (0:360:1);

    \draw[->, very thick, red] (0,0) -- ({cos(50)},{sin(50)}) node[above right] {$\vec{r}(t)$};

    \filldraw[red] ({cos(50)},{sin(50)}) circle (1.2pt);
    \node[right, font=\footnotesize] at ({cos(50)},{sin(50)}) {$(f(t), g(t))$};

    \node[below right, blue, font=\footnotesize] at (0.7,-0.7) {$\vec{r}(t) = (\cos t, \sin t)$};
    \filldraw[black] (0,0) circle (1pt) node[below left] {\footnotesize Origen};
\end{tikzpicture}
\end{center}

\subsubsection{La hélice circular en $\mathbb{R}^3$}

Uno de los ejemplos más representativos de curva en el espacio es la hélice circular:
\[ \vec{r}(t) = (a\cos t,\; a\sin t,\; bt), \quad t \in \mathbb{R} \]
donde $a$ controla el radio y $b$ el paso. Las proyecciones sobre el plano $xy$ forman un círculo, mientras que la curva asciende de forma uniforme.

\begin{center}
\shorthandoff{>}
\begin{tikzpicture}[x={(-0.5cm,-0.4cm)}, y={(1cm,0cm)}, z={(0cm,1cm)}, scale=1.8]
    \draw[->, gray, thick] (0,0,0) -- (2,0,0) node[below left] {$x$};
    \draw[->, gray, thick] (0,0,0) -- (0,2,0) node[right] {$y$};
    \draw[->, gray, thick] (0,0,0) -- (0,0,2.8) node[above] {$z$};

    \draw[very thick, blue, variable=\t, domain=0:360, samples=36, smooth]
        plot ({cos(\t)}, {sin(\t)}, {\t/360});
    \draw[very thick, blue, variable=\t, domain=360:720, samples=36, smooth]
        plot ({cos(\t)}, {sin(\t)}, {\t/360});

    \draw[dashed, gray, variable=\t, domain=0:360, samples=24, smooth]
        plot ({cos(\t)}, {sin(\t)}, 0);

    \draw[dashed, gray!50] (1,0,0) -- (1,0,2);

    \filldraw[red] ({cos(90)},{sin(90)},{0.25}) circle (1.5pt)
        node[right=2pt, black, font=\footnotesize] {$\vec{r}(t_0)$};

    \node[blue, font=\footnotesize] at (-2.0,-0.1,2.5)
        {$\vec{r}(t)=(\cos t,\,\sin t,\,\tfrac{t}{2\pi})$};
\end{tikzpicture}
\end{center}

\subsection{Derivadas de Funciones Vectoriales}

\begin{mydefinition}{Derivada de una función vectorial}
La derivada de $\vec{r}(t)$ se define de manera análoga a la derivada escalar:
\[ \vec{r}'(t) = \lim_{\Delta t \to 0} \frac{\vec{r}(t + \Delta t) - \vec{r}(t)}{\Delta t} \]
Si $\vec{r}(t) = (f_1(t), f_2(t), \dots ,f_n(t))$, entonces:
\[ \vec{r}'(t) = \bigl(f_1'(t),\; f_2'(t),\; \dots ,\; f_n'(t)\bigr) \]
siempre que las $n$ derivadas componentes existan. En ese caso se dice que $\vec{r}$ es \textbf{diferenciable} en $t$.
\end{mydefinition}

\subsubsection{Interpretación geométrica}

El vector $\vec{r}(t + \Delta t) - \vec{r}(t)$ une dos puntos de la curva. Al dividirlo entre $\Delta t$ y tomar el límite, se obtiene un vector que es tangente a la curva en el punto $\vec{r}(t)$, apuntando en la dirección del movimiento.

\begin{center}
\shorthandoff{>}
\begin{tikzpicture}[scale=1.4]
    \draw[->, gray] (-0.5,0) -- (4.5,0) node[right] {$x$};
    \draw[->, gray] (0,-0.5) -- (0,3) node[above] {$y$};

    \draw[very thick, blue, domain=0.2:3.8, samples=80]
        plot (\x, {0.4*\x*\x - 0.5*\x + 1.2});

    \coordinate (P) at (2, {0.4*4 - 1 + 1.2});
    \coordinate (Q) at (3, {0.4*9 - 1.5 + 1.2});

    \draw[->, very thick, purple] (P) -- (Q) node[midway, above left, font=\footnotesize] {$\vec{r}(t{+}\Delta t) - \vec{r}(t)$};

    \draw[->, very thick, red] (P) -- ++(1, {2*0.4*2 - 0.5}) node[right] {$\vec{r}'(t)$};

    \filldraw[black] (P) circle (1.5pt) node[below left] {$\vec{r}(t)$};
    \filldraw[black] (Q) circle (1.5pt) node[above right] {$\vec{r}(t{+}\Delta t)$};
\end{tikzpicture}
\end{center}

\subsubsection{Reglas de derivación}

\begin{mytheorem}{Reglas de derivación vectorial}
Sean $\vec{u}(t)$ y $\vec{v}(t)$ funciones vectoriales diferenciables y $c(t)$ una función escalar diferenciable. Entonces:
\begin{enumerate}
    \item $\dfrac{d}{dt}[\vec{u} + \vec{v}] = \vec{u}' + \vec{v}'$
    \item $\dfrac{d}{dt}[c\,\vec{u}] = c'\vec{u} + c\,\vec{u}'$
    \item $\dfrac{d}{dt}[\vec{u} \cdot \vec{v}] = \vec{u}' \cdot \vec{v} + \vec{u} \cdot \vec{v}'$
    \item $\dfrac{d}{dt}[\vec{u} \times \vec{v}] = \vec{u}' \times \vec{v} + \vec{u} \times \vec{v}'$
    \item $\dfrac{d}{dt}[\vec{u}(c(t))] = c'(t)\,\vec{u}'(c(t))$ \quad (Regla de la cadena)
\end{enumerate}
\end{mytheorem}
    
\subsection{Integrales de Funciones Vectoriales}

La integral de una función vectorial también se aplica componente a componente.

\begin{mydefinition}{Integral indefinida}
\[ \int \vec{r}(t)\, dt = \left(\int f_1(t)\,dt,\; \int f_2(t)\,dt,\; \dots ,\; \int f_n(t)\,dt\right) + \vec{C} \]
donde $\vec{C} = (C_1, C_2, \dots , C_n)$ es un \textbf{vector constante de integración}.
\end{mydefinition}

\begin{mydefinition}{Integral definida}
\[ \int_a^b \vec{r}(t)\, dt = \left(\int_a^b f_1(t)\,dt,\; \int_a^b f_2(t)\,dt,\; \dots ,\; \int_a^b f_n(t)\,dt\right) \]
\end{mydefinition}

El Teorema Fundamental del Cálculo también aplica: si $\vec{R}(t)$ es una antiderivada de $\vec{r}(t)$, es decir, $\vec{R}'(t) = \vec{r}(t)$, entonces:
\[ \int_a^b \vec{r}(t)\, dt = \vec{R}(b) - \vec{R}(a) \]

\subsection{Cinemática Vectorial}

La aplicación más natural de las funciones vectoriales es la descripción del movimiento de una partícula en el espacio. Si $\vec{r}(t)$ representa la \textbf{posición} de la partícula en el instante $t$, las sucesivas derivadas codifican cantidades cinemáticas con significado físico preciso.

\subsubsection{Posición, velocidad y aceleración}

\begin{mydefinition}{Cantidades cinemáticas}
Sea $\vec{r}(t) = (x(t), y(t), z(t))$ la función de posición de una partícula. Se definen:
\begin{itemize}
    \item \textbf{Vector velocidad:} $\vec{v}(t) = \vec{r}'(t) = \bigl(x'(t),\, y'(t),\, z'(t)\bigr)$
    \item \textbf{Rapidez escalar:} $v(t) = \|\vec{v}(t)\| = \|\vec{r}'(t)\|$
    \item \textbf{Vector aceleración:} $\vec{a}(t) = \vec{v}'(t) = \vec{r}''(t) = \bigl(x''(t),\, y''(t),\, z''(t)\bigr)$
\end{itemize}
\end{mydefinition}

Una propiedad importante de la cinemática vectorial es que el vector velocidad es siempre tangente a la trayectoria. Por definición de derivada, el vector $\vec{r}(t + \Delta t) - \vec{r}(t)$ une dos posiciones cercanas de la partícula. Al dividir entre $\Delta t$ y tomar el límite, el resultado apunta exactamente en la dirección en que la curva avanza en ese instante. Por lo tanto, $\vec{v}(t)$ es tangente a la curva y su módulo es la rapidez con la que la partícula la recorre.

\begin{center}
\shorthandoff{>}
\begin{tikzpicture}[scale=1.3]
    \draw[->, gray] (-0.5,0) -- (5.5,0) node[right] {$x$};
    \draw[->, gray] (0,-0.5) -- (0,3.5) node[above] {$y$};

    % Trayectoria (parábola orientada)
    \draw[very thick, blue, domain=0.4:5, samples=50]
        plot (\x, {-0.3*(\x-2.5)*(\x-2.5) + 3});

    % Punto en t=t0 (x=1.5)
    \def\px{1.5}
    \def\py{-0.3*(1.5-2.5)*(1.5-2.5)+3}  % = -0.3 + 3 = 2.7

    % Vector posición
    \draw[->, thick, gray!70] (0,0) -- (\px, 2.7) node[midway, left=2pt, font=\footnotesize, gray] {$\vec{r}(t)$};

    % Derivada de la parábola: y' = -0.6(x-2.5), en x=1.5: y' = -0.6*(-1) = 0.6
    % Vector velocidad (tangente): dirección (1, 0.6), normalizado para visual
    \draw[->, very thick, red] (\px, 2.7) -- ++(1.0, 0.6) node[above right] {$\vec{v}(t) = \vec{r}'(t)$};

    % Vector aceleración (2ª derivada, apunta hacia abajo: y''=-0.6)
    \draw[->, very thick, orange] (\px, 2.7) -- ++(0, -0.9) node[below right] {$\vec{a}(t) = \vec{r}''(t)$};

    % Punto
    \filldraw[black] (\px, 2.7) circle (2pt) node[below left] {$P = \vec{r}(t_0)$};

    % Punto futuro
    \def\px2{3.2}
    \draw[->, very thick, red, opacity=0.4] (\px2, {-0.3*(\px2-2.5)*(\px2-2.5)+3}) -- ++(1.0, {-0.6*(\px2-2.5)}) ;
    \filldraw[black, opacity=0.4] (\px2, {-0.3*(\px2-2.5)*(\px2-2.5)+3}) circle (1.5pt);

    \node[blue, font=\footnotesize] at (4.5, 0.8) {Trayectoria};
    \node[font=\footnotesize, align=left] at (4.8, 2.8) {\textcolor{red}{$\vec{v}$ siempre tangente}};
\end{tikzpicture}
\end{center}

La aceleración $\vec{a}(t)$ mide cómo cambia el \textit{vector} velocidad, no solo su magnitud. Esto tiene una consecuencia importante: una partícula puede tener aceleración no nula aunque su rapidez sea constante, simplemente porque está cambiando de dirección.

El caso más ilustrativo es el movimiento circular uniforme: Sea $\vec{r}(t) = (R\cos t,\; R\sin t)$, con radio $R$ y velocidad angular constante $1$ rad/s.
\begin{align*}
    \vec{v}(t) &= (-R\sin t,\; R\cos t) & \|\vec{v}(t)\| &= R \quad \text{(rapidez constante)} \\
    \vec{a}(t) &= (-R\cos t,\; -R\sin t) = -\vec{r}(t) & \|\vec{a}(t)\| &= R
\end{align*}

\begin{center}
\shorthandoff{>}
\begin{tikzpicture}[scale=1.4]
    \draw[->, gray] (-2.2,0) -- (2.5,0) node[right] {$x$};
    \draw[->, gray] (0,-2.2) -- (0,2.5) node[above] {$y$};

    % Círculo
    \draw[very thick, blue] (0,0) circle (1.8cm);

    % Punto en t=40 grados
    \def\ang{40}
    \def\R{1.8}
    \coordinate (P) at ({1.8*cos(\ang)}, {1.8*sin(\ang)});

    % Vector posición
    \draw[->, thick, gray!60] (0,0) -- (P) node[midway, above left, font=\footnotesize, gray] {$\vec{r}$};

    % Vector velocidad (tangente al círculo: dirección (-sin, cos))
    \draw[->, very thick, red] (P) -- ++({-\R*0.5*sin(\ang)}, {\R*0.5*cos(\ang)})
        node[above right] {$\vec{v}(t)$};

    % Vector aceleración (centrípeta: apunta hacia el centro)
    \draw[->, very thick, orange] (P) -- ++({-0.7*cos(\ang)}, {-0.7*sin(\ang)})
        node[below right] {$\vec{a}(t)$};

    \filldraw[black] (P) circle (1.5pt);
    \filldraw[black] (0,0) circle (1.5pt) node[below left] {$O$};

    \node[font=\footnotesize, align=center] at (-1.5,-1.7)
        {$\|\vec{v}\| = R$ \\$\vec{a} \perp \vec{v}$};
\end{tikzpicture}
\end{center}

La aceleración $\vec{a} = -\vec{r}$ apunta siempre hacia el centro del círculo (aceleración centrípeta), y es perpendicular a la velocidad en todo momento: $\vec{v} \cdot \vec{a} = 0$.

Por otro lado, dado que $\vec{v}(t) = \vec{r}'(t)$, se puede recuperar la posición integrando si se conoce la condición inicial:
\[ \vec{r}(t) = \vec{r}(t_0) + \int_{t_0}^{t} \vec{v}(u)\, du \]
De manera análoga, conociendo $\vec{a}(t)$ y las condiciones iniciales $\vec{v}(t_0)$ y $\vec{r}(t_0)$, se puede reconstruir toda la trayectoria.

Utilizando el ejemplo de un proyectil: Una pelota se lanza con velocidad inicial $\vec{v}_0 = (v_0\cos\alpha,\; v_0\sin\alpha)$ desde el origen. La única aceleración es la gravedad: $\vec{a}(t) = (0, -g)$.

Integrando una vez: $\vec{v}(t) = \vec{v}_0 + \int_0^t \vec{a}\, du = (v_0\cos\alpha,\; v_0\sin\alpha - gt)$.

Integrando de nuevo con $\vec{r}(0) = \vec{0}$:
\[ \vec{r}(t) = \left(v_0\cos\alpha\cdot t,\;\; v_0\sin\alpha\cdot t - \tfrac{1}{2}g t^2\right) \]

\begin{center}
\shorthandoff{>}
\begin{tikzpicture}[scale=1.1]
    \draw[->, gray] (-0.3,0) -- (5.5,0) node[right] {$x$};
    \draw[->, gray] (0,-0.3) -- (0,2.8) node[above] {$y$};

    % Parábola con v0=5, alpha=45, g=9.8: x=5cos45*t, y=5sin45*t - 4.9t^2
    % max height at t=5sin45/9.8 ≈ 0.361 s: y≈0.638
    % range at t=2*0.361≈0.722 s: x≈2.552
    % Scale so it looks good: multiply t by 1, scale x by 2
    % r(t) = (5cos45*t, 5sin45*t - 4.9t^2), t in [0, 0.722]
    % Scale x by 2 visually -> (3.536t*2, 3.536t - 4.9t^2) = (7.07t, 3.536t - 4.9t^2)
    % Let's parametrize with domain for nice plot: scale everything
    % Use: x = t, y = t*(1-t/T)*H where T=4.5, H=2.3
    \draw[very thick, blue, domain=0:4.5, samples=50]
        plot (\x, {2.3*(\x/4.5)*(1 - \x/4.5)*4});

    % Velocity vectors at a few points
    % dy/dx = 2.3*(1/4.5)*(1 - 2*x/(4.5))*4 = (2.3*4/4.5)*(1-2x/4.5)
    \foreach \xp/\sc in {0.5/0.5, 1.5/0.45, 2.25/0.4, 3.0/0.45, 4.0/0.5} {
        \pgfmathsetmacro{\yp}{2.3*(\xp/4.5)*(1-\xp/4.5)*4}
        \pgfmathsetmacro{\slope}{(2.3*4/4.5)*(1-2*\xp/4.5)}
        \pgfmathsetmacro{\norm}{sqrt(1+\slope*\slope)}
        \draw[->, red, thick]
            (\xp, \yp) -- ++({0.7/\norm}, {0.7*\slope/\norm});
    }

    % Gravity arrows
    \foreach \xp in {1.0, 2.25, 3.5} {
        \pgfmathsetmacro{\yp}{2.3*(\xp/4.5)*(1-\xp/4.5)*4}
        \draw[->, orange, thick] (\xp, \yp) -- ++(0, -0.5) node[right, font=\tiny] {};
    }

    \filldraw[black] (0,0) circle (2pt);
    \filldraw[black] (4.5,0) circle (2pt);

    \node[above right, red, font=\footnotesize] at (1.0, 2.5) {$\vec{v}(t)$ (tangente)};
    \node[right, orange, font=\footnotesize] at (3.9, 1.5) {$\vec{a} = (0,-g)$};
    \node[font=\footnotesize] at (2.25, -0.3) {Alcance};
\end{tikzpicture}
\end{center}

\subsubsection{Longitud de arco}

Intuitivamente, la longitud de una curva es el resultado de ``enderezarla''. Para calcularla de forma rigurosa, se utiliza el mismo principio que en el cálculo integral: aproximar con objetos sencillos (segmentos de recta) y tomar un límite.

Dividiendo el intervalo $[a, b]$ en $n$ subintervalos iguales mediante los puntos $a = t_0 < t_1 < \cdots < t_n = b$, con paso $\Delta t = \frac{b-a}{n}$. Cada $t_i$ define un punto $P_i = \vec{r}(t_i)$ sobre la curva.

Uniendo los puntos consecutivos $P_i$ y $P_{i+1}$ con segmentos de recta, formando una poligonal inscrita en la curva. La longitud del segmento $i$-ésimo es:
\[ \Delta L_i = \|P_{i+1} - P_i\| = \|\vec{r}(t_{i+1}) - \vec{r}(t_i)\| \]

\begin{center}
\shorthandoff{>}
\begin{tikzpicture}[scale=1.4]
    \draw[->, gray] (-0.3,0) -- (4.5,0) node[right] {$x$};
    \draw[->, gray] (0,-0.3) -- (0,3) node[above] {$y$};

    % Curva suave
    \draw[very thick, blue!40, domain=0.2:4.2, samples=60]
        plot (\x, {0.35*\x*\x - 0.9*\x + 2.3});

    % Puntos de la partición
    \foreach \xv/\lab in {0.5/t_0, 1.2/t_1, 1.9/t_2, 2.6/t_3, 3.3/t_4, 4.0/t_5} {
        \pgfmathsetmacro{\yv}{0.35*\xv*\xv - 0.9*\xv + 2.3}
        \filldraw[red] (\xv, \yv) circle (1.8pt);
        \node[font=\tiny, below=3pt] at (\xv, 0) {$\lab$};
        \draw[dashed, gray!50, thin] (\xv, 0) -- (\xv, \yv);
    }

    % Segmentos de la poligonal
    \foreach \xa/\xb in {0.5/1.2, 1.2/1.9, 1.9/2.6, 2.6/3.3, 3.3/4.0} {
        \pgfmathsetmacro{\ya}{0.35*\xa*\xa - 0.9*\xa + 2.3}
        \pgfmathsetmacro{\yb}{0.35*\xb*\xb - 0.9*\xb + 2.3}
        \draw[very thick, red!80] (\xa, \ya) -- (\xb, \yb);
    }

    % Etiqueta de un segmento
    \pgfmathsetmacro{\ymid}{(0.35*1.9*1.9-0.9*1.9+2.3+0.35*2.6*2.6-0.9*2.6+2.3)/2}
    \node[font=\footnotesize, right=2pt] at (2.25, \ymid) {$\Delta L_i = \|\vec{r}(t_{i+1}){-}\vec{r}(t_i)\|$};

\end{tikzpicture}
\end{center}

Por la definición de derivada, cuando $\Delta t$ es pequeño:
\[ \vec{r}(t_i + \Delta t) - \vec{r}(t_i) \approx \vec{r}'(t_i)\,\Delta t \]
Por lo tanto, la longitud de cada cuerda se aproxima por:
\[ \Delta L_i = \|\vec{r}(t_{i+1}) - \vec{r}(t_i)\| \approx \|\vec{r}'(t_i)\|\,\Delta t \]

La siguiente figura ilustra esto geométricamente: al acercar $P_{i+1}$ a $P_i$, la cuerda (roja) es casi indistinguible del vector tangente escalado $\vec{r}'(t_i)\Delta t$ (azul).

\begin{center}
\shorthandoff{>}
\begin{tikzpicture}[scale=2.2]
    % Arco de curva local
    \draw[very thick, blue!40, domain=0.5:2.5, samples=40]
        plot (\x, {0.18*\x*\x - 0.3*\x + 1.2});

    % Dos puntos cercanos
    \def\xa{1.0}  \def\ya{0.18*1.0*1.0-0.3*1.0+1.2}   % ≈ 1.08
    \def\xb{1.6}  \def\yb{0.18*1.6*1.6-0.3*1.6+1.2}   % ≈ 1.261

    \coordinate (Pi)  at (\xa, {0.18*1.0*1.0-0.3*1.0+1.2});
    \coordinate (Pi1) at (\xb, {0.18*1.6*1.6-0.3*1.6+1.2});

    % Cuerda (Δr)
    \draw[->, very thick, red] (Pi) -- (Pi1)
        node[midway, below=4pt] {$\Delta\vec{r}_i = \vec{r}(t_{i+1})-\vec{r}(t_i)$};

    % Tangente escalada r'(t)*Δt (pendiente en xa: 0.36*1-0.3=0.06, Δx=0.6)
    % Vector tangente real: (1, 0.06), escalar a Δx=0.6
    \draw[->, very thick, blue, dashed] (Pi) -- ++(0.6, 0.036)
        node[above=6pt] {$\vec{r}'(t_i)\,\Delta t$};

    % Puntos
    \filldraw[red] (Pi) circle (1pt) node[left=2pt] {$P_i$};
    \filldraw[red] (Pi1) circle (1pt) node[right=2pt] {$P_{i+1}$};
\end{tikzpicture}
\end{center}

La longitud total de la poligonal es:
\[ L_n = \sum_{i=0}^{n-1} \Delta L_i \approx \sum_{i=0}^{n-1} \|\vec{r}'(t_i)\|\,\Delta t \]
Esta es una suma de Riemann para la función $\|\vec{r}'(t)\|$. Al hacer $n \to \infty$ (equivalente a $\Delta t \to 0$), la poligonal converge a la curva real y la suma converge a la integral:

\begin{mydefinition}{Longitud de arco}
Sea $\vec{r}(t)$ continuamente diferenciable en $[a, b]$. La longitud de arco de la curva es:
\[ L = \lim_{n\to\infty} \sum_{i=0}^{n-1}\|\vec{r}'(t_i)\|\,\Delta t = \int_a^b \|\vec{r}'(t)\|\, dt \]
Expandiendo la norma en sus componentes:
\[ L = \int_a^b \sqrt{\,[x'(t)]^2 + [y'(t)]^2 + [z'(t)]^2}\; dt \]
\end{mydefinition}

\subsubsection{Parámetro de longitud de arco}

Es posible reparametrizar la curva usando como parámetro la longitud de arco $s$ acumulada desde $t = a$:
\[ s(t) = \int_a^t \|\vec{r}'(u)\|\, du \]
Por el Teorema Fundamental del Cálculo, $\dfrac{ds}{dt} = \|\vec{r}'(t)\| = v(t)$. Esto permite cambiar entre los dos parámetros:
\begin{itemize}
    \item $s$ como función de $t$: mide cuánta longitud se ha recorrido hasta el instante $t$.
    \item $t$ como función de $s$: invirtiendo la relación, permite localizar en qué punto de la curva se está habiendo recorrido una distancia $s$.
\end{itemize}

\subsection{Los vectores $\hat{T}$, $\hat{N}$ y $\hat{B}$ de Frenet-Serret}

Para analizar la geometría de una curva en el espacio (sin importar cómo se parametrice), se construye un sistema de referencia móvil asociado a cada punto de la curva: el triedro de Frenet-Serret, formado por tres vectores unitarios y mutuamente ortogonales $\hat{T}$, $\hat{N}$ y $\hat{B}$.

\begin{center}
\shorthandoff{>}
\begin{tikzpicture}[x={(-0.5cm,-0.4cm)}, y={(1cm,0cm)}, z={(0cm,1cm)}, scale=2.2]
    \draw[->, gray] (0,0,0) -- (1.8,0,0) node[below left, font=\footnotesize] {$x$};
    \draw[->, gray] (0,0,0) -- (0,1.8,0) node[right, font=\footnotesize] {$y$};
    \draw[->, gray] (0,0,0) -- (0,0,2.0) node[above, font=\footnotesize] {$z$};

    % Hélice r(t)=(cos t, sin t, t/360) para t en grados, una vuelta y media
    \draw[thick, blue!70, variable=\t, domain=0:540, samples=48, smooth]
        plot ({cos(\t)}, {sin(\t)}, {\t/540});

    % Punto en t=90 grados: P=(0, 1, 90/540) = (0, 1, 1/6)
    % T = (-1, 0, ~0.16) normalizado ≈ (-0.988, 0, 0.157)
    % N = (0, -1, 0)
    % B = T x N ≈ (0.157, 0, 0.988)
    \coordinate (P) at ({cos(90)}, {sin(90)}, {90/540});

    \draw[->, very thick, red]         (P) -- ++(-0.55, 0,    0.09) node[left, font=\footnotesize] {$\hat{T}$};
    \draw[->, very thick, green!60!black] (P) -- ++(0,   -0.5, 0)    node[right, font=\footnotesize] {$\hat{N}$};
    \draw[->, very thick, orange]      (P) -- ++(0.09,  0,    0.55) node[above, font=\footnotesize] {$\hat{B}$};

    \filldraw[black] (P) circle (1.2pt) node[below right, font=\footnotesize] {$P$};
\end{tikzpicture}
\end{center}

\begin{center}
{\footnotesize
\begin{tabular}{ccc}
    \textcolor{red}{$\hat{T}$: tangente (dirección de movimiento)} &
    \textcolor{green!60!black}{$\hat{N}$: normal principal (hacia donde dobla)} &
    \textcolor{orange}{$\hat{B}$: binormal ($\hat{T}\times\hat{N}$)}
\end{tabular}
}
\end{center}

\subsubsection{Vector tangente unitario $\hat{T}$}

El vector tangente unitario apunta en la dirección del movimiento y se obtiene normalizando la derivada de la posición:

\begin{mydefinition}{Vector tangente unitario}
\[ \hat{T}(t) = \frac{\vec{r}'(t)}{\|\vec{r}'(t)\|} \]
\end{mydefinition}

\subsubsection{Curvatura $\kappa$}

La curvatura mide cuánto cambia la dirección de la curva por unidad de longitud de arco, es decir, qué tan rápido se dobla la curva.

\begin{mydefinition}{Curvatura}
La curvatura $\kappa$ de una curva en $\mathbb{R}^3$ se define como:
\[ \kappa = \left\|\frac{d\hat{T}}{ds}\right\| = \frac{\|\hat{T}'(t)\|}{\|\vec{r}'(t)\|} \]
Una fórmula alternativa, más práctica para el cálculo, es:
\[ \kappa = \frac{\|\vec{r}'(t) \times \vec{r}''(t)\|}{\|\vec{r}'(t)\|^3} \]
\end{mydefinition}


\subsubsection{Vector normal principal $\hat{N}$}

El vector normal principal apunta hacia el interior de la curva. Este vector se obtiene normalizando la derivada del vector tangente unitario. Dado que $\hat{T}$ es unitario, su derivada es siempre perpendicular a $\hat{T}$, lo que garantiza que $\hat{N}$ también es perpendicular a $\hat{T}$.

\begin{mydefinition}{Vector normal principal}
\[ \hat{N}(t) = \frac{\hat{T}'(t)}{\|\hat{T}'(t)\|} \]
\end{mydefinition}

\subsubsection{Vector binormal $\hat{B}$}

El vector binormal completa el triedro de Frenet-Serret como un tercer vector ortogonal a los dos anteriores. Es el producto vectorial de $\hat{T}$ y $\hat{N}$, y apunta en la dirección perpendicular al plano osculador formado por $\hat{T}$ y $\hat{N}$.

\begin{mydefinition}{Vector binormal}
\[ \hat{B}(t) = \hat{T}(t) \times \hat{N}(t) \]
\end{mydefinition}

El plano determinado por $\hat{N}$ y $\hat{B}$ se llama plano normal, y el plano determinado por $\hat{T}$ y $\hat{B}$ se llama plano rectificante.

\subsubsection{Resumen del triedro de Frenet-Serret}

\begin{mytheorem}{Fórmulas de Frenet-Serret}
Los tres vectores del triedro satisfacen las siguientes relaciones diferenciales:
\[
\frac{d\hat{T}}{ds} = \kappa\,\hat{N}, \qquad
\frac{d\hat{N}}{ds} = -\kappa\,\hat{T} + \tau\,\hat{B}, \qquad
\frac{d\hat{B}}{ds} = -\tau\,\hat{N}
\]
donde $\kappa \geq 0$ es la curvatura y $\tau$ es la torsión.
\end{mytheorem}

\subsubsection{Descomposición de la aceleración}

Una consecuencia importante del triedro de Frenet-Serret es que la aceleración siempre yace en el plano osculador y se puede descomponer en dos componentes ortogonales:

\begin{mytheorem}{Componentes tangencial y normal de la aceleración}
\[ \vec{a}(t) = a_T\,\hat{T} + a_N\,\hat{N} \]
donde:
\[ a_T = \frac{d^2s}{dt^2} = \frac{\vec{r}'(t) \cdot \vec{r}''(t)}{\|\vec{r}'(t)\|} \qquad \text{(componente tangencial)} \]
\[ a_N = \kappa \left(\frac{ds}{dt}\right)^2 = \frac{\|\vec{r}'(t) \times \vec{r}''(t)\|}{\|\vec{r}'(t)\|} \qquad \text{(componente normal)} \]
La componente $a_T$ refleja el cambio de rapidez; la componente $a_N$ refleja el cambio de dirección.
\end{mytheorem}

\begin{center}
\shorthandoff{>}
% Curva: y = 0.5 + 0.7*sin(70*x)  [angulo en grados]
% Punto P en x=1.0:
%   y(1)  = 0.5 + 0.7*sin(70°)  ≈ 1.158
%   y'(1) = 0.7*cos(70°)*(70π/180) ≈ 0.292   →  T̂ ≈ (0.960, 0.280)
%   y''(1) < 0  (concava hacia abajo)         →  N̂ ≈ (0.280, -0.960)
% Componentes elegidas: a_T = 0.7, a_N = 0.5
%   a_T·T̂ ≈ (0.672,  0.196)
%   a_N·N̂ ≈ (0.140, -0.480)
%   a     ≈ (0.812, -0.284)   [suma exacta → paralelogramo cierra]
\begin{tikzpicture}[scale=1.5]
    \draw[->, gray, thin] (-0.2,0) -- (3.5,0) node[right] {$x$};
    \draw[->, gray, thin] (0,-0.1) -- (0,2.4) node[above] {$y$};

    \draw[very thick, blue, domain=0:3.2, samples=50]
        plot (\x, {0.5 + 0.7*sin(70*\x)});

    \coordinate (P)  at (1.0,   1.158);   % punto base
    \coordinate (AT) at (1.672, 1.354);   % P + a_T·T̂
    \coordinate (AN) at (1.140, 0.678);   % P + a_N·N̂
    \coordinate (A)  at (1.812, 0.874);   % P + a  (= AT + a_N·N̂ = AN + a_T·T̂)

    % Paralelogramo (líneas punteadas)
    \draw[dashed, gray!60, thin] (AT) -- (A);
    \draw[dashed, gray!60, thin] (AN) -- (A);

    % Aceleración total
    \draw[->, very thick, purple] (P) -- (A)
        node[right] {$\vec{a}$};

    % Componente tangencial (a lo largo de T̂, paralela a la curva)
    \draw[->, very thick, red] (P) -- (AT)
        node[above right] {$a_T\,\hat{T}$};

    % Componente normal (perpendicular a T̂, hacia el centro de curvatura)
    \draw[->, very thick, green!60!black] (P) -- (AN)
        node[below right] {$a_N\,\hat{N}$};

    % Marca de ángulo recto en P entre T̂ y N̂
    % paso ε=0.09 a lo largo de T̂=(0.960,0.280) y N̂=(0.280,-0.960)
    \draw[thick]
        (1.0864, 1.1832)            % P + ε·T̂
        -- (1.1116, 1.0968)         % P + ε·T̂ + ε·N̂
        -- (1.0252, 1.0716);        % P + ε·N̂

    \filldraw[black] (P) circle (1.5pt) node[above left] {$P$};
\end{tikzpicture}
\end{center}

\subsection{Ejercicios}

\begin{enumerate}
    \item Determinar el dominio natural de cada función vectorial:
\begin{enumerate}
    \item $\vec{r}(t) = \left(\sqrt{9 - t^2},\; \ln(t + 1),\; \dfrac{1}{t}\right)$
    \item $\vec{r}(t) = \left(e^t,\; \arcsin(2t - 1),\; \sqrt{t}\right)$
\end{enumerate}
    \item Sea $\vec{r}(t) = \left(\dfrac{t^2 - 4}{t - 2},\; \cos t,\; t^2 + 1\right)$.
\begin{enumerate}
    \item Calcular $\displaystyle\lim_{t \to 2} \vec{r}(t)$ componente a componente.
    \item Determinar si $\vec{r}(t)$ es continua en $t = 2$. Justificar.
    \item Redefinir $\vec{r}(2)$ si es necesario para hacer la función continua en ese punto.
\end{enumerate}
    \item Identificar y describir la curva trazada por cada función vectorial. En cada caso, encontrar la curva equivalente en forma cartesiana eliminando el parámetro $t$.
\begin{enumerate}
    \item $\vec{r}(t) = (2\cos t,\; 3\sin t),\quad t \in [0, 2\pi]$
    \item $\vec{r}(t) = (t,\; t^2 - 1),\quad t \in \mathbb{R}$
    \item $\vec{r}(t) = (2\cos t,\; 2\sin t,\; 3t),\quad t \in [0, 4\pi]$
\end{enumerate}
    \item Sea $\vec{r}(t) = (e^t \cos t,\; e^t \sin t,\; e^t)$. Calcular $\vec{r}'(t)$ y $\vec{r}''(t)$. Verificar la regla del producto para la función escalar $c(t) = e^t$ y el vector $\vec{u}(t) = (\cos t, \sin t, 1)$, comprobando que:
\[ \vec{r}'(t) = c'(t)\,\vec{u}(t) + c(t)\,\vec{u}'(t) \]
    \item Calcular las siguientes integrales vectoriales:
\begin{enumerate}
    \item $\displaystyle\int_0^{\pi} \bigl(\cos t,\; \sin t,\; t\bigr)\, dt$
    \item $\displaystyle\int \left(t^2,\; e^{2t},\; \frac{1}{1+t^2}\right) dt$
    \item Verificar el Teorema Fundamental del Cálculo: si $\vec{R}(t) = \left(\frac{t^3}{3},\; \frac{e^{2t}}{2},\; \arctan t\right)$, comprobar que $\vec{R}'(t) = \left(t^2,\; e^{2t},\; \frac{1}{1+t^2}\right)$.
\end{enumerate}
    \item Una partícula tiene función de posición $\vec{r}(t) = (t^2 - 1,\; 2t,\; t^3/3)$ con $t \geq 0$ (distancias en metros, $t$ en segundos).
\begin{enumerate}
    \item Encontrar el vector velocidad $\vec{v}(t)$ y la rapidez escalar $v(t) = \|\vec{v}(t)\|$.
    \item Encontrar el vector aceleración $\vec{a}(t)$.
    \item Evaluar $\vec{v}(1)$, $v(1)$ y $\vec{a}(1)$. ¿Cuál es la dirección de movimiento en $t=1$?
    \item Determinar si existe algún instante $t > 0$ en el que la aceleración sea paralela a la velocidad.
\end{enumerate}
    \item Una pelota es lanzada desde el punto $\vec{r}_0 = (0, 1.5, 0)$ m con velocidad inicial $\vec{v}_0 = (10, 0, 8)$ m/s. La única aceleración es la gravedad: $\vec{a}(t) = (0, 0, -9.8)$ m/s$^2$.
\begin{enumerate}
    \item Determinar $\vec{v}(t)$ y $\vec{r}(t)$ integrando la aceleración con las condiciones iniciales dadas.
    \item Encontrar el instante $t^*$ en que la pelota alcanza su altura máxima en $z$.
    \item Calcular la posición en $t^*$.
    \item Encontrar el instante en que la pelota golpea el plano $z = 0$ y el punto de impacto.
\end{enumerate}
    \item Calcular la longitud de arco de cada curva en el intervalo indicado:
\begin{enumerate}
    \item $\vec{r}(t) = (3t,\; 4\cos t,\; 4\sin t)$ en $[0, 2\pi]$.
    \item $\vec{r}(t) = (t^2,\; \frac{2}{3}t^3,\; t)$ en $[0, 1]$.

    \textit{Sugerencia para (b):} Simplificar $\|\vec{r}'(t)\|^2$ buscando un cuadrado perfecto.
\end{enumerate}
    \item Para la curva $\vec{r}(t) = (\cos t,\; \sin t,\; t)$:
\begin{enumerate}
    \item Calcular el vector tangente unitario $\hat{T}(t)$.
    \item Calcular la curvatura $\kappa$ usando la fórmula $\kappa = \dfrac{\|\hat{T}'(t)\|}{\|\vec{r}'(t)\|}$.
    \item Verificar el resultado usando la fórmula alternativa $\kappa = \dfrac{\|\vec{r}'(t) \times \vec{r}''(t)\|}{\|\vec{r}'(t)\|^3}$.
    \item ¿La curvatura es constante? ¿Qué dice esto sobre la forma de la curva?
\end{enumerate}
    \item Para la curva $\vec{r}(t) = (t,\; t^2,\; \frac{2}{3}t^3)$, encontrar en el punto correspondiente a $t = 1$:
\begin{enumerate}
    \item El vector tangente unitario $\hat{T}(1)$.
    \item El vector normal principal $\hat{N}(1)$.
    \item El vector binormal $\hat{B}(1) = \hat{T}(1) \times \hat{N}(1)$.
    \item Verificar que los tres vectores son mutuamente ortogonales y unitarios.
\end{enumerate}
    \item Una partícula se mueve según $\vec{r}(t) = (2\sin t,\; 2\cos t,\; 3t)$ (distancias en metros).
\begin{enumerate}
    \item Calcular la velocidad $\vec{v}(t)$, la aceleración $\vec{a}(t)$ y la rapidez $v(t)$.
    \item Descomponer la aceleración en sus componentes tangencial y normal:
    \[ a_T = \frac{\vec{r}'(t)\cdot\vec{r}''(t)}{\|\vec{r}'(t)\|}, \qquad a_N = \frac{\|\vec{r}'(t)\times\vec{r}''(t)\|}{\|\vec{r}'(t)\|} \]
    \item Verificar que $\|\vec{a}\|^2 = a_T^2 + a_N^2$.
    \item Calcular la curvatura $\kappa$ de la trayectoria. ¿Cuál es su radio de curvatura $\rho = 1/\kappa$?
\end{enumerate}
    \item Un robot industrial mueve su efector final siguiendo la trayectoria $\vec{r}(t) = (t - \sin t,\; 1 - \cos t,\; t)$ para $t \in [0, 2\pi]$ (la curva sobre un cilindro llamada \textit{cicloide cilíndrica}).
\begin{enumerate}
    \item Calcular $\vec{r}'(t)$ y $\|\vec{r}'(t)\|$. Identificar para qué valores de $t \in [0, 2\pi]$ la velocidad se anula.
    \item Calcular la longitud total de la trayectoria $L = \displaystyle\int_0^{2\pi} \|\vec{r}'(t)\|\, dt$.

    \textit{Sugerencia:} Usar la identidad $1 - \cos t = 2\sin^2\!\left(\frac{t}{2}\right)$.
    \item Calcular la curvatura $\kappa(t) = \dfrac{\|\vec{r}'(t)\times\vec{r}''(t)\|}{\|\vec{r}'(t)\|^3}$ en $t = \pi$.
    \item Determinar si la función de longitud de arco $s(t) = \displaystyle\int_0^t \|\vec{r}'(u)\|\,du$ es invertible en $(0, 2\pi)$. ¿Qué condición debe cumplirse y por qué?
\end{enumerate}
    \item Demostrar, partiendo únicamente de la definición de derivada como límite, que el producto punto satisface la regla del producto:
\[ \frac{d}{dt}\bigl[\vec{u}(t) \cdot \vec{v}(t)\bigr] = \vec{u}'(t) \cdot \vec{v}(t) + \vec{u}(t) \cdot \vec{v}'(t) \]

    \item En el análisis cinemático de mecanismos y robots, la aceleración de un punto que sigue una trayectoria curvilínea se descompone en dos componentes físicamente distintas. Demostrar rigurosamente que, para cualquier partícula que sigue una curva suave $\vec{r}(t)$, el vector aceleración admite la descomposición:
\[ \vec{a}(t) = a_T\,\hat{T} + a_N\,\hat{N} \]
y derivar las expresiones explícitas:
\[ a_T = \frac{\vec{r}'(t)\cdot\vec{r}''(t)}{\|\vec{r}'(t)\|}, \qquad a_N = \frac{\|\vec{r}'(t)\times\vec{r}''(t)\|}{\|\vec{r}'(t)\|} \]
Concluir verificando el resultado sobre el movimiento circular $\vec{r}(t) = (R\cos\omega t,\; R\sin\omega t,\; 0)$ e interpretar físicamente ambas componentes.

\textit{Sugerencia:} Partir de $\vec{v} = v\hat{T}$ (con $v = \|\vec{r}'\|$), diferenciar con la regla del producto, y expandir $\hat{T}'$ usando la regla de la cadena junto con la primera fórmula de Frenet-Serret, $\dfrac{d\hat{T}}{ds} = \kappa\hat{N}$. Para las fórmulas computacionales, derivar $v^2 = \vec{r}'\cdot\vec{r}'$ y aplicar la identidad de Lagrange.
\end{enumerate}


% \section{Funciones escalares}

La mayoría de los fenómenos físicos e ingenieriles dependen de múltiples variables simultáneamente: la temperatura en un sólido depende de las coordenadas $(x, y, z)$, la presión de un gas depende de la temperatura y el volumen, la resistencia de una viga depende de su geometría y del material. Para modelar estos fenómenos se introducen las funciones escalares de varias variables, que extienden el concepto de función real $f: \mathbb{R} \to \mathbb{R}$ a dominios de dimensión superior.

\subsection{Definición}

\begin{mydefinition}{Función escalar de $n$ variables}
Una función escalar de $n$ variables es una regla $f$ que asigna a cada punto $(x_1, x_2, \dots, x_n)$ en un subconjunto $D \subseteq \mathbb{R}^n$ un único número real:
\[ f: D \subseteq \mathbb{R}^n \longrightarrow \mathbb{R}, \qquad (x_1, x_2, \dots, x_n) \longmapsto f(x_1, x_2, \dots, x_n) \]
Al conjunto $D$ se le llama dominio de $f$, y al conjunto de todos los valores posibles de salida se le llama imagen o rango:
\[ \text{Im}(f) = \{f(x_1, \dots, x_n) \in \mathbb{R} : (x_1, \dots, x_n) \in D\} \]
\end{mydefinition}

\subsection{Dominio y Rango}

El dominio natural de $f$ es el conjunto de todos los puntos para los que la expresión de $f$ está matemáticamente definida. Las restricciones más frecuentes son:

\begin{itemize}
    \item \textbf{Denominadores:} $g(x,y) \neq 0$
    \item \textbf{Raíces de índice par:} $g(x,y) \geq 0$
    \item \textbf{Logaritmos:} $g(x,y) > 0$
    \item \textbf{Funciones inversas:} respetar el dominio de $\arcsin$, $\arccos$, etc.
\end{itemize}

\textbf{Ejemplo 1.} Sea $f(x,y) = \sqrt{9 - x^2 - y^2}$. La condición $9 - x^2 - y^2 \geq 0$ define el dominio:
\[ D = \{(x,y) \in \mathbb{R}^2 : x^2 + y^2 \leq 9\} \]
Este es el círculo de radio $3$ centrado en el origen. La imagen es $\text{Im}(f) = [0, 3]$: el mínimo $0$ se alcanza en la circunferencia $x^2+y^2=9$ y el máximo $3$ en el origen.

\begin{center}
\shorthandoff{>}
\begin{tikzpicture}[scale=0.75]
    \draw[->, gray] (-3.8,0) -- (3.8,0) node[right] {$x$};
    \draw[->, gray] (0,-3.8) -- (0,3.8) node[above] {$y$};

    \fill[blue!15] (0,0) circle (3cm);
    \draw[very thick, blue] (0,0) circle (3cm);

    \node at (0,0) [font=\footnotesize] {$D$};
    \node[above right, blue, font=\footnotesize] at (2.5, 2.0) {$x^2+y^2=9$};
    \draw[<-, blue, thin] (2.12, 2.12) -- (2.3, 1.9);
\end{tikzpicture}
\end{center}

\textbf{Ejemplo 2.} Sea $g(x,y,z) = \dfrac{\ln(1 - x^2 - y^2)}{z}$. Se necesitan dos condiciones simultáneas:
\[
    1 - x^2 - y^2 > 0 \;\Rightarrow\; x^2 + y^2 < 1
    \qquad\text{y}\qquad
    z \neq 0
\]
\[ D = \{(x,y,z) \in \mathbb{R}^3 : x^2 + y^2 < 1 \text{ y } z \neq 0\} \]

Para determinar la imagen de $f$, es útil: identificar los valores extremos en la frontera del dominio, determinar si $f$ es continua, y verificar si hay máximos o mínimos globales.

\subsection{Gráfica de funciones con dominio en $ \mathbb{R}^2 $}

Para una función de dos variables $z = f(x, y)$, la gráfica es la superficie en $\mathbb{R}^3$ formada por todos los puntos de la forma $(x, y, f(x,y))$:
\[ \text{Graf}(f) = \bigl\{(x,\, y,\, f(x,y)) \in \mathbb{R}^3 : (x,y) \in D\bigr\} \]
Cada punto del plano $xy$ se eleva o hunde según el valor que $f$ le asigna.

\textbf{Ejemplo 1.} Paraboloide elíptico: $f(x,y) = x^2 + y^2$

Esta superficie tiene un mínimo global en el origen $(0,0,0)$ y crece sin límite al alejarse. En cualquier plano vertical que pase por el eje $z$, la sección es una parábola.

\begin{center}
\begin{tikzpicture}
\begin{axis}[
    view={45}{30},
    axis lines=center,
    xlabel={$x$}, ylabel={$y$}, zlabel={$z$},
    xlabel style={anchor=north east},
    ylabel style={anchor=north west},
    zlabel style={anchor=south},
    xmin=-2.5, xmax=2.5,
    ymin=-2.5, ymax=2.5,
    zmin=0,    zmax=5,
    xtick={-2,-1,0,1,2}, ytick={-2,-1,0,1,2}, ztick={1,2,3,4},
    tick label style={font=\footnotesize},
    width=8cm, height=8cm,
    colormap/viridis,
]
    \addplot3[
        surf, opacity=0.85,
        domain=-2:2, domain y=-2:2,
        samples=18,
    ] {x^2 + y^2};
\end{axis}
\end{tikzpicture}
\end{center}

\textbf{Ejemplo 2.} Silla de montar (paraboloide hiperbólico): $f(x,y) = x^2 - y^2$

Esta superficie sube como una parábola en la dirección $x$ y baja como una parábola en la dirección $y$. El punto $(0,0,0)$ es un punto de silla: no es ni máximo ni mínimo local.

\begin{center}
\begin{tikzpicture}
\begin{axis}[
    view={45}{25},
    axis lines=center,
    xlabel={$x$}, ylabel={$y$}, zlabel={$z$},
    xlabel style={anchor=north east},
    ylabel style={anchor=north west},
    zlabel style={anchor=south},
    xmin=-2.5, xmax=2.5,
    ymin=-2.5, ymax=2.5,
    zmin=-4.5, zmax=4.5,
    xtick={-2,-1,0,1,2}, ytick={-2,-1,0,1,2},
    tick label style={font=\footnotesize},
    width=8cm, height=8cm,
    colormap/cool,
]
    \addplot3[
        surf, opacity=0.85,
        domain=-2:2, domain y=-2:2,
        samples=18,
    ] {x^2 - y^2};
    \addplot3[mark=*, mark size=2pt, black] coordinates {(0,0,0)};
\end{axis}
\node[font=\footnotesize] at (6.5, 0.8) {Punto de silla};
\end{tikzpicture}
\end{center}

\subsection{Curvas de nivel de funciones con dominio en $\mathbb{R}^2$}

Visualizar una superficie en papel es limitado. Una alternativa es proyectar sobre el plano $xy$ las intersecciones de la superficie con planos horizontales $z = c$.

\begin{mydefinition}{Curvas de nivel}
La curva de nivel de $f(x,y)$ para el valor $c \in \mathbb{R}$ es el conjunto:
\[ C_c = \{(x,y) \in D : f(x,y) = c\} \]
La colección de curvas de nivel para distintos valores de $c$ se llama mapa de contorno de $f$. Es la representación que usan los mapas topográficos: cada curva une puntos con la misma altitud.
\end{mydefinition}

La curva de nivel $C_c$ es la proyección sobre el plano $xy$ de la intersección de la superficie $z = f(x,y)$ con el plano horizontal $z = c$. Curvas de nivel muy juntas indican que la función cambia rápidamente; curvas muy separadas indican cambio lento.

\textbf{Ejemplo 1.} Curvas de nivel del paraboloide: $f(x,y) = x^2 + y^2$

La ecuación $x^2 + y^2 = c$ define círculos de radio $\sqrt{c}$ (solo para $c > 0$). Para $c = 0$ el único punto es el origen.

\begin{center}
\shorthandoff{>}
\begin{tikzpicture}[scale=1.1]
    \draw[->, gray] (-2.7,0) -- (2.7,0) node[right] {$x$};
    \draw[->, gray] (0,-2.7) -- (0,2.7) node[above] {$y$};

    % Curvas de nivel: circulos de radio sqrt(c)
    \foreach \c/\col in {1/blue!60, 2/blue!80, 4/blue}{
        \pgfmathsetmacro{\rad}{sqrt(\c)}
        \draw[\col, thick] (0,0) circle (\rad cm);
        \pgfmathsetmacro{\labelx}{\rad * 0.707}
        \node[\col, font=\footnotesize, fill=white, inner sep=1pt]
            at (\labelx+0.05, \labelx+0.05) {$c=\c$};
    }

    \filldraw[black] (0,0) circle (1.5pt) node[below left, font=\footnotesize] {$c=0$};

    \node[font=\footnotesize, align=center] at (0,-2.5)
        {$x^2+y^2 = c$ \quad (círculos de radio $\sqrt{c}$)};
\end{tikzpicture}
\end{center}

Nótese que las curvas están más separadas en el origen (la función sube lentamente) y más juntas hacia afuera (la función sube más rápido). Esto refleja que la pendiente de la parábola es mayor lejos del vértice.

\textbf{Ejemplo 2.} Curvas de nivel de la silla de montar: $f(x,y) = x^2 - y^2$

La ecuación $x^2 - y^2 = c$ define hipérbolas (para $c \neq 0$) o dos rectas cruzadas (para $c = 0$).

\begin{center}
\shorthandoff{>}
\begin{tikzpicture}[scale=0.9]
    \draw[->, gray] (-2.7,0) -- (2.7,0) node[right] {$x$};
    \draw[->, gray] (0,-2.7) -- (0,2.7) node[above] {$y$};

    % c=0: asíntotas y=x, y=-x
    \draw[gray!60, dashed] (-2.5,-2.5) -- (2.5,2.5);
    \draw[gray!60, dashed] (-2.5,2.5) -- (2.5,-2.5);
    \node[gray, font=\tiny] at (1.8, 1.5) {$c=0$};

    % c=1: hiperbola x^2 - y^2 = 1 (abre en x), rama derecha e izq
    \draw[blue!70, thick, domain=-2.2:2.2, samples=40]
        plot ({sqrt(1+\x*\x)}, {\x});
    \draw[blue!70, thick, domain=-2.2:2.2, samples=40]
        plot ({-sqrt(1+\x*\x)}, {\x});
    \node[blue!70, font=\tiny, fill=white] at (1.6, 0.3) {$c{=}1$};

    % c=2: hiperbola x^2 - y^2 = 2
    \draw[blue, thick, domain=-2.0:2.0, samples=40]
        plot ({sqrt(2+\x*\x)}, {\x});
    \draw[blue, thick, domain=-2.0:2.0, samples=40]
        plot ({-sqrt(2+\x*\x)}, {\x});
    \node[blue, font=\tiny, fill=white] at (2.3, 0.3) {$c{=}2$};

    % c=-1: hiperbola y^2 - x^2 = 1 (abre en y)
    \draw[red!70, thick, domain=-2.2:2.2, samples=40]
        plot ({\x}, {sqrt(1+\x*\x)});
    \draw[red!70, thick, domain=-2.2:2.2, samples=40]
        plot ({\x}, {-sqrt(1+\x*\x)});
    \node[red!70, font=\tiny, fill=white] at (0.3, 1.5) {$c{=}{-}1$};

    % c=-2
    \draw[red, thick, domain=-1.8:1.8, samples=40]
        plot ({\x}, {sqrt(2+\x*\x)});
    \draw[red, thick, domain=-1.8:1.8, samples=40]
        plot ({\x}, {-sqrt(2+\x*\x)});
    \node[red, font=\tiny, fill=white] at (0.3, 2.2) {$c{=}{-}2$};

    \filldraw[black] (0,0) circle (1.5pt);
    \node[font=\footnotesize, align=center] at (0,-2.6) {$x^2 - y^2 = c$ \quad (hipérbolas)};
\end{tikzpicture}
\end{center}

Las hipérbolas azules ($c > 0$) abren en la dirección $x$ y las rojas ($c < 0$) abren en la dirección $y$, reflejando que la función sube en $x$ y baja en $y$.

\subsection{Superficies de nivel de funciones con dominio en $\mathbb{R}^3$}

Para una función de tres variables $w = f(x, y, z)$, la gráfica sería un objeto en $\mathbb{R}^4$, imposible de visualizar directamente. El análogo de las curvas de nivel son las superficies de nivel.

\begin{mydefinition}{Superficie de nivel}
La superficie de nivel de $f(x,y,z)$ para el valor $c \in \mathbb{R}$ es el conjunto:
\[ S_c = \{(x,y,z) \in D : f(x,y,z) = c\} \]
Cada superficie de nivel es un subconjunto del espacio $\mathbb{R}^3$ y representa todos los puntos donde $f$ toma el mismo valor $c$.
\end{mydefinition}

\textbf{Ejemplo 1.} $f(x,y,z) = x^2 + y^2 + z^2$

Las superficies de nivel son $x^2 + y^2 + z^2 = c$, que son esferas concéntricas de radio $\sqrt{c}$ centradas en el origen (para $c > 0$).

\begin{center}
\shorthandoff{>}
\begin{tikzpicture}[x={(-0.5cm,-0.4cm)}, y={(1cm,0cm)}, z={(0cm,1cm)}, scale=1.4]
    \draw[->, gray] (0,0,0) -- (3,0,0) node[below left, font=\footnotesize] {$x$};
    \draw[->, gray] (0,0,0) -- (0,3,0) node[right, font=\footnotesize] {$y$};
    \draw[->, gray] (0,0,0) -- (0,0,3) node[above, font=\footnotesize] {$z$};

    % Esfera c=1 (radio 1): arcos representativos
    \draw[blue!50, thick, variable=\t, domain=0:360, samples=36, smooth]
        plot ({cos(\t)}, {sin(\t)}, 0);
    \draw[blue!50, thick, variable=\t, domain=0:180, samples=24, smooth]
        plot ({cos(\t)}, 0, {sin(\t)});
    \draw[blue!50, thick, variable=\t, domain=0:180, samples=24, smooth]
        plot (0, {cos(\t)}, {sin(\t)});
    \node[blue!70, font=\footnotesize] at (-1.5, -0.3, 0.8) {$c=1$};

    % Esfera c=4 (radio 2): arcos representativos
    \draw[blue, thick, variable=\t, domain=0:360, samples=36, smooth]
        plot ({2*cos(\t)}, {2*sin(\t)}, 0);
    \draw[blue, thick, variable=\t, domain=0:180, samples=24, smooth]
        plot ({2*cos(\t)}, 0, {2*sin(\t)});
    \draw[blue, thick, variable=\t, domain=0:180, samples=24, smooth]
        plot (0, {2*cos(\t)}, {2*sin(\t)});
    \node[blue, font=\footnotesize] at (-2.8, -0.3, 1.5) {$c=4$};

    \filldraw[black] (0,0,0) circle (1.2pt) node[below right, font=\footnotesize] {$c=0$};
\end{tikzpicture}
\end{center}

\textbf{Ejemplo 2.} $g(x,y,z) = x^2 + y^2$

Las superficies de nivel de $g$ son $x^2 + y^2 = c$, que son cilindros circulares de radio $\sqrt{c}$ con eje en la dirección $z$. La función no depende de $z$, de modo que para cada nivel $c$ hay una familia de circunferencias apiladas verticalmente.

\subsection{Ejercicios}

\begin{enumerate}
    \item Determinar el dominio natural de cada función escalar e indicar a qué conjunto geométrico corresponde en $\mathbb{R}^2$ o $\mathbb{R}^3$:
\begin{enumerate}
    \item $f(x,y) = \sqrt{16 - x^2 - 4y^2}$
    \item $g(x,y) = \ln(y - x^2)$
    \item $h(x,y) = \dfrac{1}{\sqrt{x^2 + y^2 - 1}}$
\end{enumerate}
    \item Determinar el dominio natural de cada función de tres variables y describir geométricamente la región en $\mathbb{R}^3$:
\begin{enumerate}
    \item $f(x,y,z) = \sqrt{4 - x^2 - y^2 - z^2}$
    \item $g(x,y,z) = \dfrac{\arcsin(z)}{\sqrt{x^2 + y^2}}$
\end{enumerate}
    \item Para la función $f(x,y) = \sqrt{x + y} + \sqrt{4 - x^2}$:
\begin{enumerate}
    \item Encontrar el dominio natural $D$ e identificarlo geométricamente (trazar o describir la región).
    \item Determinar la imagen de $f$.
    \item Verificar que el punto $(1, 0)$ pertenece a $D$ y calcular $f(1, 0)$.
\end{enumerate}
    \item Identificar y nombrar la superficie definida por cada función $z = f(x,y)$, describiendo su forma geométrica (paraboloide, silla de montar, plano, cono, cilindro, etc.) y señalando características clave como vértice, eje de simetría o punto de silla:
\begin{enumerate}
    \item $f(x,y) = 3x - 2y + 5$
    \item $f(x,y) = x^2 + 4y^2$
    \item $f(x,y) = \sqrt{x^2 + y^2}$
    \item $f(x,y) = 4 - x^2 - y^2$
\end{enumerate}
    \item Para $f(x,y) = x^2 + 4y^2$, trazar las curvas de nivel correspondientes a $c = 0, 1, 4, 9$ en el plano $xy$.
\begin{enumerate}
    \item Determinar la forma geométrica de cada curva de nivel e indicar sus semiejes.
    \item ¿Qué ocurre con el espaciado entre curvas de nivel conforme $c$ aumenta? ¿Qué dice esto sobre la ``pendiente'' de la superficie?
    \item ¿Para qué valores de $c$ existen curvas de nivel? ¿Cuál es la imagen de $f$?
\end{enumerate}
    \item Trazar las curvas de nivel de $f(x,y) = xy$ para $c = -2, -1, 0, 1, 2$ e identificar su forma.
\begin{enumerate}
    \item Mostrar que para $c \neq 0$ cada curva de nivel es una hipérbola y determinar sus asíntotas.
    \item ¿Qué ocurre con la curva de nivel $c = 0$? Describir el conjunto $\{(x,y): xy = 0\}$.
    \item Usando el mapa de contorno, describir el comportamiento de $f$ cerca del origen: ¿tiene máximo, mínimo o punto de silla?
\end{enumerate}
    \item Se dan las ecuaciones de las curvas de nivel de una función $f(x,y)$. En cada caso, identificar una posible expresión para $f(x,y)$ justificando la elección:
\begin{enumerate}
    \item Las curvas de nivel son rectas paralelas de la forma $2x - y = c$.
    \item Las curvas de nivel son circunferencias centradas en $(1, 2)$.
    \item Las curvas de nivel son parábolas de la forma $y = x^2 + c$.
    \item Las curvas de nivel para $c > 0$ son elipses, y no existen para $c < 0$.
\end{enumerate}
    \item Para cada función de tres variables, identificar y describir geométricamente sus superficies de nivel para $c = 1, 4, 9$:
\begin{enumerate}
    \item $f(x,y,z) = x^2 + y^2$ \quad (¿depende de $z$?)
    \item $g(x,y,z) = x^2 + y^2 - z^2$
    \item $h(x,y,z) = 2x + 3y - z$
\end{enumerate}
Para el inciso (b), distinguir los casos $c > 0$, $c = 0$ y $c < 0$.
    \item Un mapa topográfico de una montaña muestra curvas de nivel equiespaciadas en altura ($\Delta z = 100$ m) a distancias horizontales que varían de un punto a otro. En la zona $A$ las curvas están muy juntas; en la zona $B$ están muy separadas.
\begin{enumerate}
    \item ¿En cuál zona es más pronunciada la pendiente del terreno? Justificar usando la interpretación de las curvas de nivel.
    \item Si $z = f(x,y)$ modela la altitud, y en la zona $A$ una curva de nivel $c=500$ m y la siguiente $c=600$ m están separadas apenas $50$ m horizontalmente, mientras que en la zona $B$ están separadas $500$ m, estimar la ``pendiente media'' en cada zona como $\Delta z / \Delta d$.
    \item Proponer una función $f(x,y)$ sencilla cuyas curvas de nivel tengan la forma general de un volcán (alto en el centro, nivel descendente hacia afuera).
\end{enumerate}
    \item La distribución de temperatura en una placa metálica cuadrada $[-2,2]\times[-2,2]$ está modelada por:
\[ T(x,y) = 100 - 10x^2 - 20y^2 \quad (^\circ\text{C}) \]
\begin{enumerate}
    \item Encontrar el dominio y la imagen de $T$ sobre la placa.
    \item Calcular la temperatura en los puntos $(0,0)$, $(1,1)$ y $(2, 0)$.
    \item Trazar las curvas de nivel (isotermas) para $T = 100, 80, 60, 40, 20$ °C. ¿Qué forma tienen? ¿Cuáles de ellas intersectan la placa?
    \item Determinar el conjunto de puntos de la placa donde la temperatura es de exactamente $70$ °C e identificarlo geométricamente.
    \item Una partícula se mueve sobre la placa siguiendo la trayectoria $\vec{r}(t) = (t, t^2 - 1)$ para $t \in [-1, 1]$. Calcular $T(x(t), y(t))$ a lo largo de esta trayectoria y encontrar el punto de temperatura máxima.
\end{enumerate}
\end{enumerate}

\section{Límites y Continuidad de Funciones Escalares}

En el cálculo de una variable, al calcular $\lim_{x \to a} f(x)$, el punto $x$ solo puede acercarse a $a$ por dos caminos: por la izquierda o por la derecha. En el cálculo multivariable, la situación es radicalmente más compleja: el punto $(x, y)$ puede acercarse a $(x_0, y_0)$ por infinitos caminos distintos: rectas, parábolas, espirales, etc. Esto hace que demostrar la existencia de un límite sea más complicado.

\subsection{Definición formal del límite}

\begin{mydefinition}{Límite de una función de dos variables}
Sea $f: D \subseteq \mathbb{R}^2 \to \mathbb{R}$ y sea $(x_0, y_0)$ un punto de acumulación del dominio $D$. Se dice que el límite de $f$ cuando $(x,y)$ tiende a $(x_0, y_0)$ es el número real $L$, y se escribe:
\[ \lim_{(x,y) \to (x_0,y_0)} f(x,y) = L \]
si para todo $\varepsilon > 0$ existe $\delta > 0$ tal que:
\[ 0 < \sqrt{(x-x_0)^2 + (y-y_0)^2} < \delta \implies |f(x,y) - L| < \varepsilon \]
En otras palabras: todos los puntos del dominio dentro del disco de radio $\delta$ centrado en $(x_0, y_0)$ — excepto el propio centro — tienen imágenes dentro del intervalo $(L-\varepsilon, L+\varepsilon)$.
\end{mydefinition}

\subsubsection{Interpretación geométrica de la definición $\varepsilon$-$\delta$}

\begin{center}
\shorthandoff{>}
\begin{tikzpicture}[scale=1.2]
    % ---- Plano xy (izquierda) ----
    \draw[->, gray] (-2.5,0) -- (2.5,0) node[right] {$x$};
    \draw[->, gray] (0,-2.2) -- (0,2.2) node[above] {$y$};

    % Disco delta
    \fill[blue!12] (0.8, 0.5) circle (1.1cm);
    \draw[blue, thick] (0.8, 0.5) circle (1.1cm);
    \draw[dashed, blue] (0.8,0.5) -- ++(0.778, 0.778)
        node[above right, font=\footnotesize, blue] {$\delta$};

    % Punto (x0,y0)
    \filldraw[black] (0.8, 0.5) circle (2pt)
        node[below left, font=\footnotesize] {$(x_0,y_0)$};

    % Punto (x,y) interior — bajado para no solapar etiqueta delta
    \filldraw[red] (1.4, 0.9) circle (1.5pt)
        node[right, font=\footnotesize, red] {$(x,y)$};

    \node[blue, font=\footnotesize] at (-1.2, -1.7) {Disco: $0 < \|(x,y)-(x_0,y_0)\| < \delta$};

    % ---- Flecha de mapeo ----
    \draw[->, thick, gray!70] (2.6, 0) -- (4.3, 0)
        node[midway, above, font=\footnotesize] {$f$};

    % ---- Eje z (derecha) — desplazado a x=5 para mayor claridad ----
    \draw[->, gray] (5, -2) -- (5, 2.5) node[above] {$z$};

    % Intervalo (L-eps, L+eps) — ampliado para separar etiquetas
    \draw[orange, very thick] (5, 0.2) -- (5, 1.8);
    \filldraw[orange] (5, 1.0) circle (2pt)
        node[right, font=\footnotesize] {$L$};
    \draw[dashed, orange] (4.6, 0.2) -- (5.5, 0.2)
        node[right, font=\footnotesize] {$L - \varepsilon$};
    \draw[dashed, orange] (4.6, 1.8) -- (5.5, 1.8)
        node[right, font=\footnotesize] {$L + \varepsilon$};

    % Valor f(x,y) — en el lado izquierdo del eje, despejado de L y L+-eps
    \filldraw[red] (5, 1.35) circle (1.5pt)
        node[left, font=\footnotesize, red] {$f(x,y)$};

    \node[orange, font=\footnotesize] at (6.3, -1.5)
        {$|f(x,y) - L| < \varepsilon$};
\end{tikzpicture}
\end{center}

La definición exige que sin importar el camino por el que $(x,y)$ se acerque a $(x_0,y_0)$, el valor $f(x,y)$ siempre deba converger al mismo $L$. Este es el requisito clave que distingue el límite multivariable del univariable.

La definición se extiende a funciones de $n$ variables reemplazando el disco por una bola en $\mathbb{R}^n$: $0 < \|\mathbf{x} - \mathbf{x}_0\| < \delta \implies |f(\mathbf{x}) - L| < \varepsilon$, donde $\|\cdot\|$ es la norma euclídea en $\mathbb{R}^n$.


\subsection{Propiedades algebraicas del límite}

Al igual que en una variable, los límites multivariables satisfacen las reglas aritméticas estándar.

\begin{mytheorem}{Álgebra de límites}
Si $\displaystyle\lim_{(x,y)\to(x_0,y_0)} f(x,y) = L$ y $\displaystyle\lim_{(x,y)\to(x_0,y_0)} g(x,y) = M$, entonces:
\begin{enumerate}
    \item $\displaystyle\lim (f \pm g) = L \pm M$
    \item $\displaystyle\lim (f \cdot g) = L \cdot M$
    \item $\displaystyle\lim \frac{f}{g} = \frac{L}{M}$, \quad siempre que $M \neq 0$
    \item $\displaystyle\lim [f(x,y)]^n = L^n$, \quad para $n$ entero positivo
\end{enumerate}
\end{mytheorem}

\subsection{Cálculo de límites}

Existen tres estrategias principales para evaluar un límite multivariable.

\subsubsection{Estrategia 1: Sustitución directa}

Si $f$ es un polinomio o una función racional bien definida en $(x_0, y_0)$, el límite se calcula simplemente evaluando:
\[ \lim_{(x,y) \to (x_0,y_0)} f(x,y) = f(x_0, y_0) \]

\textbf{Ejemplo.}
\[ \lim_{(x,y) \to (1,2)} \frac{x^2 + 3y}{x + y} = \frac{1 + 6}{1 + 2} = \frac{7}{3} \]

Cuando la sustitución produce una indeterminación ($\frac{0}{0}$, $\frac{\infty}{\infty}$, etc.), es necesario usar técnicas algebraicas: factorización, racionalización o cambio de variable.

\textbf{Ejemplo con indeterminación.}
\[
\lim_{(x,y)\to(0,0)} \frac{x^2 - y^2}{x - y}
= \lim_{(x,y)\to(0,0)} \frac{(x-y)(x+y)}{x-y}
= \lim_{(x,y)\to(0,0)} (x + y) = 0
\]

\subsubsection{Estrategia 2: Método de trayectorias}

Este método se usa para demostrar que un límite no existe. Si se puede encontrar al menos dos trayectorias que se acercan al mismo punto pero producen valores límite distintos, entonces el límite no existe.

\begin{mynote}{Principio de trayectorias}
Si existen dos caminos $C_1$ y $C_2$ que convergen a $(x_0, y_0)$ tales que:
\[ \lim_{\substack{(x,y)\to(x_0,y_0)\\(x,y)\in C_1}} f(x,y) = L_1 \neq L_2 = \lim_{\substack{(x,y)\to(x_0,y_0)\\(x,y)\in C_2}} f(x,y) \]
entonces $\displaystyle\lim_{(x,y)\to(x_0,y_0)} f(x,y)$ \textbf{no existe}.
\end{mynote}

\begin{center}
\shorthandoff{>}
\begin{tikzpicture}[scale=1.3]
    \draw[->, gray] (-2.5,0) -- (2.5,0) node[right] {$x$};
    \draw[->, gray] (0,-2) -- (0,2.5) node[above] {$y$};

    % Trayectoria y=0 (eje x)
    \draw[->, very thick, blue] (-2, 0) -- (-0.1, 0);
    \node[blue, font=\footnotesize, below] at (-1.5, 0) {$y=0$};

    % Trayectoria y=x
    \draw[->, very thick, red] (-1.6, -1.6) -- (-0.1, -0.1);
    \node[red, font=\footnotesize, below left] at (-1.3, -1.3) {$y=x$};

    % Trayectoria y=x^2
    \draw[->, very thick, green!60!black, domain=-1.4:0, samples=30]
        plot (\x, {\x*\x});
    \node[green!60!black, font=\footnotesize] at (-1.7, 1.3) {$y=x^2$};

    % Trayectoria y = -2x
    \draw[->, very thick, orange] (-1.0, 2.0) -- (-0.05, 0.1);
    \node[orange, font=\footnotesize, left] at (-0.9, 1.8) {$y=-2x$};

    % Trayectoria curva
    \draw[->, very thick, purple, domain=-2:0, samples=40]
        plot ({\x}, {sin(60*\x)});
    \node[purple, font=\footnotesize] at (-1.7, -1.5) {curva};

    % Punto límite
    \filldraw[black] (0,0) circle (2.5pt) node[above right] {$(0,0)$};

    \node[font=\footnotesize, align=center] at (1.5,-1.5)
        {\textbf{Infinitos caminos}\\hacia $(0,0)$};
\end{tikzpicture}
\end{center}

\textbf{Ejemplo.} Demostrar que $\displaystyle\lim_{(x,y)\to(0,0)} \frac{xy}{x^2+y^2}$ no existe.

\textit{Trayectoria $C_1$: $y = 0$ (eje $x$).} Al sustituir $y=0$:
\[ \frac{x \cdot 0}{x^2 + 0} = 0 \xrightarrow{x\to 0} \mathbf{0} \]

\textit{Trayectoria $C_2$: $y = x$.} Al sustituir $y=x$:
\[ \frac{x \cdot x}{x^2 + x^2} = \frac{x^2}{2x^2} = \frac{1}{2} \xrightarrow{x\to 0} \mathbf{\tfrac{1}{2}} \]

Como $0 \neq \tfrac{1}{2}$, el límite no existe.

Para funciones homogéneas o cocientes de polinomios en el origen, es útil probar la familia de rectas $y = mx$. Si el límite a lo largo de $y = mx$ depende de $m$, el límite general no existe.

\textbf{Ejemplo con familia $y = mx$.} Sea $f(x,y) = \dfrac{x^2 y}{x^4 + y^2}$. A lo largo de $y = mx$:
\[ f(x, mx) = \frac{x^2 \cdot mx}{x^4 + m^2x^2} = \frac{mx^3}{x^2(x^2 + m^2)} = \frac{mx}{x^2 + m^2} \xrightarrow{x\to 0} 0 \]
Todas las rectas dan $0$. Pero a lo largo de la parábola $y = x^2$:
\[ f(x, x^2) = \frac{x^2 \cdot x^2}{x^4 + x^4} = \frac{x^4}{2x^4} = \frac{1}{2} \]
El límite no existe, aunque todas las rectas hubieran dado el mismo resultado. Esto ilustra que pasar la prueba de rectas \textbf{no es suficiente} para garantizar la existencia del límite.

\subsubsection{Estrategia 3: Coordenadas polares}

Para límites en el origen $(0,0)$, el cambio a coordenadas polares $x = r\cos\theta$, $y = r\sin\theta$ transforma el problema:
\[ (x,y) \to (0,0) \iff r \to 0^+ \quad \text{(para cualquier } \theta) \]

\begin{mynote}{Criterio polar}
Si al sustituir $x = r\cos\theta$, $y = r\sin\theta$ la expresión de $f$ se puede acotar como $|f| \leq g(r)$ donde $g(r) \to 0$ cuando $r \to 0^+$ independientemente de $\theta$, entonces:
\[ \lim_{(x,y)\to(0,0)} f(x,y) = 0 \]
Si la expresión resultante depende de $\theta$ de forma irreductible, el límite probablemente no existe.
\end{mynote}

\textbf{Ejemplo 1: límite que existe.} Sea $\displaystyle\lim_{(x,y)\to(0,0)}\frac{x^3}{x^2+y^2}$.

En polares: $x^2 + y^2 = r^2$, $x^3 = r^3\cos^3\theta$:
\[ \frac{r^3\cos^3\theta}{r^2} = r\cos^3\theta \]
Dado que $|\cos\theta| \leq 1$, se tiene $|r\cos^3\theta| \leq r \to 0$. Por lo tanto, el límite es $\mathbf{0}$.

\textbf{Ejemplo 2: límite que no existe.} Sea $\displaystyle\lim_{(x,y)\to(0,0)}\frac{x^2 - y^2}{x^2+y^2}$.

En polares:
\[ \frac{r^2\cos^2\theta - r^2\sin^2\theta}{r^2} = \cos^2\theta - \sin^2\theta = \cos(2\theta) \]
El resultado depende de $\theta$: no converge a un valor único. El límite no existe.

\subsubsection{Teorema del emparedado}

\begin{mytheorem}{Teorema del emparedado}
Si existen funciones $g$ y $h$ tales que $g(x,y) \leq f(x,y) \leq h(x,y)$ en un entorno de $(x_0,y_0)$, y además:
\[ \lim_{(x,y)\to(x_0,y_0)} g(x,y) = \lim_{(x,y)\to(x_0,y_0)} h(x,y) = L \]
entonces $\displaystyle\lim_{(x,y)\to(x_0,y_0)} f(x,y) = L$.
\end{mytheorem}

Este teorema es la herramienta formal que justifica el criterio polar: si se acota $|f| \leq r \cdot M$ donde $M$ es una constante (o función acotada en $\theta$), entonces $-rM \leq f \leq rM$ y ambas cotas tienden a $0$.

\textbf{Ejemplo.} $\displaystyle\lim_{(x,y)\to(0,0)} \frac{x^2 y}{x^2+y^2}$.

Usando $|y| \leq \sqrt{x^2 + y^2} = r$ y $x^2 \leq x^2 + y^2 = r^2$:
\[ \left|\frac{x^2 y}{x^2+y^2}\right| \leq \frac{x^2 |y|}{x^2+y^2} \leq \frac{r^2 \cdot r}{r^2} = r \to 0 \]
Por el Teorema del Emparedado, el límite es $\mathbf{0}$.

\subsection{Continuidad}

\begin{mydefinition}{Continuidad en un punto}
Una función $f: D \subseteq \mathbb{R}^2 \to \mathbb{R}$ es continua en el punto $(x_0, y_0) \in D$ si se satisfacen simultáneamente las tres condiciones:
\begin{enumerate}
    \item $f(x_0, y_0)$ está definida.
    \item $\displaystyle\lim_{(x,y)\to(x_0,y_0)} f(x,y)$ existe.
    \item $\displaystyle\lim_{(x,y)\to(x_0,y_0)} f(x,y) = f(x_0, y_0)$.
\end{enumerate}
Si $f$ es continua en todos los puntos de $D$, se dice que $f$ es continua en $D$.
\end{mydefinition}

\subsubsection{Funciones continuas elementales}

\begin{mytheorem}{Continuidad de funciones elementales}
Las siguientes funciones son continuas en todo su dominio natural:
\begin{itemize}
    \item Todo polinomio $p(x,y) = \sum a_{ij} x^i y^j$.
    \item Toda función racional $p(x,y)/q(x,y)$ donde $q(x,y) \neq 0$.
    \item Las funciones trigonométricas, exponenciales, logarítmicas y raíces aplicadas a expresiones continuas.
    \item La composición de funciones continuas: si $g: \mathbb{R}^2 \to \mathbb{R}$ es continua y $\phi: \mathbb{R} \to \mathbb{R}$ es continua, entonces $\phi \circ g$ es continua.
\end{itemize}
\end{mytheorem}

\subsubsection{Extensión continua de una función}

Cuando una función presenta una indeterminación en un punto, puede ser posible extenderla de manera continua redefiniendo su valor en ese punto.

\textbf{Ejemplo.} Sea $f(x,y) = \dfrac{x^3 + y^3}{x^2 + y^2}$ para $(x,y) \neq (0,0)$. En polares:
\[ \frac{r^3(\cos^3\theta + \sin^3\theta)}{r^2} = r(\cos^3\theta + \sin^3\theta) \]
Como $|\cos^3\theta + \sin^3\theta| \leq 2$, se tiene que $|f| \leq 2r \to 0$. Por tanto:
\[ \lim_{(x,y)\to(0,0)} f(x,y) = 0 \]
Definiendo $f(0,0) = 0$, la función extendida es continua en el origen.

\subsubsection{Discontinuidades y su clasificación}

Una función puede ser discontinua en un punto $(x_0, y_0)$ por tres razones:
\begin{itemize}
    \item \textbf{No está definida:} $(x_0,y_0) \notin D$.
    \item \textbf{El límite no existe:} la función oscila o toma valores distintos según el camino.
    \item \textbf{El límite existe pero no coincide con $f(x_0,y_0)$:} discontinuidad evitable.
\end{itemize}

\end{document}